\subsection{The hybrid factorization carrier $\operatorname{Ran}^{\mathrm{hyb}}(E)$}
\label{subsec:hybrid-ran-carrier}

The first entry of the Dirac--Igusa pro-object is the carrier on which
chiral operations live.  An ordinary Beilinson--Drinfeld factorization
algebra on $\operatorname{Ran}(E)$ \cite[\S3.1]{BD} sees only a finite
multiset of points of $E$ and the operations attached to their
collisions; this is the projection-finite stratum.  The Igusa product
\(\Delta_5\) carries a Borcherds variable $s$ whose positive powers
record positive elliptic-degree wrapping.  An effective one-cycle on
\(K3\times E\) of positive elliptic degree has a horizontal component
and therefore projects \emph{onto} \(E\).  Locality on
$\operatorname{Ran}(E)$ alone has no mechanism for it.

\begin{proposition}[Positive elliptic degree projects onto \(E\);
\textit{proved}]
\label{prop:positive-elliptic-degree-projects}
\label{lem:wrapped-projection-obstruction}
Let
\[
Z=\sum_i m_i[C_i]\in Z_1(K3\times E),\qquad m_i>0,
\]
be an effective one-cycle with irreducible curve components \(C_i\).
If
\[
(p_E)_*Z=b[E],\qquad b>0,
\]
then \(p_E(|Z|)=E\).  Equivalently, at least one component \(C_i\)
maps dominantly to \(E\).  If a retained object \(A\) has a support
realization row identifying \(\operatorname{cyc}_1(A)=Z\) with the
effective fundamental one-cycle of \(\operatorname{Supp}_1 A\), then
\[
b_R^{\mathrm{geom}}(A)>0\quad\Longrightarrow\quad
p_E(\operatorname{Supp}_1 A)=E .
\]
Thus such an object is not localized at finitely many points of \(E\).
\end{proposition}

\begin{proof}
For each irreducible component \(C_i\), the proper image \(p_E(C_i)\)
is either a point or all of the irreducible curve \(E\).  In the first
case \((p_E)_*[C_i]=0\); in the second case
\((p_E)_*[C_i]=\deg(C_i/E)[E]\) with positive degree.  Since
\[
(p_E)_*Z=\sum_i m_i(p_E)_*[C_i]=b[E]
\]
and \(b>0\), at least one component has positive degree over \(E\).
That component maps dominantly onto \(E\), hence
\[
E=p_E(C_i)\subset p_E(|Z|)\subset E.
\]
The support statement follows when the support-realization row
identifies \(Z\) with the effective fundamental cycle of the
one-dimensional support of \(A\).
\end{proof}

The packet
\[
\texttt{certificates/hybrid/positive\_elliptic\_degree\_projection}
\]
verifies
\(\texttt{POSITIVE\_ELLIPTIC\_DEGREE\_PROJECTION\_VERIFIED}\):
the effective-cycle theorem is proved, while the retained object row,
effective support-cycle realization row, positive degree row, object
projection row, and transition row are recorded as missing open
obligations.

The hybrid carrier replaces the Ran space by a two-stratum object that
records, simultaneously, finite local supports and wrapped elliptic
legs.  At every finite Harder--Narasimhan height $R$ the datum is a
finite-coloured correspondence base, not a Ran space with a scalar
attached.

\begin{definition}[Finite-stage hybrid Ran prestack]
\label{def:hybrid-ran-prestack}
\label{def:hybrid-ran-base}
Fix a finite HN height \(R\).  The hybrid carrier at height \(R\) is
first a prestack
\[
\operatorname{Ran}^{\mathrm{hyb}}_{R,\mathrm{pre}}(E):
\mathrm{Aff}_{\C}^{op}\longrightarrow\infty\textup{-}\mathrm{Gpd}.
\]
For a test affine \(T\), a \(T\)-point is represented by the following
finite data:
\begin{enumerate}
\item a finite local index set \(I\) and maps
\(x_i:T\to E\), \(i\in I\);
\item retained local colour labels
\(\alpha_i\in\Gamma_R^{\mathrm{loc}}\);
\item a finite wrapped index set \(J\), retained wrapped colour labels
\(\eta_j\in\Gamma_R^{\mathrm{wr}}\), and \(T\)-families
\[
W_j\in\mathcal M_{\eta_j,R}^{\mathrm{wr,rig}}(T)
\]
in the rigidified wrapped prequotient stacks, equipped with their
anchor maps
\[
\lambda_{\eta_j,R}(W_j):T\longrightarrow E .
\]
\end{enumerate}
Two representatives are identified by finite relabelling of \(I\) and
\(J\), and by the usual Ran collision maps on the local coordinates.
Equivalently,
\[
\operatorname{Ran}^{\mathrm{hyb}}_{R,\mathrm{pre}}(E)
=
\operatorname*{colim}_{(I,J,\alpha,\eta)\in\mathsf{Fin}^{\mathrm{hyb}}_R}
\left(
E^I\times\prod_{j\in J}
\mathcal M_{\eta_j,R}^{\mathrm{wr,rig}}
\right),
\]
where \(\mathsf{Fin}^{\mathrm{hyb}}_R\) is the finite indexing category
generated by relabellings and local collision maps, and where
\(\alpha:I\to\Gamma_R^{\mathrm{loc}}\),
\(\eta:J\to\Gamma_R^{\mathrm{wr}}\).  This is a prestack
definition.  For a topology \(\tau\), the stack
\[
\operatorname{Ran}^{\mathrm{hyb}}_{R,\tau}(E)
:=a_\tau\operatorname{Ran}^{\mathrm{hyb}}_{R,\mathrm{pre}}(E)
\]
is used only after the descent rows for \(a_\tau\) have been supplied.
The scalar object
\[
\bigsqcup_{b\ge0}\operatorname{Ran}(E)\times\{b\}
\]
is only the degree shadow; it is not the hybrid prestack.
\end{definition}

The packet
\[
\texttt{certificates/hybrid/hybrid\_ran\_prestack\_definition}
\]
verifies
\(\texttt{HYBRID\_RAN\_PRESTACK\_DEFINITION\_VERIFIED}\): the
finite-stage carrier is defined as a prestack functor, while the
local-stratum component rows, wrapped-stratum component rows, descent
rows, stackification rows, and transition rows are recorded as missing
open obligations.

\begin{definition}[Local \(b=0\) stratum as a subprestack]
\label{def:hybrid-local-stratum-prestack}
At finite HN height \(R\), the local stratum is the full subprestack
\[
\operatorname{Ran}^{\mathrm{loc}}_{R,\mathrm{pre}}(E)
\subset
\operatorname{Ran}^{\mathrm{hyb}}_{R,\mathrm{pre}}(E)
\]
cut out by the conditions \(J=\varnothing\) and
\(\alpha:I\to\Gamma_R^{\mathrm{loc}}\), where
\[
\Gamma_R^{\mathrm{loc}}=\{\gamma\in\Gamma_{\sigma,S}^{\mathrm{HN}}(R)
\mid b_R^{\mathrm{geom}}(\gamma)=0\}.
\]
Thus, for a test affine \(T\),
\[
\operatorname{Ran}^{\mathrm{loc}}_{R,\mathrm{pre}}(E)(T)
=
\operatorname*{colim}_{(I,\alpha)}
\operatorname{Map}(T,E)^I,
\qquad
\alpha:I\to\Gamma_R^{\mathrm{loc}},
\]
where the colimit is over finite sets, relabellings, and collision
maps.  The forgetful map to the ordinary Ran prestack discards the
colour map \(\alpha\).  This definition contains no wrapped family and
no anchor map.  The local closed-configuration incidence stacks
\(\mathfrak M_{\alpha,R,I}^{\mathrm{loc,cl}}\) and the sheaves
\(\mathcal F_{\alpha,R}^{\mathrm{loc}}\) are coefficient data over this
carrier, not the carrier itself.
\end{definition}

The packet
\[
\texttt{certificates/hybrid/local\_stratum\_b0\_prestack\_definition}
\]
verifies
\(\texttt{LOCAL\_STRATUM\_B0\_PRESTACK\_DEFINITION\_VERIFIED}\): the
\(b=0\) local stratum is defined as a subprestack of the hybrid
prestack, while component rows, closed-configuration rows, local sheaf
rows, descent rows, and transition rows are recorded as missing open
obligations.

\begin{definition}[Wrapped \(b>0\) stratum as a subprestack]
\label{def:hybrid-wrapped-stratum-prestack}
At finite HN height \(R\), the wrapped stratum is the full subprestack
\[
\operatorname{Ran}^{\mathrm{wr}}_{R,\mathrm{pre}}(E)
\subset
\operatorname{Ran}^{\mathrm{hyb}}_{R,\mathrm{pre}}(E)
\]
cut out by the conditions \(I=\varnothing\) and
\(\eta:J\to\Gamma_R^{\mathrm{wr}}\), where
\[
\Gamma_R^{\mathrm{wr}}=\{\gamma\in\Gamma_{\sigma,S}^{\mathrm{HN}}(R)
\mid b_R^{\mathrm{geom}}(\gamma)>0\}.
\]
Thus, for a test affine \(T\),
\[
\operatorname{Ran}^{\mathrm{wr}}_{R,\mathrm{pre}}(E)(T)
=
\operatorname*{colim}_{(J,\eta)}
\prod_{j\in J}
\mathcal M_{\eta_j,R}^{\mathrm{wr,rig}}(T),
\qquad
\eta:J\to\Gamma_R^{\mathrm{wr}},
\]
where the colimit is over finite sets and relabellings.  The anchor
maps
\[
\lambda_{\eta_j,R}(W_j):T\longrightarrow E
\]
are retained as part of each wrapped prestack point.  This definition
contains no finite local support maps \(x_i:T\to E\) and no local
colour labels.  The wrapped coefficient objects
\(\mathcal G_{\eta,R}^{\mathrm{wr}}\), anchor residuals, and
quotient-after-correspondence data are coefficient and correspondence
data over this carrier, not the carrier itself.
\end{definition}

The packet
\[
\texttt{certificates/hybrid/wrapped\_stratum\_bpositive\_prestack\_definition}
\]
verifies
\(\texttt{WRAPPED\_STRATUM\_BPOSITIVE\_PRESTACK\_DEFINITION\_VERIFIED}\):
the \(b>0\) wrapped stratum is defined as a subprestack of the hybrid
prestack, while component rows, populated wrapped prequotient rows,
anchor rows, support-realization rows, descent rows, and transition
rows are recorded as missing open obligations.

\begin{definition}[\(E\)-equivariant wrapped prequotient stacks]
\label{def:e-equivariant-wrapped-prequotient-stack}
Fix a finite HN height \(R\) and a wrapped colour
\(\eta\in\Gamma_R^{\mathrm{wr}}\).  The wrapped prequotient row is the
scalar-rigidified retained wrapped stack
\[
\mathcal M_{\eta,R}^{\mathrm{wr,rig}}
\subset \mathfrak M_{\eta,R}^{\mathrm{rig}}
\]
whose objects have geometric elliptic degree
\[
b_R^{\mathrm{geom}}(\eta)>0.
\]
The superscript \(\mathrm{rig}\) means that the scalar automorphisms
have been rigidified.  It does not mean that the elliptic
translation direction has been divided out.  The map
\[
\rho_{\eta,R}:
\mathcal M_{\eta,R}^{\mathrm{wr,rig}}
\longrightarrow
\mathcal M_{\eta,R}^{\mathrm{wr}}\subset\mathfrak M_{\eta,R}
\]
forgets this scalar rigidification and the finite wrapped bookkeeping.
The reduced quotient, when later used, is written
\[
\mathcal M_{\eta,R}^{\mathrm{wr},E}
:=
\bigl[\mathcal M_{\eta,R}^{\mathrm{wr,rig}}/E\bigr],
\]
and belongs to the quotient-descent rows, not to the prequotient
definition.

For a test affine \(T\) and a section \(a:T\to E\), let
\(\tau_a:E\times T\to E\times T\) be relative translation by \(a\).
An object \(A\in\mathcal M_{\eta,R}^{\mathrm{wr,rig}}(T)\), represented
by the universal perfect complex on \(K3\times E\times T\), is acted on
by
\[
a\cdot A
:=
\bigl(1_{K3}\times\tau_a\bigr)_*A .
\]
The scalar rigidification and the universal complex are transported by
this automorphism.  The origin \(0_E\) is fixed once and for all and
acts as the identity.  This defines the \(E\)-equivariant wrapped
prequotient stack.  It does not prove finite type, properness, final
anchor population, anchor losslessness, quotient descent, or transition
compatibility.
\end{definition}

\begin{lemma}[\(E\)-action law on the wrapped prequotient; \textit{proved}]
\label{lem:e-action-law-wrapped-prequotient}
The operation of
Definition~\ref{def:e-equivariant-wrapped-prequotient-stack} is an
\(E\)-action on \(\mathcal M_{\eta,R}^{\mathrm{wr,rig}}\).  It preserves
the wrapped colour \(\eta\) and the condition
\(b_R^{\mathrm{geom}}(\eta)>0\).
\end{lemma}

\begin{proof}
Relative translations satisfy
\(\tau_{0_E}=\mathrm{id}\) and
\(\tau_a\circ\tau_b=\tau_{a+b}\).  Functoriality of pushforward along
automorphisms gives
\[
0_E\cdot A=A,\qquad
a\cdot(b\cdot A)=(a+b)\cdot A .
\]
Translation on the elliptic factor acts trivially on the numerical
Mukai class and carries effective one-cycles of a fixed
\(p_E\)-degree to effective one-cycles of the same \(p_E\)-degree.
Hence the retained HN colour and
\(b_R^{\mathrm{geom}}\) are unchanged.  The freeness and zero
translation-rigidification defect on retained rows are exactly the
input of Proposition~\ref{prop:retained-e-translation-rigidifications};
they are imported here only for the \(E\)-action row and do not supply
finite type, properness, quotient descent, anchors, or transitions.
\end{proof}

The packet
\[
\texttt{certificates/hybrid/equivariant\_wrapped\_prequotient\_definition}
\]
verifies
\(\texttt{E\_EQUIVARIANT\_WRAPPED\_PREQUOTIENT\_DEFINITION\_VERIFIED}\):
the scalar-rigidified wrapped prequotient
\(\mathcal M_{\eta,R}^{\mathrm{wr,rig}}\) is defined before the
reduced \(E\)-quotient, the translation action with origin \(0_E\) is
defined, and the action law and wrapped-colour preservation have zero
defect.  Finite type, properness, final anchor population,
anchor losslessness, quotient descent, transition compatibility, and
aggregate hybrid-carrier population are recorded as missing open
obligations.

\begin{proposition}[Finite type of wrapped prequotients]
\label{prop:wrapped-prequotient-finite-type}
\textnormal{[proved at finite HN height]}
Assume the retained finite class-bound datum of
Definition~\ref{def:retained-finite-class-bound-datum}, the retained
finite-type semistable substack datum of
Definition~\ref{def:retained-finite-type-semistable-substack-datum},
and the scalar-rigidification datum of
Definition~\ref{def:retained-rigidification-inertia-datum}.  Then, for
every wrapped colour \(\eta\in\Gamma_R^{\mathrm{wr}}\), the
prequotient stack
\[
\mathcal M_{\eta,R}^{\mathrm{wr,rig}}
\]
of Definition~\ref{def:e-equivariant-wrapped-prequotient-stack} is of
finite type over \(\C\).  This is a finite-stage statement before the
reduced \(E\)-quotient is taken.
\end{proposition}

\begin{proof}
Let
\[
C_R(\eta)=\{c\in C_R\mid \widehat c=\eta\}.
\]
Proposition~\ref{prop:retained-finite-class-bounds} makes \(C_R\)
finite; hence \(C_R(\eta)\) is finite.  For each
\(c\in C_R(\eta)\),
Proposition~\ref{prop:retained-finite-type-semistable-substacks} gives
a finite-type retained semistable stack \(\mathfrak M^{ss}_{R,c}\).
The scalar subgroup \(\mathbb G_m\) is flat, finitely presented, and
central in the retained stabilizer sequence.  The rigidification
\[
\mathfrak M^{ss}_{R,c}\longrightarrow
\mathfrak M^{\mathrm{rig}}_{R,c}
=\mathfrak M^{ss}_{R,c}/\mathbb G_m
\]
is therefore an fppf gerbe by the rigidification theorem
\cite[Tag 04V2]{StacksProject}; finite type descends across such an
fppf cover by fpqc descent for finite type
\cite[Tag 041K, Lemma 74.11.11]{StacksProject}.  Hence every
\(\mathfrak M^{\mathrm{rig}}_{R,c}\) with \(c\in C_R(\eta)\) is finite
type.  The wrapped prequotient is the finite union of these
scalar-rigidified retained class pieces:
\[
\mathcal M_{\eta,R}^{\mathrm{wr,rig}}
=
\bigcup_{c\in C_R(\eta)}\mathfrak M^{\mathrm{rig}}_{R,c}.
\]
Finite unions of finite-type substacks are finite type.  The condition
\(b_R^{\mathrm{geom}}(\eta)>0\) is numerical on the retained class and
does not introduce an infinite support parameter.  The proof does not
establish properness, quotient descent, final anchor population,
anchor losslessness, transition compatibility, or aggregate
hybrid-carrier population.
\end{proof}

The packet
\[
\texttt{certificates/hybrid/wrapped\_prequotient\_finite\_type}
\]
verifies
\(\texttt{WRAPPED\_PREQUOTIENT\_FINITE\_TYPE\_VERIFIED}\): finite type
of \(\mathcal M_{\eta,R}^{\mathrm{wr,rig}}\) follows from the finite
retained class set, finite-type retained semistable substacks, and
scalar rigidification by \(\mathbb G_m\).  Properness, final anchor
population, anchor losslessness, quotient descent, transition
compatibility, and aggregate hybrid-carrier population are still
missing open obligations.

\begin{corollary}[Finite-support Ran locality misses the geometric elliptic-degree
direction; \textit{proved}]
\label{cor:finite-support-ranE-misses-s}
A support-local factorization algebra on $\operatorname{Ran}(E)$,
defined only by separating supports over disjoint opens of $E$,
sees the projection-finite sector.  It does not realize the sectors of
positive geometric elliptic degree.
\end{corollary}

\begin{proof}
The obstruction uses the geometric degree
$b_R^{\mathrm{geom}}=\deg_E$, not the Borcherds coordinate
$m=\operatorname{pr}_3\overline\Pi_X$.  On the rank-one
Oberdieck--Pandharipande/Donaldson--Thomas branch one has the
branch-local equality
$m_{\mathrm{Bch}}=d_E-1$ \cite{OberdieckReducedPT}; it is not a
definition of the wrapped stratum.  By
Proposition~\ref{prop:positive-elliptic-degree-projects}, an effective
support cycle with \(b_R^{\mathrm{geom}}>0\) projects onto all of
\(E\).  The object-level statement requires the support-realization row
named in that proposition.  Such wrapped supports cannot be separated
by disjoint opens.  Therefore ordinary Ran-locality has no locality
mechanism for the wrapped sector.
\end{proof}

\begin{definition}[Geometric elliptic-degree map]
\label{def:geometric-elliptic-degree-map}
Fix a finite HN height \(R\) and a retained HN numerical class
\(\gamma\in\Gamma_{\sigma,S}^{\mathrm{HN}}(R)\).  Let
\[
\operatorname{cyc}_1(\gamma)\in N_1(K3\times E)
\]
be the numerical one-cycle obtained from the codimension-two Chern
character component, equivalently the class of
\(\operatorname{ch}_2(A)\cap[K3\times E]\) for any retained
semistable object \(A\) of numerical class \(\gamma\).  Since
\(N_1(E)=\Z[E]\), there is a unique integer
\(b_R^{\mathrm{geom}}(\gamma)\) such that
\[
(p_E)_*\operatorname{cyc}_1(\gamma)
=b_R^{\mathrm{geom}}(\gamma)[E].
\]
On the effective retained one-dimensional sector this integer is
nonnegative.  The local and wrapped charge windows are
\[
\Gamma_R^{\mathrm{loc}}
=\{\gamma\mid b_R^{\mathrm{geom}}(\gamma)=0\},\qquad
\Gamma_R^{\mathrm{wr}}
=\{\gamma\mid b_R^{\mathrm{geom}}(\gamma)>0\}.
\]
The map \(b_R^{\mathrm{geom}}\) is not the Borcherds trace coordinate
\[
m_R^{\mathrm{Bch}}=\operatorname{pr}_3\overline\Pi_X .
\]
The latter is a normal-ordered Gram coordinate; the former is the
proper-pushforward degree of the geometric support along
\(p_E:K3\times E\to E\).
\end{definition}

The packet
\[
\texttt{certificates/hybrid/geometric\_elliptic\_degree\_map}
\]
verifies
\(\texttt{GEOMETRIC\_ELLIPTIC\_DEGREE\_MAP\_OBSTRUCTION\_VERIFIED}\):
the definition is fixed, while the retained HN domain rows, one-cycle
class rows, class-constancy rows, pushforward rows, integer degree rows,
nonnegativity rows, local/wrapped split rows, transition rows, and
retained-extension additivity rows are not supplied.  The last item is
the separate row-202 packet below.

\begin{definition}[Geometric elliptic-degree additivity datum]
\label{def:geometric-elliptic-degree-additivity-datum}
At finite HN height \(R\), a geometric elliptic-degree additivity
datum consists, for every retained extension row
\[
0\longrightarrow A_{\gamma'}\longrightarrow A_\gamma
\longrightarrow A_{\gamma''}\longrightarrow0
\]
in the finite Hall stage, of the following rows:
\begin{enumerate}
\item the retained extension-pair row, with
\(\gamma',\gamma'',\gamma\in
\Gamma_{\sigma,S}^{\mathrm{HN}}(R)\);
\item the one-cycle equality
\[
\operatorname{cyc}_1(\gamma)
=\operatorname{cyc}_1(\gamma')
+\operatorname{cyc}_1(\gamma'')
\quad\text{in }N_1(K3\times E);
\]
\item the pushed-forward equality
\[
(p_E)_*\operatorname{cyc}_1(\gamma)
=(p_E)_*\operatorname{cyc}_1(\gamma')
+(p_E)_*\operatorname{cyc}_1(\gamma'')
\quad\text{in }N_1(E);
\]
\item the integer degree equality
\[
b_R^{\mathrm{geom}}(\gamma)
=b_R^{\mathrm{geom}}(\gamma')
+b_R^{\mathrm{geom}}(\gamma'');
\]
\item compatibility of these rows with the retained transition maps
\(R\le R'\).
\end{enumerate}
The same definition applies to a distinguished triangle representing
the Hall extension correspondence, with the corresponding equality in
\(K_0(K3\times E)\).
\end{definition}

\begin{proposition}[Additivity criterion for the geometric elliptic degree]
\label{prop:geometric-elliptic-degree-additivity-criterion}
Assume a finite stage carries the datum of
Definition~\ref{def:geometric-elliptic-degree-additivity-datum}.  Then
\[
b_R^{\mathrm{geom}}(\gamma)
=b_R^{\mathrm{geom}}(\gamma')
+b_R^{\mathrm{geom}}(\gamma'')
\]
for every retained extension row
\[
0\to A_{\gamma'}\to A_\gamma\to A_{\gamma''}\to0.
\]
\end{proposition}

\begin{proof}
In \(K_0(K3\times E)\) the retained extension gives
\([A_\gamma]=[A_{\gamma'}]+[A_{\gamma''}]\).  The Chern character is
additive on \(K_0\), hence the codimension-two component gives the
one-cycle equality in \(N_1(K3\times E)\).  Proper pushforward along
\(p_E\) is linear on numerical one-cycles, so the equality remains true
in \(N_1(E)\).  Since \(N_1(E)=\Z[E]\), equality of the pushed-forward
cycles is equality of the coefficients of \([E]\), which are precisely
the integers \(b_R^{\mathrm{geom}}\) of
Definition~\ref{def:geometric-elliptic-degree-map}.  The transition
rows state that this coefficient equality is preserved under
\(R\le R'\).
\end{proof}

The packet
\[
\texttt{certificates/hybrid/geometric\_elliptic\_degree\_additivity}
\]
verifies
\(\texttt{GEOMETRIC\_ELLIPTIC\_DEGREE\_ADDITIVITY\_OBSTRUCTION\_VERIFIED}\):
the additivity criterion is fixed, the degree-map definition packet and
the retained-extension-closure packet are imported, and the missing
retained extension-pair rows, one-cycle additivity rows, pushforward
additivity rows, integer degree additivity rows, and transition rows are
recorded fail-closed.  Retained word closure by itself proves only that
the middle HN word is retained; it does not prove additivity of
\((p_E)_*\operatorname{cyc}_1\).

\begin{definition}[Geometric/Borcherds degree comparison datum]
\label{def:geometric-borcherds-degree-comparison-datum}
At finite HN height \(R\), set
\[
m_R^{\mathrm{Bch}}(\widehat\gamma)
:=\operatorname{pr}_3\overline\Pi_X(\widehat\gamma),
\qquad
\overline\Pi_X(\widehat\gamma)=(n,l,m).
\]
A comparison datum between \(b_R^{\mathrm{geom}}\) and
\(m_R^{\mathrm{Bch}}\) consists of:
\begin{enumerate}
\item a specified branch of retained HN classes;
\item a lift of each retained class on that branch to a charge
\(\widehat\gamma\) in the domain of \(\overline\Pi_X\);
\item populated geometric degree rows computing
\(b_R^{\mathrm{geom}}\) from
\((p_E)_*\operatorname{cyc}_1\);
\item populated Gram-coordinate rows computing
\(\operatorname{pr}_3\overline\Pi_X(\widehat\gamma)\);
\item a declared equality or affine comparison formula with zero defect
on the branch;
\item compatibility of the comparison with retained transitions
\(R\le R'\).
\end{enumerate}
\end{definition}

\begin{proposition}[Geometric/Borcherds degree separation]
\label{prop:geometric-borcherds-degree-separation}
The finite hybrid datum distinguishes the support degree
\[
b_R^{\mathrm{geom}}(\gamma)
=\operatorname{coeff}_{[E]}(p_E)_*\operatorname{cyc}_1(\gamma)
\]
from the Borcherds trace coordinate
\[
m_R^{\mathrm{Bch}}(\widehat\gamma)
=\operatorname{pr}_3\overline\Pi_X(\widehat\gamma).
\]
The local/wrapped split is made by \(b_R^{\mathrm{geom}}\).  The
protected trace is graded by \(m_R^{\mathrm{Bch}}\).  No formula
identifying the two coordinates, or replacing one by the other, is part
of the hybrid datum unless a comparison datum of
Definition~\ref{def:geometric-borcherds-degree-comparison-datum} is
supplied.
\end{proposition}

\begin{proof}
The two maps have different constructions.  The map
\(b_R^{\mathrm{geom}}\) is obtained from the numerical one-cycle class
of a retained object by proper pushforward to \(N_1(E)=\Z[E]\); it is
defined before the local/wrapped factorization carrier is formed.  The
coordinate \(m_R^{\mathrm{Bch}}\) is the third coordinate of the
normal-ordered Gram homomorphism \(\overline\Pi_X\); it is used after
the Gram descent to grade protected integration in the Borcherds
variables.  A branch-local identity, such as the rank-one
Oberdieck--Pandharipande comparison \(m_{\mathrm{Bch}}=d_E-1\)
\cite{OberdieckReducedPT}, becomes a theorem only after the branch,
lift, geometric degree rows,
Gram-coordinate rows, and transition compatibility are supplied.  In
the absence of these rows there is no typed morphism that turns
\(\operatorname{pr}_3\overline\Pi_X\) into
\(\operatorname{coeff}_{[E]}(p_E)_*\operatorname{cyc}_1\).
\end{proof}

The packet
\[
\texttt{certificates/hybrid/geometric\_borcherds\_degree\_separation}
\]
verifies
\(\texttt{GEOMETRIC\_BORCHERDS\_DEGREE\_SEPARATION\_OBSTRUCTION\_VERIFIED}\):
the two coordinates are separated, substitution of
\(\operatorname{pr}_3\overline\Pi_X\) for \(b_R^{\mathrm{geom}}\) is
excluded, and branch-comparison rows are recorded as missing open
obligations.

\begin{proposition}[Middle type-II wall routes through wrapped and mixed geometry]
\label{prop:middle-type-ii-wall-wrapped-mixed}
\textnormal{[routing proved; retained wall object not constructed]}
Let
\[
\gamma_{\mathrm{wall}}(\delta_2)=(0,1,1)
\]
be the middle simple type-II wall image of
Proposition~\ref{prop:type-ii-wall-images}.  On the type-II wall branch
where the geometric elliptic degree is written in the OP/Borcherds
coordinate as \(d=m+1\), this wall has \(d=2>0\).  Consequently, any
retained finite-stage class \(\eta\) with a supplied branch-comparison
row
\[
\overline\Pi_X(\eta)=(0,1,1),\qquad
b_R^{\mathrm{geom}}(\eta)=d=2,
\]
lies in \(\Gamma_R^{\mathrm{wr}}\), not in \(\Gamma_R^{\mathrm{loc}}\).
Its one-wall carrier is the wrapped prequotient side
\(\mathcal M_{\eta,R}^{\mathrm{wr,rig}}\) of
Definition~\ref{def:e-equivariant-wrapped-prequotient-stack}.  Any Hall
operation inserting projection-finite local colours against this wall
is typed by the one-sided \(LW/WL\) mixed correspondences of
Definition~\ref{def:mixed-local-wrapped-correspondence-datum}, and by
the ordered two-sided mixed correspondence of
Definition~\ref{def:two-sided-mixed-correspondence-datum} when local
insertions occur on both sides.  It is not an \(LL\) operation on the
ordinary Ran stratum.

This is a routing theorem.  It does not supply the retained wall object
\(x_{\delta_2,R}\), an effective support-realization row, a populated
wrapped anchor, quotient descent, the O2 wall atlas, or protected
integration.
\end{proposition}

\begin{proof}
Proposition~\ref{prop:type-ii-wall-images} gives the middle wall
triple \((n,l,m)=(0,1,1)\).  The wall-branch convention
\(d=m+1\) gives \(d=2\).  After the retained branch-comparison row
identifies this integer with \(b_R^{\mathrm{geom}}(\eta)\), the
definition
\[
\Gamma_R^{\mathrm{loc}}=\{b_R^{\mathrm{geom}}=0\},\qquad
\Gamma_R^{\mathrm{wr}}=\{b_R^{\mathrm{geom}}>0\}
\]
puts \(\eta\) in the wrapped window and excludes it from the local
window.  By Proposition~\ref{prop:positive-elliptic-degree-projects},
if a retained object realizing \(\eta\) is supplied together with its
effective support-cycle row, its one-dimensional support projects onto
all of \(E\), so finite-support Ran locality cannot contain it.
Definition~\ref{def:hybrid-wrapped-stratum-prestack} therefore places
the one-wall carrier in the wrapped \(b>0\) subprestack, with the
prequotient stack of
Definition~\ref{def:e-equivariant-wrapped-prequotient-stack}.  The
local/wrapped and two-sided mixed correspondence definitions are exactly
the before-\(E\)-quotient operations whose source contains one wrapped
input and projection-finite local inputs.  The final sentence records
the rows still absent from this routing calculation.
\end{proof}

The packet
\[
\texttt{certificates/hybrid/middle\_type\_ii\_wall\_wrapped\_mixed}
\]
verifies
\(\texttt{MIDDLE\_TYPE\_II\_WALL\_WRAPPED\_MIXED\_ROUTING\_VERIFIED}\):
the middle wall image has \(m=1\), the wall-branch degree \(d=m+1=2\)
is positive, the routing is to the wrapped \(b>0\) side, and local
insertions are typed by mixed correspondences.  Its obstruction rows
record that retained branch-comparison, wall-object, support,
anchor, O2 atlas, quotient, protected-integration, aggregate-carrier,
and transition rows are still missing.

\begin{proposition}[Hybrid operation profiles recover the three simple type-II walls]
\label{prop:type-ii-walls-hybrid-operation-recovery}
\textnormal{[operation-profile recovery proved; carrier population not constructed]}
Let
\[
\gamma_{\mathrm{wall}}(\delta_1)=(1,1,0),\qquad
\gamma_{\mathrm{wall}}(\delta_2)=(0,1,1),\qquad
\gamma_{\mathrm{wall}}(\delta_3)=(0,-1,0)
\]
be the three simple type-II wall images of
Proposition~\ref{prop:type-ii-wall-images}.  On the same type-II wall
branch \(d=m+1\), their elliptic degrees are
\[
d_{\delta_1}=1,\qquad d_{\delta_2}=2,\qquad d_{\delta_3}=1.
\]
Thus, after retained branch-comparison rows identify
\[
b_R^{\mathrm{geom}}(\eta_i)=d_{\delta_i},\qquad
\overline\Pi_X(\eta_i)=\gamma_{\mathrm{wall}}(\delta_i),
\]
the three simple type-II walls are recovered by the hybrid vocabulary
as the three singleton wrapped colour profiles
\[
\eta_i\in\Gamma_R^{\mathrm{wr}},\qquad i=1,2,3.
\]
For each \(i\), a finite local insertion on the left or right of this
wall is typed by the \(LW\) or \(WL\) mixed correspondence, and local
insertions on both sides are typed by the ordered \(LWL\) correspondence.
No \(LL\)-only local operation recovers a type-II wall profile.

This is recovery of operation \emph{types}.  It does not populate the
wrapped prequotient stacks, mixed extension stacks, anchors, quotient
descent, O2 wall atlas, protected integration, or transition system.
\end{proposition}

\begin{proof}
The triples are those of Proposition~\ref{prop:type-ii-wall-images}.
Applying the wall-branch formula \(d=m+1\) gives
\[
(1,1,0)\mapsto1,\qquad
(0,1,1)\mapsto2,\qquad
(0,-1,0)\mapsto1.
\]
All three integers are positive.  Once a retained branch-comparison row
identifies these integers with \(b_R^{\mathrm{geom}}\), the definition
\[
\Gamma_R^{\mathrm{wr}}=\{b_R^{\mathrm{geom}}>0\},\qquad
\Gamma_R^{\mathrm{loc}}=\{b_R^{\mathrm{geom}}=0\}
\]
places all three \(\eta_i\) in the wrapped window and none in the local
window.  Definition~\ref{def:hybrid-wrapped-stratum-prestack} supplies
the singleton wrapped carrier profile.  Definitions
\ref{def:mixed-local-wrapped-correspondence-datum} and
\ref{def:two-sided-mixed-correspondence-datum} are exactly the before
\(E\)-quotient operation profiles with one wrapped input and local
inputs on one or both sides.  Since an \(LL\) correspondence has only
local inputs and local target colour, it cannot have target colour
\(\eta_i\in\Gamma_R^{\mathrm{wr}}\).  The final sentence separates
operation-profile recovery from the missing populated finite-stage
rows.
\end{proof}

The packet
\[
\texttt{certificates/hybrid/type\_ii\_wall\_hybrid\_operation\_recovery}
\]
verifies
\(\texttt{TYPE\_II\_WALL\_HYBRID\_OPERATION\_RECOVERY\_VERIFIED}\):
the three imported wall triples are distinct, have positive
wall-branch degrees \(1,2,1\), are routed to singleton wrapped colour
profiles, and carry exactly the \(LW\), \(WL\), and \(LWL\) mixed
operation profiles.  Its obstruction rows record that retained
branch-comparison, wall-object, support, anchor, mixed-stack, O2 atlas,
\(E\)-quotient losslessness, quotient, protected-integration,
aggregate-carrier, and transition rows are still missing.

\begin{proposition}[Type-II wall profiles survive the \(E\)-quotient assignment]
\label{prop:type-ii-walls-survive-e-quotient}
\textnormal{[profile-level survival proved; object-level nonvanishing not proved]}
Let \(\eta_i\), \(i=1,2,3\), be the singleton wrapped type-II wall
profiles of Proposition~\ref{prop:type-ii-walls-hybrid-operation-recovery},
and assume the finite quotient-after-correspondence datum of
Definition~\ref{def:quotient-after-correspondence-pseudofunctor}.  Then
the quotient assignment sends each retained unreduced wrapped profile to
\[
\mathcal Q^{\mathrm{corr}}_{E,R}
\bigl(\mathcal M_{\eta_i,R}^{\mathrm{wr,rig}}\bigr)
=
\bigl[\mathcal M_{\eta_i,R}^{\mathrm{wr,rig}}/E\bigr].
\]
For every retained finite local insertion on one or both sides of
\(\eta_i\), the associated \(LW\), \(WL\), or \(LWL\) operation profile is
sent to the quotient of the already formed unreduced correspondence
stack.  Hence none of the three simple type-II wall profiles is deleted
by the reduced \(E\)-quotient at the quotient-after-correspondence
assignment level.

This does not prove that
\(\mathcal M_{\eta_i,R}^{\mathrm{wr,rig}}\) is nonempty, that its quotient
cycle has nonzero Borel--Moore class, that vanishing-cycle coefficients
descend nontrivially, or that the reduced object enters the O2 wall
atlas, protected Pfaffian, aggregate hybrid carrier, or transition
system.
\end{proposition}

\begin{proof}
By Proposition~\ref{prop:type-ii-walls-hybrid-operation-recovery}, the
three simple type-II wall images have positive wall-branch elliptic
degrees \(1,2,1\).  After the retained branch-comparison rows they are
therefore the singleton wrapped colours \(\eta_i\in\Gamma_R^{\mathrm{wr}}\).
The same proposition identifies the finite local insertions on these
wall profiles with the one-sided \(LW\) and \(WL\) mixed operation
profiles and the two-sided \(LWL\) operation profile, all before the
reduced \(E\)-quotient.

Definition~\ref{def:quotient-after-correspondence-pseudofunctor}
assigns \([Y/E]\) to every retained object \(Y\) of
\(\mathsf{Corr}^{E,\mathrm{hyb}}_R\), including wrapped rigidified
prequotients.  Taking \(Y=\mathcal M_{\eta_i,R}^{\mathrm{wr,rig}}\)
gives the displayed quotient object for each \(i\).  The same definition
assigns to every generating unreduced correspondence
\[
e=(Y_0\xleftarrow{p_e}\mathfrak E_e\xrightarrow{q_e}Y_1)
\]
the already-formed quotient correspondence
\[
\overline e=
(\overline Y_0\xleftarrow{\overline p_e}
[\mathfrak E_e/E]\xrightarrow{\overline q_e}\overline Y_1).
\]
Applying this to the \(LW\), \(WL\), and \(LWL\) operation profiles
therefore preserves the operation labels as quotient-after-correspondence
arrows.  Proposition~\ref{prop:quotient-after-correspondence-excludes-quotient-first}
records that this is not a quotient-first reconstruction from reduced
object stacks, and
Proposition~\ref{prop:quotient-after-correspondence-hn-transition-compatibility}
keeps any supplied transition row in the same quotient-presented form.
The final paragraph states the unsupplied rows and prevents this
profile-level assignment theorem from being read as object-level
nonvanishing, O2 atlas construction, or protected integration.
\end{proof}

The packet
\[
\texttt{certificates/hybrid/type\_ii\_wall\_e\_quotient\_survival}
\]
verifies
\(\texttt{TYPE\_II\_WALL\_E\_QUOTIENT\_SURVIVAL\_VERIFIED}\): the three
wrapped type-II wall profiles are assigned quotient objects, their nine
\(LW\), \(WL\), and \(LWL\) operation profiles are assigned
quotient-after-correspondence arrows, quotient-first reconstruction is
excluded, and HN transition compatibility is imported.  Its firewall rows
record that object-level nonzero survival, coefficient descent,
orientation descent, O2 wall-atlas insertion, protected integration, and
aggregate hybrid-carrier population remain open.

\begin{proposition}[No type-I wall enters the type-II hybrid wall inlet]
\label{prop:no-type-i-wall-enters-hybrid-carrier}
\textnormal{[profile-level exclusion proved; aggregate carrier not populated]}
Let
\[
\vartheta_{12}=f_{-2}-f_2,\qquad
\vartheta_{13}=f_2-f_3,\qquad
\vartheta_{23}=f_{-2}-f_3
\]
be the three type-I chamber-automorphism roots of
Lemma~\ref{lem:bkm-chamber-s3}.  At finite HN height \(R\), let the
type-II wall-profile inlet be the finite set of retained profiles
\[
\Gamma_R^{\mathrm{wall},II}
=
\left\{\eta_i\;\middle|\;
\overline\Pi_X(\eta_i)=\gamma_{\mathrm{wall}}(\delta_i),
b_R^{\mathrm{geom}}(\eta_i)=d_{\delta_i},
i=1,2,3
\right\}
\]
appearing in Propositions
\ref{prop:type-ii-walls-hybrid-operation-recovery} and
\ref{prop:type-ii-walls-survive-e-quotient}.  Then every element of
\(\Gamma_R^{\mathrm{wall},II}\) is labelled by one of the three
type-II divisor roots \(\delta_1,\delta_2,\delta_3\).  No type-I root
\(\vartheta_{12},\vartheta_{13},\vartheta_{23}\) supplies a wall profile
in \(\Gamma_R^{\mathrm{wall},II}\), and none is created by the
quotient-after-correspondence assignment.

This is an exclusion theorem for the wall-profile inlet.  It is not a
global exhaustion of all retained hybrid colours, not a population
theorem for the aggregate hybrid carrier, and not an O2 wall-atlas
type-I exclusion row.
\end{proposition}

\begin{proof}
Remark~\ref{rem:bkm-divisibility-types} and the finite
\(\texttt{type\_ii\_chamber\_isotropic}\) packet record the two
square-two divisibility classes.  The roots
\(\delta_1,\delta_2,\delta_3\) have ambient divisibility \(2\), lie in
the type-II sublattice, and occur in \(\operatorname{div}(\Delta_5)\)
with order \(1\).  The roots
\(\vartheta_{12},\vartheta_{13},\vartheta_{23}\) have ambient
divisibility \(1\), generate the \(S_3\) chamber automorphism action of
Lemma~\ref{lem:bkm-chamber-s3}, and have zero
\(\Delta_5\)-divisor order.

Proposition~\ref{prop:type-ii-walls-hybrid-operation-recovery} imports
only the three type-II wall images
\[
\gamma_{\mathrm{wall}}(\delta_1),\qquad
\gamma_{\mathrm{wall}}(\delta_2),\qquad
\gamma_{\mathrm{wall}}(\delta_3)
\]
and routes them to singleton wrapped profiles with positive
wall-branch degrees \(1,2,1\).  Hence the wall-profile inlet just
defined is indexed by the divisor-supported type-II roots and by no
other square-two root.  The type-I reflections act only by permuting
the three \(\delta_i\); they do not add a fourth wall label, and they do
not turn a zero-order type-I chamber wall into a
\(\Delta_5\)-divisor wall.

Proposition~\ref{prop:type-ii-walls-survive-e-quotient} then applies
\(\mathcal Q^{\mathrm{corr}}_{E,R}\) to the same three already formed
wall profiles and to their \(LW\), \(WL\), and \(LWL\) operation
profiles.  Since quotient-after-correspondence assigns quotient objects
and quotient arrows to retained inputs, it preserves this three-element
indexing set and introduces no new type-I profile.  The final paragraph
separates this profile-level firewall from the unsupplied global carrier
and O2-atlas rows.
\end{proof}

The packet
\[
\texttt{certificates/hybrid/type\_i\_false\_wall\_exclusion}
\]
verifies
\(\texttt{TYPE\_I\_FALSE\_WALL\_EXCLUSION\_VERIFIED}\): the admitted
wall-profile rows are exactly the three type-II profiles
\(\delta_1,\delta_2,\delta_3\), the three type-I chamber roots have
divisibility \(1\), zero \(\Delta_5\)-divisor support, and zero
\(\Delta_5\)-divisor order, and the \(E\)-quotient survival row does
not resurrect a type-I wall profile.  Its firewall rows record that
global colour exhaustion, aggregate hybrid-carrier population, retained
wall objects, coefficient descent, O2 wall-atlas insertion, protected
integration, and transition to O2 remain open.

\begin{definition}[Local/local correspondence datum]
\label{def:local-local-correspondence-datum}
Fix a finite HN height \(R\).  Let
\(\alpha:I\to\Gamma_R^{\mathrm{loc}}\) and
\(\beta:J\to\Gamma_R^{\mathrm{loc}}\) be finite local colour maps, and
write \(|\alpha|=\sum_{i\in I}\alpha_i\),
\(|\beta|=\sum_{j\in J}\beta_j\).  A local/local correspondence row is
defined only for a retained target colour
\(\gamma\in\Gamma_R^{\mathrm{loc}}\) satisfying
\[
\gamma=|\alpha|+|\beta|.
\]
The target index set \(K\) is \(I\sqcup J\), modulo the ordinary Ran
collision maps.  The ordered local/local correspondence is the
projection-finite diagram
\[
\mathfrak E^{\mathrm{LL}}_{\alpha,\beta;\gamma,R,I,J,K}
\xrightarrow{p^{\mathrm{LL}}}
\mathfrak M^{\mathrm{loc,cl}}_{\alpha,R,I}
\times
\mathfrak M^{\mathrm{loc,cl}}_{\beta,R,J},
\qquad
\mathfrak E^{\mathrm{LL}}_{\alpha,\beta;\gamma,R,I,J,K}
\xrightarrow{q^{\mathrm{LL}}}
\mathfrak M^{\mathrm{loc,cl}}_{\gamma,R,K},
\]
classifying retained extensions
\[
0\longrightarrow A_{\beta,J}\longrightarrow B_{\gamma,K}
\longrightarrow A_{\alpha,I}\longrightarrow0.
\]
This diagram lives over
\(\operatorname{Ran}^{\mathrm{loc}}_{R,\mathrm{pre}}(E)\times
\operatorname{Ran}^{\mathrm{loc}}_{R,\mathrm{pre}}(E)\) and targets
\(\operatorname{Ran}^{\mathrm{loc}}_{R,\mathrm{pre}}(E)\).  It has no
wrapped input and no reduced \(E\)-quotient.  Symmetric descent,
collision compatibility, properness of the local extension stack,
Thom--Sebastiani transport, and transition compatibility are separate
rows.  The diagram is not a mixed correspondence, not a wrapped/wrapped
correspondence, not a two-sided mixed correspondence, and not an
ordinary uncoloured Ran product.
\end{definition}

The packet
\[
\texttt{certificates/hybrid/local\_local\_correspondence\_definition}
\]
verifies
\(\texttt{LOCAL\_LOCAL\_CORRESPONDENCE\_DEFINITION\_VERIFIED}\):
the ordered \(LL\) correspondence type is defined, while
extension-stack rows, source/target-map rows, descent rows, properness
rows, Thom--Sebastiani rows, collision rows, and transition rows are
not supplied by that definition packet.

\begin{proposition}[Target properness for local/local extensions]
\label{prop:local-local-target-properness}
\textnormal{[proved for the Hall target map]}
Assume the local closed-incidence stacks
\(\mathfrak M^{\mathrm{loc,cl}}_{\alpha,R,I}\),
\(\mathfrak M^{\mathrm{loc,cl}}_{\beta,R,J}\), and
\(\mathfrak M^{\mathrm{loc,cl}}_{\gamma,R,K}\) carry the universal
perfect complexes and the retained extension-closure datum of
Definition~\ref{def:retained-extension-closure-datum}.  Then the
target map
\[
q^{\mathrm{LL}}:
\mathfrak E^{\mathrm{LL}}_{\alpha,\beta;\gamma,R,I,J,K}
\longrightarrow
\mathfrak M^{\mathrm{loc,cl}}_{\gamma,R,K}
\]
is proper.  The source map \(p^{\mathrm{LL}}\) is not asserted to be
proper.
\end{proposition}

\begin{proof}
Let \(B_{\gamma,K}\) be the universal target object on
\(\mathfrak M^{\mathrm{loc,cl}}_{\gamma,R,K}\).  The stack
\(\mathfrak E^{\mathrm{LL}}_{\alpha,\beta;\gamma,R,I,J,K}\) is the
closed substack of the relative Quot scheme over
\(\mathfrak M^{\mathrm{loc,cl}}_{\gamma,R,K}\) parametrising
quotients
\[
B_{\gamma,K}\twoheadrightarrow A_{\alpha,I}
\]
whose kernel is an object \(A_{\beta,J}\), with exact sequence
\[
0\longrightarrow A_{\beta,J}\longrightarrow B_{\gamma,K}
\longrightarrow A_{\alpha,I}\longrightarrow0.
\]
The retained extension-closure datum says that the middle term stays
inside the retained finite HN window.  The local support conditions
over \(E^I\), \(E^J\), and \(E^K\), and the retained numerical colour
conditions \(\alpha,\beta,\gamma\), are closed conditions inside the
relative Quot scheme.  Grothendieck projectivity of Quot
\cite[Theorem~2.2.4]{HuybrechtsLehn} makes the relative Quot scheme
projective over \(\mathfrak M^{\mathrm{loc,cl}}_{\gamma,R,K}\), hence
proper.  A closed substack of a proper stack over the same base is
proper, so \(q^{\mathrm{LL}}\) is proper.

The source map has a different fibre.  Over fixed endpoints
\((A_{\alpha,I},A_{\beta,J})\), the extension parameters contain
\(\operatorname{Ext}^1(A_{\alpha,I},A_{\beta,J})\) modulo the usual
\(\operatorname{Hom}\)-automorphisms.  When this Ext group is nonzero,
the fibre contains affine directions and is not proper in general.
Thus the properness needed for the Hall pushforward is the
target-properness of \(q^{\mathrm{LL}}\), not source-properness of
\(p^{\mathrm{LL}}\).
\end{proof}

The packet
\[
\texttt{certificates/hybrid/local\_local\_extension\_properness}
\]
verifies
\(\texttt{LOCAL\_LOCAL\_TARGET\_PROPERNESS\_VERIFIED}\): the
local/local extension stack is represented as a closed relative Quot
locus over the target local incidence stack, the target map
\(q^{\mathrm{LL}}\) is proper, and the false source-properness
substitute is excluded.  Descent, collision compatibility,
Thom--Sebastiani transport, transition compatibility, and aggregate
hybrid-carrier population remain separate obligations.  The
compact-support exceptional pushforward model is supplied separately by
Definition~\ref{def:compact-support-exceptional-pushforward-model}.

\begin{definition}[One-sided mixed local/wrapped correspondence datum]
\label{def:mixed-local-wrapped-correspondence-datum}
Fix a finite HN height \(R\).  Let
\(\alpha:I\to\Gamma_R^{\mathrm{loc}}\) be a finite local colour map and
write \(|\alpha|=\sum_{i\in I}\alpha_i\).  Let
\(\eta\in\Gamma_R^{\mathrm{wr}}\).  A one-sided mixed correspondence
row is defined only for a retained target colour
\(\zeta\in\Gamma_R^{\mathrm{wr}}\) satisfying the typed equality
\[
\zeta=|\alpha|+\eta
\quad\text{or}\quad
\zeta=\eta+|\alpha|,
\]
with the positive wrapped condition read through \(b_R^{\mathrm{geom}}\).
The two ordered mixed diagrams are the unreduced correspondence stacks
\[
\mathfrak E^{\mathrm{LW}}_{\alpha,\eta;\zeta,R,I}
\xrightarrow{p^{\mathrm{LW}}}
\mathfrak M^{\mathrm{loc,cl}}_{\alpha,R,I}
\times
\mathcal M_{\eta,R}^{\mathrm{wr,rig}},
\qquad
\mathfrak E^{\mathrm{LW}}_{\alpha,\eta;\zeta,R,I}
\xrightarrow{q^{\mathrm{LW}}}
\mathcal M_{\zeta,R}^{\mathrm{wr,rig}},
\]
classifying retained extensions
\[
0\longrightarrow W_\eta\longrightarrow B_\zeta
\longrightarrow A_\alpha\longrightarrow0,
\]
and
\[
\mathfrak E^{\mathrm{WL}}_{\eta,\alpha;\zeta,R,I}
\xrightarrow{p^{\mathrm{WL}}}
\mathcal M_{\eta,R}^{\mathrm{wr,rig}}
\times
\mathfrak M^{\mathrm{loc,cl}}_{\alpha,R,I},
\qquad
\mathfrak E^{\mathrm{WL}}_{\eta,\alpha;\zeta,R,I}
\xrightarrow{q^{\mathrm{WL}}}
\mathcal M_{\zeta,R}^{\mathrm{wr,rig}},
\]
classifying retained extensions
\[
0\longrightarrow A_\alpha\longrightarrow B_\zeta
\longrightarrow W_\eta\longrightarrow0.
\]
The symbols
\(\mathfrak M^{\mathrm{loc,cl}}_{\alpha,R,I}\) and
\(\mathcal M_{\eta,R}^{\mathrm{wr,rig}}\) are the local closed
incidence stack and wrapped rigidified prequotient supplied by the
finite hybrid datum.  These diagrams live before the reduced
\(E\)-quotient.  They are not local-local correspondences, not
wrapped-wrapped correspondences, not ordered two-sided mixed
correspondences, and not a proof of admissibility, base change,
projection formula, Thom--Sebastiani transport, or source/target-anchor
losslessness.
\end{definition}

The packet
\[
\texttt{certificates/hybrid/mixed\_local\_wrapped\_correspondence\_definition}
\]
verifies
\(\texttt{MIXED\_LOCAL\_WRAPPED\_CORRESPONDENCE\_DEFINITION\_VERIFIED}\):
the one-sided \(LW\) and \(WL\) mixed correspondence types are defined,
while extension-stack rows, source/target-map rows, populated
anchor-memory rows, admissibility rows, base-change rows,
projection-formula rows, Thom--Sebastiani rows, and transition rows are
recorded as missing open obligations in that definition packet.

\begin{proposition}[Target admissibility for mixed local/wrapped extensions]
\label{prop:mixed-local-wrapped-target-admissibility}
\textnormal{[proved for the Hall target maps]}
Assume the local closed-incidence stacks
\(\mathfrak M^{\mathrm{loc,cl}}_{\alpha,R,I}\), the wrapped
prequotients \(\mathcal M_{\eta,R}^{\mathrm{wr,rig}}\) and
\(\mathcal M_{\zeta,R}^{\mathrm{wr,rig}}\), the universal perfect
complexes, the retained closed-substack datum, and the retained
extension-closure datum are supplied at the fixed finite HN height
\(R\).  Then the target maps in the one-sided unreduced mixed
correspondences
\[
q^{\mathrm{LW}}:
\mathfrak E^{\mathrm{LW}}_{\alpha,\eta;\zeta,R,I}
\longrightarrow
\mathcal M_{\zeta,R}^{\mathrm{wr,rig}},
\qquad
q^{\mathrm{WL}}:
\mathfrak E^{\mathrm{WL}}_{\eta,\alpha;\zeta,R,I}
\longrightarrow
\mathcal M_{\zeta,R}^{\mathrm{wr,rig}}
\]
are proper.  Hence the target leg needed for the mixed Hall
pushforward is admissible as a proper geometric map before the reduced
\(E\)-quotient.  The source maps \(p^{\mathrm{LW}}\) and
\(p^{\mathrm{WL}}\) are not asserted to be proper.
\end{proposition}

\begin{proof}
Let \(B_\zeta\) be the universal target object on
\(\mathcal M_{\zeta,R}^{\mathrm{wr,rig}}\).  For the \(LW\) order, the
stack
\(\mathfrak E^{\mathrm{LW}}_{\alpha,\eta;\zeta,R,I}\) is the closed
substack of the relative Quot scheme over
\(\mathcal M_{\zeta,R}^{\mathrm{wr,rig}}\) parametrising quotients
\[
B_\zeta\twoheadrightarrow A_{\alpha,I}
\]
whose kernel is an object \(W_\eta\), with exact sequence
\[
0\longrightarrow W_\eta\longrightarrow B_\zeta
\longrightarrow A_{\alpha,I}\longrightarrow0.
\]
The quotient is required to lie in the local closed-incidence stack,
and the kernel is required to lie in the retained wrapped prequotient.
The retained closed-substack datum makes these endpoint conditions
closed on the relative Quot scheme.  The retained colour condition
\(\zeta=|\alpha|+\eta\) and the finite HN window condition are fixed
row conditions; the retained extension-closure datum records that the
allowed middle term remains in the finite window.  Grothendieck
projectivity of Quot \cite[Theorem~2.2.4]{HuybrechtsLehn} makes the
relative Quot scheme projective, hence proper, over
\(\mathcal M_{\zeta,R}^{\mathrm{wr,rig}}\).  A closed substack of this
relative Quot scheme is proper over the same target, so
\(q^{\mathrm{LW}}\) is proper.

For the \(WL\) order one uses the relative Quot scheme over the same
target, but now with quotient
\[
B_\zeta\twoheadrightarrow W_\eta
\]
and kernel \(A_{\alpha,I}\):
\[
0\longrightarrow A_{\alpha,I}\longrightarrow B_\zeta
\longrightarrow W_\eta\longrightarrow0.
\]
The wrapped quotient condition, the local-kernel condition, and the
retained colour condition \(\zeta=\eta+|\alpha|\) are closed row
conditions by the same retained closed-substack and extension-closure
inputs.  The relative Quot scheme is projective over
\(\mathcal M_{\zeta,R}^{\mathrm{wr,rig}}\), and the \(WL\) extension
stack is a closed substack of it.  Hence \(q^{\mathrm{WL}}\) is proper.

The source maps have the usual Hall extension fibres.  Over fixed
endpoints in the \(LW\) order one sees extension classes in
\(\operatorname{Ext}^1(A_{\alpha,I},W_\eta)\) modulo
\(\operatorname{Hom}\)-automorphisms; in the \(WL\) order one sees
\(\operatorname{Ext}^1(W_\eta,A_{\alpha,I})\) modulo the same
automorphism relation.  These fibres contain affine directions when
the corresponding \(\operatorname{Ext}^1\)-group is nonzero, so source
properness is not a consequence of the argument and is not claimed.
The compact-support exceptional pushforward model, base change,
projection formula, Thom--Sebastiani transport, anchor losslessness,
quotient descent, and transition compatibility remain separate rows.
\end{proof}

The packet
\[
\texttt{certificates/hybrid/mixed\_local\_wrapped\_extension\_admissibility}
\]
verifies
\(\texttt{MIXED\_LOCAL\_WRAPPED\_ADMISSIBILITY\_VERIFIED}\): the
one-sided mixed extension stacks are represented as closed relative
Quot loci over the wrapped target prequotient, the target maps
\(q^{\mathrm{LW}}\) and \(q^{\mathrm{WL}}\) are proper, and the false
source-properness substitute is excluded.  The compact-support
exceptional pushforward model is supplied separately by
Definition~\ref{def:compact-support-exceptional-pushforward-model}.
Populated anchor-memory rows, base change, projection formula,
Thom--Sebastiani transport, transition compatibility, quotient descent,
and aggregate hybrid-carrier population remain separate obligations.

\begin{definition}[Wrapped/wrapped correspondence datum]
\label{def:wrapped-wrapped-correspondence-datum}
Fix a finite HN height \(R\).  Let
\(\eta_1,\eta_2\in\Gamma_R^{\mathrm{wr}}\).  A wrapped/wrapped
correspondence row is defined only for a retained target colour
\(\zeta\in\Gamma_R^{\mathrm{wr}}\) satisfying
\[
\zeta=\eta_1+\eta_2,
\]
with the positive wrapped condition read through \(b_R^{\mathrm{geom}}\).
The ordered wrapped/wrapped correspondence is the unreduced diagram
\[
\mathfrak E^{\mathrm{WW}}_{\eta_1,\eta_2;\zeta,R}
\xrightarrow{p^{\mathrm{WW}}}
\mathcal M_{\eta_1,R}^{\mathrm{wr,rig}}
\times
\mathcal M_{\eta_2,R}^{\mathrm{wr,rig}},
\qquad
\mathfrak E^{\mathrm{WW}}_{\eta_1,\eta_2;\zeta,R}
\xrightarrow{q^{\mathrm{WW}}}
\mathcal M_{\zeta,R}^{\mathrm{wr,rig}},
\]
classifying retained extensions
\[
0\longrightarrow W_{\eta_2}\longrightarrow B_\zeta
\longrightarrow W_{\eta_1}\longrightarrow0.
\]
This diagram lives over
\(\operatorname{Ran}^{\mathrm{wr}}_{R,\mathrm{pre}}(E)\times
\operatorname{Ran}^{\mathrm{wr}}_{R,\mathrm{pre}}(E)\) and targets
\(\operatorname{Ran}^{\mathrm{wr}}_{R,\mathrm{pre}}(E)\), before the
reduced \(E\)-quotient.  The source pair is ordered; symmetric descent
for swapping \(\eta_1\) and \(\eta_2\) is a separate row.  Source and
target anchor memory, admissibility, Thom--Sebastiani transport, and
transition compatibility are also separate rows.  The diagram is not a
mixed correspondence, not a local-local correspondence, not a
two-sided mixed correspondence, and not a scalar or quotient-first
replacement.
\end{definition}

The packet
\[
\texttt{certificates/hybrid/wrapped\_wrapped\_correspondence\_definition}
\]
verifies
\(\texttt{WRAPPED\_WRAPPED\_CORRESPONDENCE\_DEFINITION\_VERIFIED}\):
the ordered \(WW\) correspondence type is defined, while extension-stack
rows, source/target-map rows, populated anchor-memory rows, admissibility rows,
Thom--Sebastiani rows, symmetric-descent rows, and transition rows are
recorded as missing open obligations in that definition packet.

\begin{proposition}[Target admissibility for wrapped/wrapped extensions]
\label{prop:wrapped-wrapped-target-admissibility}
\textnormal{[proved for the Hall target map]}
Assume the wrapped prequotients
\(\mathcal M_{\eta_1,R}^{\mathrm{wr,rig}}\),
\(\mathcal M_{\eta_2,R}^{\mathrm{wr,rig}}\), and
\(\mathcal M_{\zeta,R}^{\mathrm{wr,rig}}\), the universal perfect
complexes, the retained closed-substack datum, and the retained
extension-closure datum are supplied at the fixed finite HN height
\(R\).  Then the target map in the ordered unreduced wrapped/wrapped
correspondence
\[
q^{\mathrm{WW}}:
\mathfrak E^{\mathrm{WW}}_{\eta_1,\eta_2;\zeta,R}
\longrightarrow
\mathcal M_{\zeta,R}^{\mathrm{wr,rig}}
\]
is proper.  Hence the target leg needed for the wrapped/wrapped Hall
pushforward is admissible as a proper geometric map before the reduced
\(E\)-quotient.  The source map \(p^{\mathrm{WW}}\) is not asserted to
be proper.
\end{proposition}

\begin{proof}
Let \(B_\zeta\) be the universal target object on
\(\mathcal M_{\zeta,R}^{\mathrm{wr,rig}}\).  The stack
\(\mathfrak E^{\mathrm{WW}}_{\eta_1,\eta_2;\zeta,R}\) is the closed
substack of the relative Quot scheme over
\(\mathcal M_{\zeta,R}^{\mathrm{wr,rig}}\) parametrising quotients
\[
B_\zeta\twoheadrightarrow W_{\eta_1}
\]
whose kernel is an object \(W_{\eta_2}\), with exact sequence
\[
0\longrightarrow W_{\eta_2}\longrightarrow B_\zeta
\longrightarrow W_{\eta_1}\longrightarrow0.
\]
The quotient and kernel are required to lie in the retained wrapped
prequotients of colours \(\eta_1\) and \(\eta_2\).  The retained
closed-substack datum makes these endpoint conditions closed on the
relative Quot scheme.  The retained colour condition
\(\zeta=\eta_1+\eta_2\) and the finite HN window condition are fixed
row conditions; the retained extension-closure datum records that the
allowed middle term remains in the finite window.  Grothendieck
projectivity of Quot \cite[Theorem~2.2.4]{HuybrechtsLehn} makes the
relative Quot scheme projective, hence proper, over
\(\mathcal M_{\zeta,R}^{\mathrm{wr,rig}}\).  A closed substack of this
relative Quot scheme is proper over the same target, so
\(q^{\mathrm{WW}}\) is proper.

The source map has the usual Hall extension fibre.  Over fixed
endpoints \((W_{\eta_1},W_{\eta_2})\), the extension parameters contain
\(\operatorname{Ext}^1(W_{\eta_1},W_{\eta_2})\) modulo the usual
\(\operatorname{Hom}\)-automorphisms.  When this
\(\operatorname{Ext}^1\)-group is nonzero, the fibre contains affine
directions and is not proper in general.  Thus the properness supplied
here is target properness of \(q^{\mathrm{WW}}\), not source
properness of \(p^{\mathrm{WW}}\).  The compact-support exceptional
pushforward model, symmetric descent, populated anchor-memory rows,
Thom--Sebastiani transport, quotient descent, transition compatibility,
and aggregate hybrid-carrier population remain separate rows.
\end{proof}

The packet
\[
\texttt{certificates/hybrid/wrapped\_wrapped\_extension\_admissibility}
\]
verifies
\(\texttt{WRAPPED\_WRAPPED\_ADMISSIBILITY\_VERIFIED}\): the ordered
wrapped/wrapped extension stack is represented as a closed relative
Quot locus over the wrapped target prequotient, the target map
\(q^{\mathrm{WW}}\) is proper, and the false source-properness
substitute is excluded.  The compact-support exceptional pushforward
model is supplied separately by
Definition~\ref{def:compact-support-exceptional-pushforward-model}.
Symmetric descent, populated anchor-memory rows, Thom--Sebastiani
transport, quotient descent, transition compatibility, and aggregate
hybrid-carrier population remain separate obligations.

\begin{definition}[Compact-support exceptional pushforward model]
\label{def:compact-support-exceptional-pushforward-model}
Let
\[
Y\xleftarrow{\ p\ }\mathfrak E\xrightarrow{\ f\ }Z
\]
be one of the retained finite correspondence maps in the local,
mixed, wrapped, or flag correspondence systems at HN height \(R\), and
let
\[
K_{\mathfrak E}\in D^b_c(\mathfrak E)
\]
be the constructible vanishing-cycle complex with orientation line on
the source of \(f\).  A compact-support exceptional pushforward model
for \((f,K_{\mathfrak E})\) is a tuple
\[
\mathfrak Q^{\mathrm{cs}}(f,K_{\mathfrak E})
=
(j_f,\overline{\mathfrak E}_f,\overline f,K_{\mathfrak E},
f^{\mathrm{cs}}_!K_{\mathfrak E})
\]
such that:
\begin{enumerate}
\item \(j_f:\mathfrak E\hookrightarrow\overline{\mathfrak E}_f\) is an
open immersion into a finite-type Artin stack with finite residual
inertia on the retained strata;
\item \(\overline f:\overline{\mathfrak E}_f\to Z\) is proper and
\[
f=\overline f\circ j_f;
\]
\item \(j_{f,!}K_{\mathfrak E}\) exists in the chosen constructible
six-functor theory on the retained d-critical stack, and
\[
f^{\mathrm{cs}}_!K_{\mathfrak E}
:=
\overline f_*\,j_{f,!}K_{\mathfrak E}
\in D^b_c(Z)
\]
is the compact-support exceptional pushforward along \(f\);
\item if \(f\) is proper, the identity compactification
\[
j_f=\mathrm{id}_{\mathfrak E},\qquad
\overline{\mathfrak E}_f=\mathfrak E,\qquad
\overline f=f
\]
is allowed, and then \(f^{\mathrm{cs}}_!=f_*\).
\end{enumerate}
The datum is a chosen compact-support model, not a theorem that two
compactifications give canonically identical functors.  It does not
prove base change, the projection formula, Thom--Sebastiani
compatibility, quotient descent, transition compatibility, symmetric
descent, or associativity on the two-step flag atlas.  Those are the
separate functorial rows required before
\[
q_!\operatorname{TS}^{\mathrm{red}}p^*
\]
is a coherent Hall multiplication on the finite hybrid carrier.
\end{definition}

The packet
\[
\texttt{certificates/hybrid/compact\_support\_exceptional\_pushforward\_model}
\]
verifies
\(\texttt{COMPACT\_SUPPORT\_EXCEPTIONAL\_PUSHFORWARD\_MODEL\_DEFINED}\):
Definition~\ref{def:compact-support-exceptional-pushforward-model}
defines the compactification, extension-by-zero, and proper
pushforward formula for a retained finite correspondence map, and it
records the proper-map specialization \(f^{\mathrm{cs}}_!=f_*\).  It
does not populate all correspondence maps, prove independence of the
compactification, prove base change, prove the projection formula,
transport Thom--Sebastiani, descend through the \(E\)-quotient, prove
transition compatibility, or populate the aggregate hybrid carrier.

\begin{proposition}[Base change for one-sided mixed correspondences]
\label{prop:mixed-correspondence-base-change}
\textnormal{[proved for Cartesian compact-support squares]}
Let \(o\in\{\mathrm{LW},\mathrm{WL}\}\), and let
\[
q^o:\mathfrak E^o_{\alpha,\eta;\zeta,R,I}
\longrightarrow Z^o_{\zeta,R}
:=\mathcal M_{\zeta,R}^{\mathrm{wr,rig}}
\]
be the target map of the corresponding one-sided mixed correspondence.
Let \(K^o\in D^b_c(\mathfrak E^o_{\alpha,\eta;\zeta,R,I})\) be the
constructible vanishing-cycle complex with orientation line.  Suppose
that a retained target morphism \(g:Z'\to Z^o_{\zeta,R}\) is supplied,
and that the mixed correspondence and the compact-support model are
pulled back by Cartesian squares
\[
\begin{tikzcd}
\mathfrak E^{o\prime}\ar[r,"\widetilde g"]\ar[d,"j'_o"']&
\mathfrak E^o_{\alpha,\eta;\zeta,R,I}\ar[d,"j_o"]\\
\overline{\mathfrak E}^{o\prime}\ar[r,"\overline g"]\ar[d,"\overline q^{o\prime}"']&
\overline{\mathfrak E}^{o}\ar[d,"\overline q^{o}"]\\
Z'\ar[r,"g"']&Z^o_{\zeta,R}
\end{tikzcd}
\]
with \(\overline q^o\) and \(\overline q^{o\prime}\) proper, and with
a coefficient identification
\[
K^{o\prime}\simeq \widetilde g^{\,*}K^o.
\]
Then the compact-support exceptional pushforward of
Definition~\ref{def:compact-support-exceptional-pushforward-model}
carries a canonical base-change isomorphism
\[
g^*(q^o)^{\mathrm{cs}}_!K^o
\;\xrightarrow{\;\sim\;}\;
(q^{o\prime})^{\mathrm{cs}}_!K^{o\prime}.
\]
For the proper target maps of
Proposition~\ref{prop:mixed-local-wrapped-target-admissibility}, the
identity compactification gives the ordinary proper-base-change
isomorphism for \(q^{\mathrm{LW}}\) and \(q^{\mathrm{WL}}\).
\end{proposition}

\begin{proof}
By definition,
\[
(q^o)^{\mathrm{cs}}_!K^o
=\overline q^o_*\,j_{o,!}K^o.
\]
The lower square is Cartesian and \(\overline q^o\) is proper, so the
proper base-change isomorphism in the chosen constructible
six-functor theory gives
\[
g^*\overline q^o_*\,j_{o,!}K^o
\simeq
\overline q^{o\prime}_*\,\overline g^{\,*}j_{o,!}K^o.
\]
The upper square is Cartesian and \(j_o\) is an open immersion; hence
extension by zero commutes with this base change:
\[
\overline g^{\,*}j_{o,!}K^o
\simeq
j'_{o,!}\,\widetilde g^{\,*}K^o.
\]
Using the supplied coefficient identification
\(K^{o\prime}\simeq\widetilde g^{\,*}K^o\), the last term becomes
\[
\overline q^{o\prime}_*\,j'_{o,!}K^{o\prime}
=(q^{o\prime})^{\mathrm{cs}}_!K^{o\prime}.
\]
This proves the displayed base-change isomorphism.  The argument
proves only the Cartesian compact-support base-change row for the
chosen model.  It does not prove the projection formula,
Thom--Sebastiani compatibility, quotient descent, transition
compatibility, choice independence for compactifications, or
associativity on the two-step flag atlas.
\end{proof}

The packet
\[
\texttt{certificates/hybrid/mixed\_correspondence\_base\_change}
\]
verifies
\(\texttt{MIXED\_CORRESPONDENCE\_BASE\_CHANGE\_VERIFIED}\):
Proposition~\ref{prop:mixed-correspondence-base-change} supplies the
Cartesian compact-support base-change isomorphisms for the \(LW\) and
\(WL\) target maps, including the proper-target specialization.  It
does not prove the projection formula, Thom--Sebastiani transport,
quotient descent, transition compatibility, compactification-choice
independence, or aggregate hybrid-carrier population.

\begin{proposition}[Projection formula for one-sided mixed correspondences]
\label{prop:mixed-correspondence-projection-formula}
\textnormal{[proved for retained finite-Tor target coefficients]}
Let \(o\in\{\mathrm{LW},\mathrm{WL}\}\), and let
\[
q^o:\mathfrak E^o_{\alpha,\eta;\zeta,R,I}
\longrightarrow Z^o_{\zeta,R}
:=\mathcal M_{\zeta,R}^{\mathrm{wr,rig}}
\]
be the target map of the corresponding one-sided mixed correspondence,
equipped with a chosen compact-support model
\[
\mathfrak E^o_{\alpha,\eta;\zeta,R,I}
\xrightarrow{\,j_o\,}
\overline{\mathfrak E}^{o}
\xrightarrow{\,\overline q^o\,}
Z^o_{\zeta,R},
\qquad q^o=\overline q^o\circ j_o,
\]
where \(j_o\) is an open immersion and \(\overline q^o\) is proper.
Let \(K^o\in D^b_c(\mathfrak E^o_{\alpha,\eta;\zeta,R,I})\) be the
constructible vanishing-cycle complex with orientation line.  Let
\(L\in D^b_c(Z^o_{\zeta,R})\) be a retained target coefficient of
finite Tor-amplitude, so that the open-immersion and proper projection
formulas hold in the chosen constructible six-functor theory.  Then
there is a canonical isomorphism
\[
(q^o)^{\mathrm{cs}}_!
\bigl(K^o\otimes^{\mathbf L}(q^o)^*L\bigr)
\;\xrightarrow{\;\sim\;}\;
(q^o)^{\mathrm{cs}}_!K^o\otimes^{\mathbf L}L .
\]
For the proper target maps of
Proposition~\ref{prop:mixed-local-wrapped-target-admissibility}, the
identity compactification specializes this to the ordinary proper
projection formula for \(q^{\mathrm{LW}}\) and \(q^{\mathrm{WL}}\).
\end{proposition}

\begin{proof}
Since \(q^o=\overline q^o\circ j_o\), one has
\[
(q^o)^*L=j_o^*(\overline q^o)^*L.
\]
Expanding the compact-support model gives
\[
(q^o)^{\mathrm{cs}}_!
\bigl(K^o\otimes^{\mathbf L}(q^o)^*L\bigr)
=
\overline q^o_*\,
j_{o,!}\bigl(K^o\otimes^{\mathbf L}j_o^*(\overline q^o)^*L\bigr).
\]
The projection formula for the open immersion \(j_o\) identifies the
right-hand side with
\[
\overline q^o_*
\bigl(j_{o,!}K^o\otimes^{\mathbf L}(\overline q^o)^*L\bigr).
\]
The map \(\overline q^o\) is proper; the proper projection formula
therefore gives
\[
\overline q^o_*
\bigl(j_{o,!}K^o\otimes^{\mathbf L}(\overline q^o)^*L\bigr)
\simeq
\overline q^o_*j_{o,!}K^o\otimes^{\mathbf L}L
=
(q^o)^{\mathrm{cs}}_!K^o\otimes^{\mathbf L}L.
\]
This proves the stated isomorphism.  The proof uses the chosen
compactification and the retained finite-Tor target coefficient.  It
does not prove compactification-choice independence,
Thom--Sebastiani compatibility, quotient descent, transition
compatibility, or associativity on the two-step flag atlas.
\end{proof}

The packet
\[
\texttt{certificates/hybrid/mixed\_correspondence\_projection\_formula}
\]
verifies
\(\texttt{MIXED\_CORRESPONDENCE\_PROJECTION\_FORMULA\_VERIFIED}\):
Proposition~\ref{prop:mixed-correspondence-projection-formula}
supplies the compact-support projection formula for the \(LW\) and
\(WL\) target maps with retained finite-Tor target coefficients.  It
does not prove Thom--Sebastiani transport, quotient descent, transition
compatibility, compactification-choice independence, or aggregate
hybrid-carrier population.

\begin{proposition}[Thom--Sebastiani transport for one-sided mixed correspondences]
\label{prop:mixed-correspondence-thom-sebastiani}
\textnormal{[proved for supplied reduced additive d-critical charts]}
Let \(o\in\{\mathrm{LW},\mathrm{WL}\}\), and let
\[
Y^o_{\mathrm{src}}
=
\begin{cases}
\mathfrak M^{\mathrm{loc,cl}}_{\alpha,R,I}
\times\mathcal M^{\mathrm{wr,rig}}_{\eta,R},&o=\mathrm{LW},\\[2mm]
\mathcal M^{\mathrm{wr,rig}}_{\eta,R}
\times\mathfrak M^{\mathrm{loc,cl}}_{\alpha,R,I},&o=\mathrm{WL}
\end{cases}
\]
be the ordered source product of the one-sided mixed correspondence
with source map \(p^o:\mathfrak E^o_{\alpha,\eta;\zeta,R,I}\to
Y^o_{\mathrm{src}}\).  Write
\[
K^o_{\mathrm{src}}
=
\Phi^{\mathrm{red}}_{Y^o_{\mathrm{src}}}\otimes
o^{\mathrm{red}}_{Y^o_{\mathrm{src}}},
\qquad
K^o_{\mathfrak E}
=
\Phi^{\mathrm{red}}_{\mathfrak E^o_{\alpha,\eta;\zeta,R,I}}
\otimes
o^{\mathrm{red}}_{\mathfrak E^o_{\alpha,\eta;\zeta,R,I}}.
\]
Assume that the retained mixed extension stack is covered by
Joyce--BBDJS d-critical charts for which:
\begin{enumerate}
\item the extension potential is the pullback of the ordered sum of
the two source potentials, up to the standard nondegenerate quadratic
stabilization;
\item the K3 semiregularity cosection is additive along the mixed
extension correspondence, so the reduced obstruction theories and
reduced vanishing-cycle complexes are pulled from the same quotient
of the ordinary obstruction theory;
\item the Joyce--Upmeier orientation line on
\(\mathfrak E^o_{\alpha,\eta;\zeta,R,I}\) is identified with
\((p^o)^*o^{\mathrm{red}}_{Y^o_{\mathrm{src}}}\), and the
orientation Thom--Sebastiani defect is zero on chart overlaps.
\end{enumerate}
Then the mixed correspondence carries a canonical reduced
Thom--Sebastiani transport isomorphism
\[
\operatorname{TS}^{\mathrm{red},o}_{\alpha,\eta;R,I}:
(p^o)^*K^o_{\mathrm{src}}
\;\xrightarrow{\;\sim\;}\;
K^o_{\mathfrak E}
\]
in \(D^b_c(\mathfrak E^o_{\alpha,\eta;\zeta,R,I})\).  The construction
is ordered: the \(LW\) and \(WL\) isomorphisms differ by the Koszul
symmetry of the source product and are recorded as separate rows.
\end{proposition}

\begin{proof}
On a retained d-critical chart, condition (1) puts the mixed extension
potential in the BBDJS Thom--Sebastiani normal form for the ordered
sum of the two source potentials, after harmless quadratic
stabilization \cite{BBDJS2015}.  Hence the BBDJS
Thom--Sebastiani theorem gives an isomorphism between the pullback of
the exterior product of the two source vanishing-cycle complexes and
the vanishing-cycle complex of the extension chart.  Condition (2)
identifies the ordinary and cosection-reduced sides of this isomorphism
after quotienting by the common additive semiregularity cosection.
Condition (3) tensors the vanishing-cycle isomorphism with the
Joyce--Upmeier orientation transport and makes it independent of the
chosen chart overlap.  The descent data therefore glue the chartwise
isomorphisms to the displayed object of
\(D^b_c(\mathfrak E^o_{\alpha,\eta;\zeta,R,I})\).

The proof supplies only the binary mixed Thom--Sebastiani transport
for the chosen reduced d-critical and orientation atlas.  It does not
construct the orientation line or discharge \((O1)\).  The
quotient-after-correspondence descent of this transport is
Proposition~\ref{prop:quotient-after-correspondence-preserves-thom-sebastiani}.
Transition compatibility and the eight-word two-step flag pentagon
remain separate obligations.
\end{proof}

The packet
\[
\texttt{certificates/hybrid/mixed\_correspondence\_thom\_sebastiani}
\]
verifies
\(\texttt{MIXED\_THOM\_SEBASTIANI\_TRANSPORT\_CONDITIONAL\_VERIFIED}\):
Proposition~\ref{prop:mixed-correspondence-thom-sebastiani} supplies
the binary mixed Thom--Sebastiani transport for \(LW\) and \(WL\) from
supplied additive reduced d-critical charts and zero-defect orientation
transport.  It does not prove \((O1)\), quotient descent, transition
compatibility, two-step flag associativity, or aggregate
hybrid-carrier population.

\begin{definition}[Eight local/wrapped binary words]
\label{def:eight-local-wrapped-binary-words}
At a fixed finite HN height \(R\), let
\[
\mathbb A_{\mathrm{hyb}}=\{\mathrm L,\mathrm W\}
\]
be the ordered hybrid alphabet, where \(\mathrm L\) denotes a
projection-finite local input of elliptic degree \(0\), and
\(\mathrm W\) denotes a rigidified wrapped input of positive elliptic
degree.  The binary type product is
\[
\mathrm L\star\mathrm L=\mathrm L,\qquad
\mathrm L\star\mathrm W=\mathrm W,\qquad
\mathrm W\star\mathrm L=\mathrm W,\qquad
\mathrm W\star\mathrm W=\mathrm W.
\]
Thus a length-three word
\[
w=\epsilon_1\epsilon_2\epsilon_3\in\mathbb A_{\mathrm{hyb}}^3
\]
has left and right intermediate types
\[
\epsilon_{12}=\epsilon_1\star\epsilon_2,\qquad
\epsilon_{23}=\epsilon_2\star\epsilon_3,
\]
and final type
\[
\epsilon_{123}=(\epsilon_1\star\epsilon_2)\star\epsilon_3
=\epsilon_1\star(\epsilon_2\star\epsilon_3).
\]
The eight ordered words are
\[
\mathcal W^{(3)}_{\mathrm{hyb}}
=
\{\mathrm{LLL},\mathrm{LLW},\mathrm{LWL},\mathrm{WLL},
\mathrm{LWW},\mathrm{WLW},\mathrm{WWL},\mathrm{WWW}\}.
\]
For compatible ordered charges \((\xi_1,\xi_2,\xi_3)\), the two
parenthesization symbols attached to \(w\) are
\[
((\epsilon_1\epsilon_2)\epsilon_3):
(\xi_1,\xi_2,\xi_3)\mapsto((\xi_1+\xi_2),\xi_3),
\]
\[
(\epsilon_1(\epsilon_2\epsilon_3)):
(\xi_1,\xi_2,\xi_3)\mapsto(\xi_1,(\xi_2+\xi_3)).
\]
This is only the word and type convention.  It does not construct the
two-step flag stacks, prove that the stacks are retained, or prove
associativity, pentagon coherence, quotient descent, transition
compatibility, or aggregate hybrid-carrier population.
\end{definition}

The packet
\[
\texttt{certificates/hybrid/eight\_word\_binary\_vocabulary}
\]
verifies
\(\texttt{EIGHT\_WORD\_BINARY\_VOCABULARY\_DEFINED}\):
Definition~\ref{def:eight-local-wrapped-binary-words} fixes the
ordered alphabet, the type product, the eight length-three words, and
the two parenthesization symbols.  It does not construct the two-step
flag stacks or prove associativity.

\begin{proposition}[Eight two-step local/wrapped flag stacks]
\label{prop:eight-two-step-flag-stacks}
\textnormal{[constructed as retained flag-Quot substacks]}
Assume the retained local and wrapped object stacks, the four ordered
binary correspondence stacks \(LL,LW,WL,WW\), the retained universal
complexes, closed endpoint conditions, and retained extension closure
at height \(R\).  For every word
\[
w=\epsilon_1\epsilon_2\epsilon_3\in
\mathcal W^{(3)}_{\mathrm{hyb}}
\]
and every compatible ordered colour triple
\((\xi_1,\xi_2,\xi_3)\) of those types, there is a retained
finite-type Artin stack
\[
\mathfrak F^{(2),w}_{\xi_1,\xi_2,\xi_3,R}
\]
classifying two-step filtrations
\[
0=C_0\subset C_1\subset C_2\subset C_3
\]
in the finite HN heart such that
\[
C_i/C_{i-1}\in\mathfrak M^{\epsilon_i}_{\xi_i,R}
\qquad (i=1,2,3),
\]
the intermediate objects \(C_2/C_0\), \(C_3/C_1\), and \(C_3/C_0\)
have the types
\[
\epsilon_{12},\qquad \epsilon_{23},\qquad \epsilon_{123}
\]
of Definition~\ref{def:eight-local-wrapped-binary-words}, and every
intermediate colour is retained.  If one of these types is
\(\mathrm W\), the corresponding object carries the fixed wrapped
rigidification and the before-quotient anchor memory.

Let
\[
\mathfrak E^{\epsilon\delta}_{\xi,\upsilon;\xi+\upsilon,R}
\]
denote the ordered binary correspondence stack of type
\(\epsilon\delta\), interpreted as \(LL,LW,WL\), or \(WW\) by
Definitions~\ref{def:local-local-correspondence-datum},
\ref{def:mixed-local-wrapped-correspondence-datum}, and
\ref{def:wrapped-wrapped-correspondence-datum}.  The flag stack carries
canonical comparison morphisms
\[
\lambda^w:
\mathfrak F^{(2),w}_{\xi_1,\xi_2,\xi_3,R}
\longrightarrow
\mathfrak E^{\epsilon_{12}\epsilon_3}_{\xi_1+\xi_2,\xi_3;\xi_1+\xi_2+\xi_3,R}
\times_{\mathfrak M^{\epsilon_{12}}_{\xi_1+\xi_2,R}}
\mathfrak E^{\epsilon_1\epsilon_2}_{\xi_1,\xi_2;\xi_1+\xi_2,R}
\]
and
\[
\rho^w:
\mathfrak F^{(2),w}_{\xi_1,\xi_2,\xi_3,R}
\longrightarrow
\mathfrak E^{\epsilon_1\epsilon_{23}}_{\xi_1,\xi_2+\xi_3;\xi_1+\xi_2+\xi_3,R}
\times_{\mathfrak M^{\epsilon_{23}}_{\xi_2+\xi_3,R}}
\mathfrak E^{\epsilon_2\epsilon_3}_{\xi_2,\xi_3;\xi_2+\xi_3,R}.
\]
The maps are defined by the two forgetful operations on the filtration:
\(\lambda^w\) remembers \(C_1\subset C_2\subset C_3\) as
\(((\xi_1,\xi_2),\xi_3)\), while \(\rho^w\) remembers it as
\((\xi_1,(\xi_2,\xi_3))\).
\end{proposition}

\begin{proof}
Fix \(w\) and \((\xi_1,\xi_2,\xi_3)\).  Over the retained terminal
object stack of colour \(\xi_1+\xi_2+\xi_3\), take the relative
two-step flag-Quot stack parametrising quotients
\[
C_3\twoheadrightarrow C_3/C_2,\qquad
C_2\twoheadrightarrow C_2/C_1
\]
or equivalently subobjects
\(C_1\subset C_2\subset C_3\).  Grothendieck projectivity for
relative Quot schemes and its flag version give finite type over the
retained terminal stack \cite[Theorem~2.2.4]{HuybrechtsLehn}.  The
conditions that the three graded pieces lie in the prescribed local or
wrapped retained substacks, that the intermediate colours equal the
displayed sums, and that the wrapped pieces carry the fixed
rigidifications and anchor memory are closed retained row conditions.
The retained extension-closure packet ensures that \(C_2/C_0\),
\(C_3/C_1\), and \(C_3/C_0\) remain in the finite HN window.  Hence
the required \(\mathfrak F^{(2),w}_{\xi_1,\xi_2,\xi_3,R}\) is a
retained finite-type closed substack of the flag-Quot stack.

Forgetting the first or the second extension step gives the two
displayed morphisms to the iterated correspondence fibre products.
The construction is carried out before the reduced \(E\)-quotient and
only constructs the eight flag stacks and their comparison maps.  It
does not prove that \(\lambda^w\) and \(\rho^w\) induce the same Hall
operation, does not prove associativity defects vanish, and does not
prove the four-input pentagon, quotient descent, transition
compatibility, or aggregate hybrid-carrier population.
\end{proof}

The packet
\[
\texttt{certificates/hybrid/eight\_word\_two\_step\_flag\_stacks}
\]
verifies
\(\texttt{EIGHT\_WORD\_TWO\_STEP\_FLAG\_STACKS\_CONSTRUCTED}\):
Proposition~\ref{prop:eight-two-step-flag-stacks} constructs the
retained two-step flag stack and the two comparison maps for each of
the eight words.  It does not prove associativity, pentagon coherence,
quotient descent, transition compatibility, or aggregate population.

\begin{proposition}[Associativity for the \(LLL\) word]
\label{prop:lll-associativity}
\textnormal{[proved for the pure local two-step flag stack]}
Let \(\alpha,\beta,\delta\) be retained local colours at height \(R\),
and write
\[
\gamma_{12}=\alpha+\beta,\qquad
\gamma_{23}=\beta+\delta,\qquad
\gamma_{123}=\alpha+\beta+\delta .
\]
Assume the \(LLL\) flag stack
\(\mathfrak F^{(2),\mathrm{LLL}}_{\alpha,\beta,\delta,R}\) of
Proposition~\ref{prop:eight-two-step-flag-stacks}, the local/local
target-proper correspondence rows of
Proposition~\ref{prop:local-local-target-properness}, the chosen
compact-support pushforward model, and the local base-change,
projection-formula, and reduced Thom--Sebastiani coefficient
identifications on the \(LLL\) flag square.  Then the pure local Hall
product satisfies
\[
m^{\mathrm{LL}}_{\gamma_{12},\delta,R}
\circ
\bigl(m^{\mathrm{LL}}_{\alpha,\beta,R}\otimes 1\bigr)
=
m^{\mathrm{LL}}_{\alpha,\gamma_{23},R}
\circ
\bigl(1\otimes m^{\mathrm{LL}}_{\beta,\delta,R}\bigr)
\]
as morphisms on the retained compactly supported reduced
vanishing-cycle complexes of the three local inputs.  Equivalently,
the \(LLL\) associativity defect class
\[
\mathfrak o^{\mathrm{assoc}}_{\mathrm{LLL},R}
\]
vanishes.
\end{proposition}

\begin{proof}
Both composites are pull-push operations attached to the same
two-step local filtration stack
\[
0\subset C_1\subset C_2\subset C_3
\]
with graded pieces of colours \(\alpha,\beta,\delta\).  The left
parenthesization uses the comparison map \(\lambda^{\mathrm{LLL}}\) to
first form the \(LL\) extension of \(\alpha\) by \(\beta\), then the
extension of \(\gamma_{12}\) by \(\delta\).  The right
parenthesization uses \(\rho^{\mathrm{LLL}}\) to first form the
extension of \(\beta\) by \(\delta\), then the extension of \(\alpha\)
by \(\gamma_{23}\).  These are two descriptions of the same flag
object, so their source and target maps factor through the same
retained flag stack.

Proper base change identifies the two iterated target pushforwards
after pullback to the flag stack; the projection formula moves the
source coefficients through the same compact-support pushforward; and
the reduced Thom--Sebastiani identification on the \(LLL\) flag square
identifies the two tensor products of vanishing-cycle and orientation
lines.  Hence the two pull-push composites agree.  This proves
\(\mathfrak o^{\mathrm{assoc}}_{\mathrm{LLL},R}=0\).

The argument is purely local.  It proves no associativity statement
for \(LLW,LWL,WLL,LWW,WLW,WWL\), or \(WWW\), and it does not prove the
four-input pentagon, quotient descent, transition compatibility, or
aggregate hybrid-carrier population.
\end{proof}

The packet
\[
\texttt{certificates/hybrid/lll\_associativity}
\]
verifies
\(\texttt{LLL\_ASSOCIATIVITY\_CONDITIONAL\_VERIFIED}\):
Proposition~\ref{prop:lll-associativity} proves the \(LLL\) word
associativity defect vanishes under the stated local flag-stack,
base-change, projection-formula, and reduced Thom--Sebastiani
coefficient rows.  It does not prove the seven remaining word
associativity rows or the four-input pentagon.

\begin{proposition}[Associativity for the \(LLW\) word]
\label{prop:llw-associativity}
\textnormal{[proved for the local/local--mixed flag stack]}
Let \(\alpha,\beta\) be retained local colours and let \(\eta\) be a
retained wrapped colour at height \(R\).  Write
\[
\gamma_{12}=\alpha+\beta,\qquad
\zeta_{23}=\beta+\eta,\qquad
\zeta_{123}=\alpha+\beta+\eta .
\]
Assume the \(LLW\) flag stack
\(\mathfrak F^{(2),\mathrm{LLW}}_{\alpha,\beta,\eta,R}\) of
Proposition~\ref{prop:eight-two-step-flag-stacks}, the \(LL\) and
\(LW\) target-admissible correspondence rows, the chosen
compact-support pushforward model, and the required local and mixed
base-change, projection-formula, and reduced Thom--Sebastiani
coefficient identifications on the \(LLW\) flag square.  Then
\[
m^{\mathrm{LW}}_{\gamma_{12},\eta,R}
\circ
\bigl(m^{\mathrm{LL}}_{\alpha,\beta,R}\otimes 1\bigr)
=
m^{\mathrm{LW}}_{\alpha,\zeta_{23},R}
\circ
\bigl(1\otimes m^{\mathrm{LW}}_{\beta,\eta,R}\bigr)
\]
as morphisms on the retained compactly supported reduced
vanishing-cycle complexes of the ordered inputs.  Equivalently,
\[
\mathfrak o^{\mathrm{assoc}}_{\mathrm{LLW},R}=0 .
\]
\end{proposition}

\begin{proof}
Both composites are represented by the same retained flag stack
\[
0\subset C_1\subset C_2\subset C_3,
\qquad
C_1/C_0\in L_\alpha,\quad
C_2/C_1\in L_\beta,\quad
C_3/C_2\in W_\eta .
\]
The left comparison map \(\lambda^{\mathrm{LLW}}\) records first the
local/local extension of \(\alpha\) by \(\beta\), and then the
local/wrapped extension of \(\gamma_{12}\) by \(\eta\).  The right
comparison map \(\rho^{\mathrm{LLW}}\) records first the
local/wrapped extension of \(\beta\) by \(\eta\), and then the
local/wrapped extension of \(\alpha\) by \(\zeta_{23}\).  These two
iterated descriptions have the same underlying filtered object.

The local and mixed base-change isomorphisms identify the two iterated
compact-support pushforwards after pullback to
\(\mathfrak F^{(2),\mathrm{LLW}}_{\alpha,\beta,\eta,R}\).  The
projection formulas move the external source coefficients through the
same target pushforward, and the reduced Thom--Sebastiani transports
identify the two coefficient systems on the common flag stack.  Hence
the two composites agree, so
\(\mathfrak o^{\mathrm{assoc}}_{\mathrm{LLW},R}=0\).

The proof uses no wrapped/wrapped binary product.  It proves no
associativity statement for \(LWL,WLL,LWW,WLW,WWL\), or \(WWW\), and it
does not prove the four-input pentagon, quotient descent, transition
compatibility, or aggregate hybrid-carrier population.
\end{proof}

The packet
\[
\texttt{certificates/hybrid/llw\_associativity}
\]
verifies
\(\texttt{LLW\_ASSOCIATIVITY\_CONDITIONAL\_VERIFIED}\):
Proposition~\ref{prop:llw-associativity} proves the \(LLW\) word
associativity defect vanishes under the stated \(LL/LW\) flag-stack,
base-change, projection-formula, and reduced Thom--Sebastiani
coefficient rows.  It does not prove the six remaining word
associativity rows or the four-input pentagon.

\begin{proposition}[Associativity for the \(LWL\) word]
\label{prop:lwl-associativity}
\textnormal{[proved for the mixed--mixed flag stack]}
Let \(\alpha,\beta\) be retained local colours and let \(\eta\) be a
retained wrapped colour at height \(R\), ordered as
\((\alpha,\eta,\beta)\).  Write
\[
\zeta_{12}=\alpha+\eta,\qquad
\zeta_{23}=\eta+\beta,\qquad
\zeta_{123}=\alpha+\eta+\beta .
\]
Assume the \(LWL\) flag stack
\(\mathfrak F^{(2),\mathrm{LWL}}_{\alpha,\eta,\beta,R}\) of
Proposition~\ref{prop:eight-two-step-flag-stacks}, the \(LW\) and
\(WL\) target-admissible correspondence rows, the chosen
compact-support pushforward model, and the required mixed base-change,
projection-formula, and reduced Thom--Sebastiani coefficient
identifications on the \(LWL\) flag square.  Then
\[
m^{\mathrm{WL}}_{\zeta_{12},\beta,R}
\circ
\bigl(m^{\mathrm{LW}}_{\alpha,\eta,R}\otimes 1\bigr)
=
m^{\mathrm{LW}}_{\alpha,\zeta_{23},R}
\circ
\bigl(1\otimes m^{\mathrm{WL}}_{\eta,\beta,R}\bigr)
\]
as morphisms on the retained compactly supported reduced
vanishing-cycle complexes of the ordered inputs.  Equivalently,
\[
\mathfrak o^{\mathrm{assoc}}_{\mathrm{LWL},R}=0 .
\]
\end{proposition}

\begin{proof}
The common refinement is the retained \(LWL\) flag stack
\[
0\subset C_1\subset C_2\subset C_3,
\qquad
C_1/C_0\in L_\alpha,\quad
C_2/C_1\in W_\eta,\quad
C_3/C_2\in L_\beta .
\]
The left comparison \(\lambda^{\mathrm{LWL}}\) reads this filtration as
an \(LW\) extension followed by a \(WL\) extension.  The right
comparison \(\rho^{\mathrm{LWL}}\) reads it as a \(WL\) extension
followed by an \(LW\) extension.  Both descriptions have the same
terminal filtered object of wrapped colour \(\zeta_{123}\).

The mixed base-change isomorphisms identify the two iterated
compact-support pushforwards on the common flag stack.  The mixed
projection formula moves the ordered source coefficients through the
same pushforward, and the mixed reduced Thom--Sebastiani transports
identify the two coefficient systems.  Therefore the two composites
agree, and
\(\mathfrak o^{\mathrm{assoc}}_{\mathrm{LWL},R}=0\).

The proof is the ordered \(LWL\) case only.  It proves no
associativity statement for \(WLL,LWW,WLW,WWL\), or \(WWW\), and it
does not prove the four-input pentagon, quotient descent, transition
compatibility, or aggregate hybrid-carrier population.
\end{proof}

The packet
\[
\texttt{certificates/hybrid/lwl\_associativity}
\]
verifies
\(\texttt{LWL\_ASSOCIATIVITY\_CONDITIONAL\_VERIFIED}\):
Proposition~\ref{prop:lwl-associativity} proves the \(LWL\) word
associativity defect vanishes under the stated \(LW/WL\) flag-stack,
base-change, projection-formula, and reduced Thom--Sebastiani
coefficient rows.  It does not prove the five remaining word
associativity rows or the four-input pentagon.

\begin{proposition}[Associativity for the \(WLL\) word]
\label{prop:wll-associativity}
\textnormal{[proved for the mixed--local/local flag stack]}
Let \(\alpha,\beta\) be retained local colours and let \(\eta\) be a
retained wrapped colour at height \(R\), ordered as
\((\eta,\alpha,\beta)\).  Write
\[
\zeta_{12}=\eta+\alpha,\qquad
\gamma_{23}=\alpha+\beta,\qquad
\zeta_{123}=\eta+\alpha+\beta .
\]
Assume the \(WLL\) flag stack
\(\mathfrak F^{(2),\mathrm{WLL}}_{\eta,\alpha,\beta,R}\) of
Proposition~\ref{prop:eight-two-step-flag-stacks}, the \(WL\) and
\(LL\) target-admissible correspondence rows, the chosen
compact-support pushforward model, and the required local and mixed
base-change, projection-formula, and reduced Thom--Sebastiani
coefficient identifications on the \(WLL\) flag square.  Then
\[
m^{\mathrm{WL}}_{\zeta_{12},\beta,R}
\circ
\bigl(m^{\mathrm{WL}}_{\eta,\alpha,R}\otimes 1\bigr)
=
m^{\mathrm{WL}}_{\eta,\gamma_{23},R}
\circ
\bigl(1\otimes m^{\mathrm{LL}}_{\alpha,\beta,R}\bigr)
\]
as morphisms on the retained compactly supported reduced
vanishing-cycle complexes of the ordered inputs.  Equivalently,
\[
\mathfrak o^{\mathrm{assoc}}_{\mathrm{WLL},R}=0 .
\]
\end{proposition}

\begin{proof}
The common refinement is the retained \(WLL\) flag stack
\[
0\subset C_1\subset C_2\subset C_3,
\qquad
C_1/C_0\in W_\eta,\quad
C_2/C_1\in L_\alpha,\quad
C_3/C_2\in L_\beta .
\]
The left comparison \(\lambda^{\mathrm{WLL}}\) records first the
wrapped/local extension of \(\eta\) by \(\alpha\), then the
wrapped/local extension of \(\zeta_{12}\) by \(\beta\).  The right
comparison \(\rho^{\mathrm{WLL}}\) records first the local/local
extension of \(\alpha\) by \(\beta\), then the wrapped/local extension
of \(\eta\) by \(\gamma_{23}\).  Both descriptions have terminal
wrapped colour \(\zeta_{123}\).

The mixed and local base-change isomorphisms identify the two iterated
compact-support pushforwards on the common flag stack.  The projection
formulas move the ordered source coefficients through the same
pushforward, and the reduced Thom--Sebastiani transports identify the
two coefficient systems.  Therefore the two composites agree, and
\(\mathfrak o^{\mathrm{assoc}}_{\mathrm{WLL},R}=0\).

The proof uses no wrapped/wrapped binary product.  It proves no
associativity statement for \(LWW,WLW,WWL\), or \(WWW\), and it does
not prove the four-input pentagon, quotient descent, transition
compatibility, or aggregate hybrid-carrier population.
\end{proof}

The packet
\[
\texttt{certificates/hybrid/wll\_associativity}
\]
verifies
\(\texttt{WLL\_ASSOCIATIVITY\_CONDITIONAL\_VERIFIED}\):
Proposition~\ref{prop:wll-associativity} proves the \(WLL\) word
associativity defect vanishes under the stated \(WL/LL\) flag-stack,
base-change, projection-formula, and reduced Thom--Sebastiani
coefficient rows.  It does not prove the four remaining word
associativity rows or the four-input pentagon.

\begin{proposition}[Associativity for the \(LWW\) word]
\label{prop:lww-associativity}
\textnormal{[proved conditionally for the mixed--wrapped/wrapped flag stack]}
Let \(\alpha\) be a retained local colour and let \(\eta_1,\eta_2\)
be retained wrapped colours at height \(R\), ordered as
\((\alpha,\eta_1,\eta_2)\).  Write
\[
\zeta_{12}=\alpha+\eta_1,\qquad
\zeta_{23}=\eta_1+\eta_2,\qquad
\zeta_{123}=\alpha+\eta_1+\eta_2 .
\]
Assume the \(LWW\) flag stack
\(\mathfrak F^{(2),\mathrm{LWW}}_{\alpha,\eta_1,\eta_2,R}\) of
Proposition~\ref{prop:eight-two-step-flag-stacks}, the \(LW\) and
\(WW\) target-admissible correspondence rows, the chosen
compact-support pushforward model, and the required mixed and
wrapped/wrapped base-change, projection-formula, and reduced
Thom--Sebastiani coefficient identifications on the \(LWW\) flag
square.  Then
\[
m^{\mathrm{WW}}_{\zeta_{12},\eta_2,R}
\circ
\bigl(m^{\mathrm{LW}}_{\alpha,\eta_1,R}\otimes 1\bigr)
=
m^{\mathrm{LW}}_{\alpha,\zeta_{23},R}
\circ
\bigl(1\otimes m^{\mathrm{WW}}_{\eta_1,\eta_2,R}\bigr)
\]
as morphisms on the retained compactly supported reduced
vanishing-cycle complexes of the ordered inputs.  Equivalently,
\[
\mathfrak o^{\mathrm{assoc}}_{\mathrm{LWW},R}=0 .
\]
The wrapped/wrapped base-change, projection-formula, and reduced
Thom--Sebastiani rows are hypotheses of this proposition; they are not
supplied by the wrapped/wrapped target-admissibility theorem.
\end{proposition}

\begin{proof}
The common refinement is the retained \(LWW\) flag stack
\[
0\subset C_1\subset C_2\subset C_3,
\qquad
C_1/C_0\in L_\alpha,\quad
C_2/C_1\in W_{\eta_1},\quad
C_3/C_2\in W_{\eta_2}.
\]
The left comparison \(\lambda^{\mathrm{LWW}}\) records first the
local/wrapped extension of \(\alpha\) by \(\eta_1\), then the
wrapped/wrapped extension of \(\zeta_{12}\) by \(\eta_2\).  The right
comparison \(\rho^{\mathrm{LWW}}\) records first the wrapped/wrapped
extension of \(\eta_1\) by \(\eta_2\), then the local/wrapped
extension of \(\alpha\) by \(\zeta_{23}\).  Both descriptions have
terminal wrapped colour \(\zeta_{123}\).

The mixed and wrapped/wrapped base-change isomorphisms identify the
two iterated compact-support pushforwards on the common flag stack.
The corresponding projection formulas move the ordered source
coefficients through the same pushforward, and the reduced
Thom--Sebastiani transports identify the two coefficient systems.
Therefore the two composites agree, and
\(\mathfrak o^{\mathrm{assoc}}_{\mathrm{LWW},R}=0\).

The proof is wordwise and conditional on the wrapped/wrapped
functorial coefficient rows.  It proves no associativity statement for
\(WLW,WWL\), or \(WWW\), and it does not prove the four-input
pentagon, quotient descent, transition compatibility, or aggregate
hybrid-carrier population.
\end{proof}

The packet
\[
\texttt{certificates/hybrid/lww\_associativity}
\]
verifies
\(\texttt{LWW\_ASSOCIATIVITY\_CONDITIONAL\_VERIFIED}\):
Proposition~\ref{prop:lww-associativity} proves the \(LWW\) word
associativity defect vanishes under the stated \(LW/WW\) flag-stack,
mixed functorial rows, and wrapped/wrapped functorial hypotheses.  It
does not prove the three remaining word associativity rows or the
four-input pentagon.

\begin{proposition}[Associativity for the \(WLW\) word]
\label{prop:wlw-associativity}
\textnormal{[proved conditionally for the mixed--wrapped/wrapped flag stack]}
Let \(\alpha\) be a retained local colour and let \(\eta_1,\eta_2\)
be retained wrapped colours at height \(R\), ordered as
\((\eta_1,\alpha,\eta_2)\).  Write
\[
\zeta_{12}=\eta_1+\alpha,\qquad
\zeta_{23}=\alpha+\eta_2,\qquad
\zeta_{123}=\eta_1+\alpha+\eta_2 .
\]
Assume the \(WLW\) flag stack
\(\mathfrak F^{(2),\mathrm{WLW}}_{\eta_1,\alpha,\eta_2,R}\) of
Proposition~\ref{prop:eight-two-step-flag-stacks}, the \(WL\), \(LW\),
and \(WW\) target-admissible correspondence rows, the chosen
compact-support pushforward model, and the required mixed and
wrapped/wrapped base-change, projection-formula, and reduced
Thom--Sebastiani coefficient identifications on the \(WLW\) flag
square.  Then
\[
m^{\mathrm{WW}}_{\zeta_{12},\eta_2,R}
\circ
\bigl(m^{\mathrm{WL}}_{\eta_1,\alpha,R}\otimes 1\bigr)
=
m^{\mathrm{WW}}_{\eta_1,\zeta_{23},R}
\circ
\bigl(1\otimes m^{\mathrm{LW}}_{\alpha,\eta_2,R}\bigr)
\]
as morphisms on the retained compactly supported reduced
vanishing-cycle complexes of the ordered inputs.  Equivalently,
\[
\mathfrak o^{\mathrm{assoc}}_{\mathrm{WLW},R}=0 .
\]
The wrapped/wrapped base-change, projection-formula, and reduced
Thom--Sebastiani rows are hypotheses of this proposition; they are not
supplied by the wrapped/wrapped target-admissibility theorem.
\end{proposition}

\begin{proof}
The common refinement is the retained \(WLW\) flag stack
\[
0\subset C_1\subset C_2\subset C_3,
\qquad
C_1/C_0\in W_{\eta_1},\quad
C_2/C_1\in L_\alpha,\quad
C_3/C_2\in W_{\eta_2}.
\]
The left comparison \(\lambda^{\mathrm{WLW}}\) records first the
wrapped/local extension of \(\eta_1\) by \(\alpha\), then the
wrapped/wrapped extension of \(\zeta_{12}\) by \(\eta_2\).  The right
comparison \(\rho^{\mathrm{WLW}}\) records first the local/wrapped
extension of \(\alpha\) by \(\eta_2\), then the wrapped/wrapped
extension of \(\eta_1\) by \(\zeta_{23}\).  Both descriptions have
terminal wrapped colour \(\zeta_{123}\).

The mixed and wrapped/wrapped base-change isomorphisms identify the
two iterated compact-support pushforwards on the common flag stack.
The corresponding projection formulas move the ordered source
coefficients through the same pushforward, and the reduced
Thom--Sebastiani transports identify the two coefficient systems.
Therefore the two composites agree, and
\(\mathfrak o^{\mathrm{assoc}}_{\mathrm{WLW},R}=0\).

The proof is wordwise and conditional on the wrapped/wrapped
functorial coefficient rows.  It proves no associativity statement for
\(WWL\) or \(WWW\), and it does not prove the four-input pentagon,
quotient descent, transition compatibility, or aggregate
hybrid-carrier population.
\end{proof}

The packet
\[
\texttt{certificates/hybrid/wlw\_associativity}
\]
verifies
\(\texttt{WLW\_ASSOCIATIVITY\_CONDITIONAL\_VERIFIED}\):
Proposition~\ref{prop:wlw-associativity} proves the \(WLW\) word
associativity defect vanishes under the stated \(WL/LW/WW\)
flag-stack, mixed functorial rows, and wrapped/wrapped functorial
hypotheses.  It does not prove the two remaining word associativity
rows or the four-input pentagon.

\begin{proposition}[Associativity for the \(WWL\) word]
\label{prop:wwl-associativity}
\textnormal{[proved conditionally for the wrapped/wrapped--mixed flag stack]}
Let \(\alpha\) be a retained local colour and let \(\eta_1,\eta_2\)
be retained wrapped colours at height \(R\), ordered as
\((\eta_1,\eta_2,\alpha)\).  Write
\[
\zeta_{12}=\eta_1+\eta_2,\qquad
\zeta_{23}=\eta_2+\alpha,\qquad
\zeta_{123}=\eta_1+\eta_2+\alpha .
\]
Assume the \(WWL\) flag stack
\(\mathfrak F^{(2),\mathrm{WWL}}_{\eta_1,\eta_2,\alpha,R}\) of
Proposition~\ref{prop:eight-two-step-flag-stacks}, the \(WW\) and
\(WL\) target-admissible correspondence rows, the chosen
compact-support pushforward model, and the required mixed and
wrapped/wrapped base-change, projection-formula, and reduced
Thom--Sebastiani coefficient identifications on the \(WWL\) flag
square.  Then
\[
m^{\mathrm{WL}}_{\zeta_{12},\alpha,R}
\circ
\bigl(m^{\mathrm{WW}}_{\eta_1,\eta_2,R}\otimes 1\bigr)
=
m^{\mathrm{WW}}_{\eta_1,\zeta_{23},R}
\circ
\bigl(1\otimes m^{\mathrm{WL}}_{\eta_2,\alpha,R}\bigr)
\]
as morphisms on the retained compactly supported reduced
vanishing-cycle complexes of the ordered inputs.  Equivalently,
\[
\mathfrak o^{\mathrm{assoc}}_{\mathrm{WWL},R}=0 .
\]
The wrapped/wrapped base-change, projection-formula, and reduced
Thom--Sebastiani rows are hypotheses of this proposition; they are not
supplied by the wrapped/wrapped target-admissibility theorem.
\end{proposition}

\begin{proof}
The common refinement is the retained \(WWL\) flag stack
\[
0\subset C_1\subset C_2\subset C_3,
\qquad
C_1/C_0\in W_{\eta_1},\quad
C_2/C_1\in W_{\eta_2},\quad
C_3/C_2\in L_\alpha .
\]
The left comparison \(\lambda^{\mathrm{WWL}}\) records first the
wrapped/wrapped extension of \(\eta_1\) by \(\eta_2\), then the
wrapped/local extension of \(\zeta_{12}\) by \(\alpha\).  The right
comparison \(\rho^{\mathrm{WWL}}\) records first the wrapped/local
extension of \(\eta_2\) by \(\alpha\), then the wrapped/wrapped
extension of \(\eta_1\) by \(\zeta_{23}\).  Both descriptions have
terminal wrapped colour \(\zeta_{123}\).

The wrapped/wrapped and mixed base-change isomorphisms identify the
two iterated compact-support pushforwards on the common flag stack.
The corresponding projection formulas move the ordered source
coefficients through the same pushforward, and the reduced
Thom--Sebastiani transports identify the two coefficient systems.
Therefore the two composites agree, and
\(\mathfrak o^{\mathrm{assoc}}_{\mathrm{WWL},R}=0\).

The proof is wordwise and conditional on the wrapped/wrapped
functorial coefficient rows.  It proves no associativity statement for
\(WWW\), and it does not prove the four-input pentagon, quotient
descent, transition compatibility, or aggregate hybrid-carrier
population.
\end{proof}

The packet
\[
\texttt{certificates/hybrid/wwl\_associativity}
\]
verifies
\(\texttt{WWL\_ASSOCIATIVITY\_CONDITIONAL\_VERIFIED}\):
Proposition~\ref{prop:wwl-associativity} proves the \(WWL\) word
associativity defect vanishes under the stated \(WW/WL\) flag-stack,
mixed functorial rows, and wrapped/wrapped functorial hypotheses.  It
does not prove the remaining \(WWW\) word associativity row or the
four-input pentagon.

\begin{proposition}[Associativity for the \(WWW\) word]
\label{prop:www-associativity}
\textnormal{[proved conditionally for the wrapped/wrapped flag stack]}
Let \(\eta_1,\eta_2,\eta_3\) be retained wrapped colours at height
\(R\), ordered as \((\eta_1,\eta_2,\eta_3)\).  Write
\[
\zeta_{12}=\eta_1+\eta_2,\qquad
\zeta_{23}=\eta_2+\eta_3,\qquad
\zeta_{123}=\eta_1+\eta_2+\eta_3 .
\]
Assume the \(WWW\) flag stack
\(\mathfrak F^{(2),\mathrm{WWW}}_{\eta_1,\eta_2,\eta_3,R}\) of
Proposition~\ref{prop:eight-two-step-flag-stacks}, the \(WW\)
target-admissible correspondence rows, the chosen compact-support
pushforward model, and the required wrapped/wrapped base-change,
projection-formula, and reduced Thom--Sebastiani coefficient
identifications on the \(WWW\) flag square.  Then
\[
m^{\mathrm{WW}}_{\zeta_{12},\eta_3,R}
\circ
\bigl(m^{\mathrm{WW}}_{\eta_1,\eta_2,R}\otimes 1\bigr)
=
m^{\mathrm{WW}}_{\eta_1,\zeta_{23},R}
\circ
\bigl(1\otimes m^{\mathrm{WW}}_{\eta_2,\eta_3,R}\bigr)
\]
as morphisms on the retained compactly supported reduced
vanishing-cycle complexes of the ordered wrapped inputs.  Equivalently,
\[
\mathfrak o^{\mathrm{assoc}}_{\mathrm{WWW},R}=0 .
\]
The wrapped/wrapped base-change, projection-formula, and reduced
Thom--Sebastiani rows are hypotheses of this proposition; they are not
supplied by the wrapped/wrapped target-admissibility theorem.
\end{proposition}

\begin{proof}
The common refinement is the retained \(WWW\) flag stack
\[
0\subset C_1\subset C_2\subset C_3,
\qquad
C_1/C_0\in W_{\eta_1},\quad
C_2/C_1\in W_{\eta_2},\quad
C_3/C_2\in W_{\eta_3}.
\]
The left comparison \(\lambda^{\mathrm{WWW}}\) records first the
wrapped/wrapped extension of \(\eta_1\) by \(\eta_2\), then the
wrapped/wrapped extension of \(\zeta_{12}\) by \(\eta_3\).  The right
comparison \(\rho^{\mathrm{WWW}}\) records first the wrapped/wrapped
extension of \(\eta_2\) by \(\eta_3\), then the wrapped/wrapped
extension of \(\eta_1\) by \(\zeta_{23}\).  Both descriptions have
terminal wrapped colour \(\zeta_{123}\).

The wrapped/wrapped base-change isomorphisms identify the two iterated
compact-support pushforwards on the common flag stack.  The
wrapped/wrapped projection formulas move the ordered source
coefficients through the same pushforward, and the wrapped/wrapped
reduced Thom--Sebastiani transports identify the two coefficient
systems.  Therefore the two composites agree, and
\(\mathfrak o^{\mathrm{assoc}}_{\mathrm{WWW},R}=0\).

The proof is wordwise and conditional on the wrapped/wrapped
functorial coefficient rows.  It proves the last length-three
wordwise associativity defect, but it does not prove the four-input
pentagon, quotient descent, transition compatibility, or aggregate
hybrid-carrier population.
\end{proof}

The packet
\[
\texttt{certificates/hybrid/www\_associativity}
\]
verifies
\(\texttt{WWW\_ASSOCIATIVITY\_CONDITIONAL\_VERIFIED}\):
Proposition~\ref{prop:www-associativity} proves the \(WWW\) word
associativity defect vanishes under the stated \(WW/WW\) flag-stack
and wrapped/wrapped functorial hypotheses.  It completes the
length-three wordwise associativity list, but it does not prove the
four-input pentagon.

\begin{proposition}[Four-input pentagon coherence]
\label{prop:four-input-pentagon-coherence}
\textnormal{[proved conditionally before the reduced \(E\)-quotient]}
Let
\[
u=\epsilon_1\epsilon_2\epsilon_3\epsilon_4\in\{\mathrm L,\mathrm W\}^4
\]
and let \((\xi_1,\xi_2,\xi_3,\xi_4)\) be retained ordered colours of
these types at height \(R\).  Assume the binary correspondence rows
for \(LL,LW,WL,WW\), the compact-support pushforward model, the
length-three associativity rows of
Propositions~\ref{prop:lll-associativity}--\ref{prop:www-associativity},
and the required local, mixed, and wrapped/wrapped base-change,
projection-formula, and reduced Thom--Sebastiani coefficient
identifications.  Assume also the BBDJS--Joyce--Upmeier
Thom--Sebastiani orientation pentagon on the retained four-step
d-critical charts.  Then there is a retained finite-type flag-Quot
stack
\[
\mathfrak F^{(3),u}_{\xi_1,\xi_2,\xi_3,\xi_4,R}
\]
classifying filtrations
\[
0=C_0\subset C_1\subset C_2\subset C_3\subset C_4
\]
with \(C_i/C_{i-1}\) of type \(\epsilon_i\) and colour \(\xi_i\),
all intermediate colours retained, and wrapped rigidifications and
anchors remembered before the \(E\)-quotient.  The five planar binary
parenthesizations of the four inputs give five iterated
correspondence fibre products receiving canonical maps from
\(\mathfrak F^{(3),u}_{\xi_1,\xi_2,\xi_3,\xi_4,R}\).  The two
composites around the Mac Lane pentagon of associators agree:
\[
a_{\xi_1,\xi_2,\xi_3+\xi_4,R}\,
a_{\xi_1+\xi_2,\xi_3,\xi_4,R}
=
(1\otimes a_{\xi_2,\xi_3,\xi_4,R})\,
a_{\xi_1,\xi_2+\xi_3,\xi_4,R}\,
(a_{\xi_1,\xi_2,\xi_3,R}\otimes1).
\]
Equivalently, for every \(u\in\{\mathrm L,\mathrm W\}^4\),
\[
\mathfrak o^{\mathrm{pent}}_{u,R}=0 .
\]
\end{proposition}

\begin{proof}
Fix \(u\) and the ordered colours.  Over the retained terminal object
stack of colour \(\xi_1+\xi_2+\xi_3+\xi_4\), take the relative
three-step flag-Quot stack parametrising the flag
\[
0=C_0\subset C_1\subset C_2\subset C_3\subset C_4
\]
with the prescribed four graded pieces.  The retained closed endpoint
conditions, universal complexes, and extension-closure datum cut out
a finite-type retained substack.  Every intermediate object
\(C_j/C_i\) carries the type forced by the product rule
\(\mathrm L\cdot\mathrm L=\mathrm L\) and otherwise
\(\mathrm W\); if the type is \(\mathrm W\), the wrapped
rigidification and anchor memory are part of the retained datum.  This
constructs \(\mathfrak F^{(3),u}_{\xi_1,\xi_2,\xi_3,\xi_4,R}\)
before quotienting by \(E\).

Each planar binary tree on four leaves forgets the same flag in a
different order and therefore gives a map from the common four-step
flag stack to the corresponding iterated binary correspondence fibre
product.  The associator attached to an internal triple is the
two-step flag comparison of
Propositions~\ref{prop:lll-associativity}--\ref{prop:www-associativity}
pulled back along the appropriate contraction of the four-step flag.
Thus both paths around the pentagon are represented on the same
source stack
\(\mathfrak F^{(3),u}_{\xi_1,\xi_2,\xi_3,\xi_4,R}\), with the same
terminal target colour.

Proper base change for the retained compact-support model identifies
the target pushforwards after pullback to the common four-step flag
stack.  The projection formula moves the four ordered source
coefficients through the same pushforward.  The reduced
Thom--Sebastiani coefficient systems on the two paths are the two
parenthesizations of the same fourfold external product.  The
BBDJS Thom--Sebastiani pentagon
\cite[Theorem~5.10(iv)]{BBDJS2015}, together with the
Joyce--Upmeier multiplicative orientation cocycle
\cite[\S4]{JoyceUpmeier}, identifies these two coefficient systems.
Hence the two pull-push functors around the pentagon agree, and
\(\mathfrak o^{\mathrm{pent}}_{u,R}=0\).

The proof is before the reduced \(E\)-quotient and at fixed height
\(R\).  It proves no quotient descent, transition compatibility,
unit compatibility, higher-coloured tree category, or aggregate
hybrid-carrier population.  Whenever a \(WW\) binary correspondence
appears in the four-input word, the wrapped/wrapped base-change,
projection-formula, and reduced Thom--Sebastiani rows remain
hypotheses, as in
Propositions~\ref{prop:lww-associativity}--\ref{prop:www-associativity}.
\end{proof}

The packet
\[
\texttt{certificates/hybrid/four\_input\_pentagon\_coherence}
\]
verifies
\(\texttt{FOUR\_INPUT\_PENTAGON\_CONDITIONAL\_VERIFIED}\):
Proposition~\ref{prop:four-input-pentagon-coherence} constructs the
retained four-step flag stack for each word in
\(\{\mathrm L,\mathrm W\}^4\) and proves the associator pentagon on
that common refinement under the stated functorial hypotheses.  It
does not define the higher-coloured tree category or prove quotient
descent, transition compatibility, or aggregate population.

\begin{definition}[Higher-coloured hybrid tree category]
\label{def:higher-coloured-hybrid-tree-category}
Fix \(R\).  The nonunital higher-coloured hybrid tree category
\[
\mathsf{Tree}^{E,\mathrm{hyb}}_R
\]
is the free coloured category generated by the retained local and
wrapped profiles below.

An object is a finite ordered hybrid profile
\[
\underline\xi=((\epsilon_i,\xi_i,\mathfrak s_i))_{i=1}^n,
\qquad \epsilon_i\in\{\mathrm L,\mathrm W\},
\]
where \(\xi_i\) is a retained colour of the indicated type.  If
\(\epsilon_i=\mathrm L\), the symbol \(\mathfrak s_i\) is a retained
closed local configuration chart, including its finite indexing set
and diagonal-refinement datum.  If \(\epsilon_i=\mathrm W\), the
symbol \(\mathfrak s_i\) is the retained wrapped rigidification with
source anchor memory before the \(E\)-quotient.  The output colour is
\[
|\underline\xi|=\sum_i\xi_i,
\]
with output type \(\mathrm L\) only for an all-local profile and
\(\mathrm W\) otherwise.

A generating morphism
\[
T:\underline\xi\longrightarrow |\underline\xi|
\]
is a planar rooted tree with leaves ordered by \(\underline\xi\).
Every internal edge carries the retained intermediate colour given by
the sum of the leaves below that edge and the type rule
\(\mathrm L\cdot\mathrm L=\mathrm L\), otherwise \(\mathrm W\).  Every
internal vertex of arity \(k\ge2\) is decorated by the retained
extension correspondence or flag-Quot stack for the ordered subprofile
below that vertex.  Local vertices remember closed-configuration
refinements and symmetric relabellings; mixed and wrapped vertices
remember source, intermediate, and target anchors before quotient.
Composition is grafting of planar rooted trees followed by contraction
of unary identity edges.

The defect of realizing this category by compactly supported reduced
vanishing-cycle pull-push functors is the higher-coloured residual
\[
\mathfrak o^{\mathrm{col}}_R
=\bigl(
\mathfrak o^{\mathrm{tree}}_R,\,
\mathfrak o^{\mathrm{unit}}_R,\,
\mathfrak o^{\mathrm{sym}}_R,\,
\mathfrak o^{\mathrm{ref}}_R,\,
\mathfrak o^{\mathrm{des}}_R,\,
\mathfrak o^{\mathrm{ov}}_R
\bigr).
\]
The first component measures compatibility of tree contractions beyond
the four-input pentagon.  The remaining components measure,
respectively, vacuum units, symmetric relabelling of local
configurations, refinement of closed local charts, descent from
ordered charts, and overlap compatibility of retained local and
wrapped charts.  This definition supplies the category and the names
of its residual components; it does not assert
\(\mathfrak o^{\mathrm{col}}_R=0\).
\end{definition}

The packet
\[
\texttt{certificates/hybrid/higher\_coloured\_tree\_category\_definition}
\]
verifies
\(\texttt{HIGHER\_COLOURED\_TREE\_CATEGORY\_DEFINED}\):
Definition~\ref{def:higher-coloured-hybrid-tree-category} defines the
nonunital higher-coloured tree category and the six components of
\(\mathfrak o^{\mathrm{col}}_R\).  It does not prove the residual
vanishes, and it does not supply unit compatibility, symmetric descent,
wrapped order conventions, quotient descent, transition compatibility,
or aggregate population.

\begin{proposition}[Higher-coloured residual vanishing criterion]
\label{prop:higher-coloured-residual-vanishing-criterion}
\textnormal{[conditional finite-stage criterion]}
At fixed height \(R\), assume
Definition~\ref{def:higher-coloured-hybrid-tree-category} and
Proposition~\ref{prop:four-input-pentagon-coherence}.  Suppose the
following finite rows are supplied in the retained
\(E\)-equivariant correspondence category before quotient:
\begin{enumerate}
\item vacuum object and unit correspondences with left and right unit
defect zero for local and wrapped inputs;
\item symmetric relabelling descent for projection-finite local
configuration charts;
\item refinement compatibility for closed local configuration charts;
\item descent from ordered local charts to the symmetric finite-colour
colimit;
\item wrapped insertion order and braid-convention compatibility for
ordered wrapped profiles;
\item overlap compatibility for retained local and wrapped charts on
common refinements.
\end{enumerate}
Then the higher-coloured residual of
Definition~\ref{def:higher-coloured-hybrid-tree-category} vanishes:
\[
\mathfrak o^{\mathrm{col}}_R=0 .
\]
Equivalently, all six components
\[
\mathfrak o^{\mathrm{tree}}_R,\quad
\mathfrak o^{\mathrm{unit}}_R,\quad
\mathfrak o^{\mathrm{sym}}_R,\quad
\mathfrak o^{\mathrm{ref}}_R,\quad
\mathfrak o^{\mathrm{des}}_R,\quad
\mathfrak o^{\mathrm{ov}}_R
\]
have defect rank zero.
\end{proposition}

\begin{proof}
The tree-contraction component is the coherence problem for binary
parenthesizations in
\(\mathsf{Tree}^{E,\mathrm{hyb}}_R\).  The length-three associators
are the rows of
Propositions~\ref{prop:lll-associativity}--\ref{prop:www-associativity},
and their four-input pentagon is
Proposition~\ref{prop:four-input-pentagon-coherence}.  Hence the
standard coherence theorem for parenthesized products identifies any
two binary contraction paths with the same planar leaves by a composite
of four-input pentagons.  Pulling those pentagons to the retained
common flag-Quot refinement gives
\(\mathfrak o^{\mathrm{tree}}_R=0\).

The remaining five components are not consequences of the pentagon.
By the supplied rows in the statement, the left and right vacuum
correspondences act as units on local and wrapped inputs, so
\(\mathfrak o^{\mathrm{unit}}_R=0\).  Symmetric relabelling descent
sets \(\mathfrak o^{\mathrm{sym}}_R=0\).  Refinement compatibility of
closed local charts sets \(\mathfrak o^{\mathrm{ref}}_R=0\).  Descent
from ordered charts to the symmetric finite-colour colimit sets
\(\mathfrak o^{\mathrm{des}}_R=0\).  Overlap compatibility on common
local/wrapped refinements sets \(\mathfrak o^{\mathrm{ov}}_R=0\).
The residual is the direct tuple of these six components; therefore
\(\mathfrak o^{\mathrm{col}}_R=0\).

This criterion does not construct the non-tree input rows.  In the
present correction ledger the unit, symmetric-relabelling, and wrapped
order rows are supplied separately; refinement, ordered-chart descent,
and overlap rows remain separate obligations needed for aggregate
population.
\end{proof}

The packet
\[
\texttt{certificates/hybrid/higher\_coloured\_residual\_vanishing\_criterion}
\]
verifies
\(\texttt{HIGHER\_COLOURED\_RESIDUAL\_CONDITIONAL\_CRITERION\_VERIFIED}\):
Proposition~\ref{prop:higher-coloured-residual-vanishing-criterion}
proves the finite-stage implication from the six named component rows
to \(\mathfrak o^{\mathrm{col}}_R=0\).  It proves the tree-contraction
component from the associativity rows and four-input pentagon, but it
does not supply the refinement, ordered-chart descent, or overlap rows
required to invoke the criterion unconditionally.  The unit,
symmetric-relabelling, and wrapped order rows are supplied by separate
packets.

\begin{definition}[Hybrid unit object and vacuum correspondences]
\label{def:hybrid-unit-object-vacuum-correspondences}
\textnormal{[definition only]}
Fix a finite HN height \(R\).  A hybrid unit datum consists of the
following retained rows before the reduced \(E\)-quotient.

First, the zero object of the retained Hall heart determines a local
zero colour
\[
0_R\in\Gamma_R^{\mathrm{loc}},\qquad b_R^{\mathrm{geom}}(0_R)=0,
\]
with empty local support chart \(\varnothing\).  The corresponding
unit stack is the one-point closed local incidence stack
\[
\mathbf 1^{\mathrm{hyb}}_R
:=
\mathfrak M^{\mathrm{loc,cl}}_{0_R,R,\varnothing}.
\]
The same symbol \(0_R:\varnothing\to\Gamma_R^{\mathrm{loc}}\) denotes
the empty local colour profile whose total colour is \(0_R\).
There is no wrapped zero colour: the unit is local because the wrapped
sector is defined by \(b_R^{\mathrm{geom}}>0\).

Second, for every finite local colour map
\(\alpha:I\to\Gamma_R^{\mathrm{loc}}\), the datum contains the split
local unit substacks
\[
\mathfrak U^{\mathrm{LL},\ell}_{0_R,\alpha;\alpha,R,I}
\subset
\mathfrak E^{\mathrm{LL}}_{0_R,\alpha;\alpha,R,\varnothing,I,I},
\qquad
\mathfrak U^{\mathrm{LL},r}_{\alpha,0_R;\alpha,R,I}
\subset
\mathfrak E^{\mathrm{LL}}_{\alpha,0_R;\alpha,R,I,\varnothing,I}.
\]
They classify, respectively, the split exact sequences
\[
0\longrightarrow A_{\alpha,I}\longrightarrow A_{\alpha,I}
\longrightarrow 0\longrightarrow0,
\qquad
0\longrightarrow0\longrightarrow A_{\alpha,I}
\longrightarrow A_{\alpha,I}\longrightarrow0 .
\]
The source maps remember the ordered pair
\((\mathbf 1^{\mathrm{hyb}}_R,A_{\alpha,I})\) or
\((A_{\alpha,I},\mathbf 1^{\mathrm{hyb}}_R)\), and the target maps
send the split sequence to \(A_{\alpha,I}\).

Third, for every wrapped colour \(\eta\in\Gamma_R^{\mathrm{wr}}\), the
datum contains the split mixed unit substacks
\[
\mathfrak U^{\mathrm{LW},\ell}_{0_R,\eta;\eta,R}
\subset
\mathfrak E^{\mathrm{LW}}_{0_R,\eta;\eta,R,\varnothing},
\qquad
\mathfrak U^{\mathrm{WL},r}_{\eta,0_R;\eta,R}
\subset
\mathfrak E^{\mathrm{WL}}_{\eta,0_R;\eta,R,\varnothing}.
\]
They classify the split exact sequences
\[
0\longrightarrow W_\eta\longrightarrow W_\eta
\longrightarrow0\longrightarrow0,
\qquad
0\longrightarrow0\longrightarrow W_\eta
\longrightarrow W_\eta\longrightarrow0 .
\]
The wrapped rigidification and source/target anchor memory are those
of \(W_\eta\); the zero local factor carries the empty support chart
and no anchor.  These rows are the unit correspondences used by
\(\mathfrak o^{\mathrm{unit}}_R\).  This definition does not prove
that the induced compact-support pull-push maps are the identity on
local or wrapped coefficients; those are the local and wrapped unit
compatibility rows.
\end{definition}

The packet
\[
\texttt{certificates/hybrid/hybrid\_unit\_object\_and\_vacuum\_correspondences}
\]
verifies
\(\texttt{HYBRID\_UNIT\_OBJECT\_AND\_VACUUM\_CORRESPONDENCES\_DEFINED}\):
Definition~\ref{def:hybrid-unit-object-vacuum-correspondences} defines
the local zero unit object and the four split vacuum correspondence
types \(LL_\ell\), \(LL_r\), \(LW_\ell\), and \(WL_r\).  It does not
prove the local or wrapped unit identity laws, quotient descent,
transition compatibility, or aggregate population.

\begin{proposition}[Local unit compatibility for the hybrid product]
\label{prop:local-unit-compatibility}
\textnormal{[proved before quotient at finite height]}
Assume the hybrid unit datum of
Definition~\ref{def:hybrid-unit-object-vacuum-correspondences}, the
local/local correspondence datum of
Definition~\ref{def:local-local-correspondence-datum}, the
target-properness row of
Proposition~\ref{prop:local-local-target-properness}, and the
compact-support model of
Definition~\ref{def:compact-support-exceptional-pushforward-model}.
Let
\[
K^{\mathrm{loc}}_{\alpha,R,I}
:=
\Phi^{\mathrm{loc}}_{\alpha,R,I}\otimes
o^{\mathrm{loc}}_{\alpha,R,I}
\]
be the reduced local coefficient on
\(\mathfrak M^{\mathrm{loc,cl}}_{\alpha,R,I}\), and fix the standard
zero-object normalization
\[
K^{\mathrm{loc}}_{0_R,R,\varnothing}
\simeq
\C_{\mathbf 1^{\mathrm{hyb}}_R}.
\]
Let \(p^{\mathrm{LL},\ell},q^{\mathrm{LL},\ell}\) and
\(p^{\mathrm{LL},r},q^{\mathrm{LL},r}\) denote the restrictions of
the \(LL\) source and target maps to
\(\mathfrak U^{\mathrm{LL},\ell}_{0_R,\alpha;\alpha,R,I}\) and
\(\mathfrak U^{\mathrm{LL},r}_{\alpha,0_R;\alpha,R,I}\).
Then the compact-support pull-push maps attached to the split local
unit correspondences satisfy
\[
\mu^{\mathrm{LL},\ell}_{0_R,\alpha;R,I}
=
(q^{\mathrm{LL},\ell})^{\mathrm{cs}}_!\,
\operatorname{TS}^{\mathrm{red}}_{0_R,\alpha}\,
(p^{\mathrm{LL},\ell})^*
=
\operatorname{id}_{K^{\mathrm{loc}}_{\alpha,R,I}},
\]
and
\[
\mu^{\mathrm{LL},r}_{\alpha,0_R;R,I}
=
(q^{\mathrm{LL},r})^{\mathrm{cs}}_!\,
\operatorname{TS}^{\mathrm{red}}_{\alpha,0_R}\,
(p^{\mathrm{LL},r})^*
=
\operatorname{id}_{K^{\mathrm{loc}}_{\alpha,R,I}}.
\]
Equivalently, the local unit defect component has defect rank zero:
\[
\mathfrak o^{\mathrm{unit,loc}}_R=0 .
\]
\end{proposition}

\begin{proof}
For the left unit row, the split substack
\[
\mathfrak U^{\mathrm{LL},\ell}_{0_R,\alpha;\alpha,R,I}
\]
is isomorphic to
\(\mathfrak M^{\mathrm{loc,cl}}_{\alpha,R,I}\).  The isomorphism sends
\(A_{\alpha,I}\) to the split sequence
\[
0\longrightarrow A_{\alpha,I}\longrightarrow A_{\alpha,I}
\longrightarrow0\longrightarrow0,
\]
and its inverse takes the target object of the sequence.  Under this
identification the source map is the closed insertion
\[
A_{\alpha,I}\longmapsto
(\mathbf 1^{\mathrm{hyb}}_R,A_{\alpha,I}),
\]
and the target map is the identity on
\(\mathfrak M^{\mathrm{loc,cl}}_{\alpha,R,I}\).  Since the target map
is proper, the compact-support model uses the identity
compactification and \((q^{\mathrm{LL},\ell})^{\mathrm{cs}}_!\) is
ordinary pushforward along the identity.

The pullback of the ordered source coefficient is
\[
(p^{\mathrm{LL},\ell})^*
\bigl(K^{\mathrm{loc}}_{0_R,R,\varnothing}
\boxtimes K^{\mathrm{loc}}_{\alpha,R,I}\bigr)
\simeq
K^{\mathrm{loc}}_{\alpha,R,I}.
\]
The zero-object normalization identifies the first factor with the
constant unit complex on the point.  The BBDJS Thom--Sebastiani
isomorphism for the sum of a zero potential and the local potential
\cite[Theorem~5.10]{BBDJS2015}, tensored with the Joyce--Upmeier
orientation unit \cite[\S4]{JoyceUpmeier}, is the identity on this
coefficient.  Hence
\[
(q^{\mathrm{LL},\ell})^{\mathrm{cs}}_!\,
\operatorname{TS}^{\mathrm{red}}_{0_R,\alpha}\,
(p^{\mathrm{LL},\ell})^*
=\operatorname{id}_{K^{\mathrm{loc}}_{\alpha,R,I}}.
\]

The right unit row is identical with the factors reversed.  The
substack
\(\mathfrak U^{\mathrm{LL},r}_{\alpha,0_R;\alpha,R,I}\) is the graph
of the split sequence
\[
0\longrightarrow0\longrightarrow A_{\alpha,I}
\longrightarrow A_{\alpha,I}\longrightarrow0,
\]
so its source map inserts
\((A_{\alpha,I},\mathbf 1^{\mathrm{hyb}}_R)\) and its target map is
again the identity.  The ordered Thom--Sebastiani isomorphism for the
local potential plus the zero potential is the right-unit form of the
same zero-factor identity.  Thus the right pull-push is
\(\operatorname{id}_{K^{\mathrm{loc}}_{\alpha,R,I}}\).

The proof is before the reduced \(E\)-quotient and at fixed height
\(R\).  It proves no wrapped unit identity, quotient descent,
transition compatibility, symmetric descent, refinement compatibility,
or aggregate hybrid-carrier population.
\end{proof}

The packet
\[
\texttt{certificates/hybrid/local\_unit\_compatibility}
\]
verifies
\(\texttt{LOCAL\_UNIT\_COMPATIBILITY\_VERIFIED}\):
Proposition~\ref{prop:local-unit-compatibility} proves the left and
right \(LL\) vacuum pull-push maps are the identity on local
coefficients.  It does not prove wrapped unit compatibility, quotient
descent, transition compatibility, or aggregate population.

\begin{proposition}[Wrapped unit compatibility for the hybrid product]
\label{prop:wrapped-unit-compatibility}
\textnormal{[proved before quotient at finite height]}
Assume the hybrid unit datum of
Definition~\ref{def:hybrid-unit-object-vacuum-correspondences}, the
one-sided mixed correspondence datum of
Definition~\ref{def:mixed-local-wrapped-correspondence-datum}, the
mixed target-admissibility row of
Proposition~\ref{prop:mixed-local-wrapped-target-admissibility}, the
compact-support model of
Definition~\ref{def:compact-support-exceptional-pushforward-model},
and the mixed Thom--Sebastiani transport row of
Proposition~\ref{prop:mixed-correspondence-thom-sebastiani}.  Let
\[
K^{\mathrm{wr}}_{\eta,R}
:=
\Phi^{\mathrm{wr}}_{\eta,R}\otimes o^{\mathrm{wr}}_{\eta,R}
\]
be the reduced wrapped coefficient on
\(\mathcal M^{\mathrm{wr,rig}}_{\eta,R}\), and keep the zero-object
normalization
\[
K^{\mathrm{loc}}_{0_R,R,\varnothing}
\simeq
\C_{\mathbf 1^{\mathrm{hyb}}_R}.
\]
Let \(p^{\mathrm{LW},\ell},q^{\mathrm{LW},\ell}\) and
\(p^{\mathrm{WL},r},q^{\mathrm{WL},r}\) denote the restrictions of the
mixed source and target maps to
\(\mathfrak U^{\mathrm{LW},\ell}_{0_R,\eta;\eta,R}\) and
\(\mathfrak U^{\mathrm{WL},r}_{\eta,0_R;\eta,R}\).  Then the
compact-support pull-push maps attached to the split wrapped unit
correspondences satisfy
\[
\mu^{\mathrm{LW},\ell}_{0_R,\eta;R}
=
(q^{\mathrm{LW},\ell})^{\mathrm{cs}}_!\,
\operatorname{TS}^{\mathrm{red},\mathrm{LW}}_{0_R,\eta}\,
(p^{\mathrm{LW},\ell})^*
=
\operatorname{id}_{K^{\mathrm{wr}}_{\eta,R}},
\]
and
\[
\mu^{\mathrm{WL},r}_{\eta,0_R;R}
=
(q^{\mathrm{WL},r})^{\mathrm{cs}}_!\,
\operatorname{TS}^{\mathrm{red},\mathrm{WL}}_{\eta,0_R}\,
(p^{\mathrm{WL},r})^*
=
\operatorname{id}_{K^{\mathrm{wr}}_{\eta,R}}.
\]
Equivalently, the wrapped unit defect component has defect rank zero:
\[
\mathfrak o^{\mathrm{unit,wr}}_R=0 .
\]
\end{proposition}

\begin{proof}
For the left wrapped unit row, the split substack
\[
\mathfrak U^{\mathrm{LW},\ell}_{0_R,\eta;\eta,R}
\]
is isomorphic to \(\mathcal M^{\mathrm{wr,rig}}_{\eta,R}\).  The
isomorphism sends \(W_\eta\) to the split sequence
\[
0\longrightarrow W_\eta\longrightarrow W_\eta
\longrightarrow0\longrightarrow0,
\]
and its inverse takes the target wrapped object.  Under this
identification the source map is the insertion
\[
W_\eta\longmapsto
(\mathbf 1^{\mathrm{hyb}}_R,W_\eta),
\]
and the target map is the identity on
\(\mathcal M^{\mathrm{wr,rig}}_{\eta,R}\).  The unit definition
records that the zero local factor carries no anchor; hence the source
and target anchor memory on the split row is exactly the anchor memory
of \(W_\eta\).  No \(E\)-quotient is taken.  Since the target map is
proper, the compact-support model uses the identity compactification
and \((q^{\mathrm{LW},\ell})^{\mathrm{cs}}_!\) is ordinary pushforward
along the identity.

The pullback of the ordered source coefficient is
\[
(p^{\mathrm{LW},\ell})^*
\bigl(K^{\mathrm{loc}}_{0_R,R,\varnothing}
\boxtimes K^{\mathrm{wr}}_{\eta,R}\bigr)
\simeq
K^{\mathrm{wr}}_{\eta,R}.
\]
The zero-object normalization identifies the first factor with the
constant unit complex on the point.  The mixed Thom--Sebastiani
transport of Proposition~\ref{prop:mixed-correspondence-thom-sebastiani}
specializes, for the zero local potential plus the wrapped potential,
to the BBDJS zero-factor Thom--Sebastiani identity
\cite[Theorem~5.10]{BBDJS2015}, tensored with the Joyce--Upmeier
orientation unit \cite[\S4]{JoyceUpmeier}.  Therefore
\[
(q^{\mathrm{LW},\ell})^{\mathrm{cs}}_!\,
\operatorname{TS}^{\mathrm{red},\mathrm{LW}}_{0_R,\eta}\,
(p^{\mathrm{LW},\ell})^*
=\operatorname{id}_{K^{\mathrm{wr}}_{\eta,R}}.
\]

The right wrapped unit row is the same argument with the factors
reversed.  The split substack
\(\mathfrak U^{\mathrm{WL},r}_{\eta,0_R;\eta,R}\) is the graph of
\[
0\longrightarrow0\longrightarrow W_\eta
\longrightarrow W_\eta\longrightarrow0,
\]
the source map inserts
\((W_\eta,\mathbf 1^{\mathrm{hyb}}_R)\), the target map is the
identity on \(\mathcal M^{\mathrm{wr,rig}}_{\eta,R}\), and the wrapped
anchor memory is unchanged.  The ordered \(WL\) Thom--Sebastiani
transport is the right zero-factor identity.  Hence the right
pull-push is \(\operatorname{id}_{K^{\mathrm{wr}}_{\eta,R}}\).

The proof is before the reduced \(E\)-quotient and at fixed height
\(R\).  It proves no quotient descent, transition compatibility,
symmetric descent, refinement compatibility, overlap compatibility, or
aggregate hybrid-carrier population.
\end{proof}

The packet
\[
\texttt{certificates/hybrid/wrapped\_unit\_compatibility}
\]
verifies
\(\texttt{WRAPPED\_UNIT\_COMPATIBILITY\_VERIFIED}\):
Proposition~\ref{prop:wrapped-unit-compatibility} proves the left
\(LW\) and right \(WL\) vacuum pull-push maps are the identity on
wrapped coefficients, with wrapped anchor memory unchanged before the
reduced \(E\)-quotient.  It does not prove quotient descent,
transition compatibility, or aggregate population.

\begin{proposition}[Symmetric descent for local configuration charts]
\label{prop:local-configuration-symmetric-descent}
\textnormal{[proved before quotient at finite height]}
Assume the local \(b=0\) carrier of
Definition~\ref{def:hybrid-local-stratum-prestack} and the
projection-finite closed local configuration sector of the finite
stage.  Let
\(\alpha:I\to\Gamma_R^{\mathrm{loc}}\) be a finite local colour map
and let \(\sigma:I\xrightarrow{\sim}I'\) be a bijection.  Define
\[
(\sigma_*\alpha)_{i'}=\alpha_{\sigma^{-1}(i')},
\qquad
\rho_\sigma:E^I\xrightarrow{\sim}E^{I'},
\qquad
(\rho_\sigma(\mathbf e))_{i'}=e_{\sigma^{-1}(i')}.
\]
Then
\[
(\operatorname{id}_E\times\rho_\sigma)(Z_I)=Z_{I'}
\]
and \(\rho_\sigma\) induces an isomorphism of closed local
configuration charts
\[
r_\sigma:
\mathfrak M^{\mathrm{loc,cl}}_{\alpha,R,I}
\xrightarrow{\sim}
\mathfrak M^{\mathrm{loc,cl}}_{\sigma_*\alpha,R,I'},
\qquad
(A,\mathbf e)\longmapsto (A,\rho_\sigma(\mathbf e)),
\]
with
\[
\pi_{\sigma_*\alpha,R,I'}\circ r_\sigma
=
\rho_\sigma\circ\pi_{\alpha,R,I}.
\]
For the reduced local coefficient
\[
K^{\mathrm{loc}}_{\alpha,R,I}
:=
\Phi^{\mathrm{loc}}_{\alpha,R,I}\otimes
o^{\mathrm{loc}}_{\alpha,R,I},
\]
there is a functorial relabelling isomorphism
\[
r_\sigma^*K^{\mathrm{loc}}_{\sigma_*\alpha,R,I'}
\simeq
K^{\mathrm{loc}}_{\alpha,R,I},
\]
and these isomorphisms satisfy the cocycle identity for composable
bijections.  Hence the stabilizer
\[
\mathfrak S_\alpha=\{\sigma\in\operatorname{Aut}(I)\mid
\sigma_*\alpha=\alpha\}
\]
acts on
\((\mathfrak M^{\mathrm{loc,cl}}_{\alpha,R,I},
K^{\mathrm{loc}}_{\alpha,R,I})\), and the coefficient descends
effectively to the quotient stack
\[
[\mathfrak M^{\mathrm{loc,cl}}_{\alpha,R,I}/\mathfrak S_\alpha].
\]
Equivalently, the fixed-chart symmetric relabelling component of the
higher-coloured residual has defect rank zero:
\[
\mathfrak o^{\mathrm{sym}}_R=0
\]
on projection-finite local configuration charts before quotient.
\end{proposition}

\begin{proof}
The incidence locus \(Z_I\subset E\times E^I\) is the union of the
graphs \(e=e_i\).  Under
\(\operatorname{id}_E\times\rho_\sigma\) the graph indexed by
\(i\in I\) is sent to the graph indexed by
\(\sigma(i)\in I'\); hence
\((\operatorname{id}_E\times\rho_\sigma)(Z_I)=Z_{I'}\).
Thus the condition
\[
p_E(\operatorname{Supp}A)\subset Z_I(\mathbf e)
\]
is equivalent to
\[
p_E(\operatorname{Supp}A)\subset Z_{I'}(\rho_\sigma(\mathbf e)).
\]
Relabelling the finite colour profile by \(\sigma\) changes no
underlying sheaf \(A\) and preserves the total retained local colour.
Therefore
\[
(A,\mathbf e)\longmapsto(A,\rho_\sigma(\mathbf e))
\]
defines the stated isomorphism \(r_\sigma\), with inverse
\(r_{\sigma^{-1}}\), and the displayed identity with the configuration
maps follows from the definition of \(\rho_\sigma\).

The local d-critical chart, its potential, and its cosection-reduced
vanishing-cycle complex are attached to the object \(A\) together with
the ordered support coordinates.  The map \(r_\sigma\) is an isomorphism
of these charts: it only relabels the support coordinates and the
matching colour labels.  Functoriality of vanishing cycles under
isomorphism, tensored with the transported Joyce--Upmeier orientation
line, gives
\[
r_\sigma^*
(\Phi^{\mathrm{loc}}_{\sigma_*\alpha,R,I'}
\otimes o^{\mathrm{loc}}_{\sigma_*\alpha,R,I'})
\simeq
\Phi^{\mathrm{loc}}_{\alpha,R,I}\otimes
o^{\mathrm{loc}}_{\alpha,R,I}.
\]
Since \(r_\tau r_\sigma=r_{\tau\sigma}\), these pullback
isomorphisms satisfy the cocycle identity.

For \(\sigma\in\mathfrak S_\alpha\) the target profile is again
\(\alpha\), so the preceding isomorphisms define a finite group action
on the stack and an equivariant structure on the coefficient complex.
Constructible complexes on a quotient stack by a finite group are the
same data as equivariant constructible complexes on the source stack;
therefore the descent to
\([\mathfrak M^{\mathrm{loc,cl}}_{\alpha,R,I}/\mathfrak S_\alpha]\)
is effective.  This proves the symmetric relabelling component for a
fixed projection-finite local chart.

The proof does not prove collision compatibility, refinement of closed
local charts, descent from ordered charts to the symmetric finite-colour
colimit, quotient descent, transition compatibility, or aggregate
hybrid-carrier population.
\end{proof}

The packet
\[
\texttt{certificates/hybrid/local\_configuration\_symmetric\_descent}
\]
verifies
\(\texttt{LOCAL\_CONFIGURATION\_SYMMETRIC\_DESCENT\_VERIFIED}\):
Proposition~\ref{prop:local-configuration-symmetric-descent} proves
finite symmetric-group descent for each projection-finite closed local
configuration chart and its reduced coefficient before the reduced
\(E\)-quotient.  It does not prove collision compatibility, refinement
compatibility, descent from ordered charts to the symmetric
finite-colour colimit, quotient descent, transition compatibility, or
aggregate population.

\begin{proposition}[Wrapped insertion order conventions]
\label{prop:wrapped-insertion-order-conventions}
\textnormal{[proved before quotient at finite height]}
Fix \(R\) and the ordered hybrid tree category of
Definition~\ref{def:higher-coloured-hybrid-tree-category}.  For an
ordered hybrid profile
\[
\underline\xi=((\epsilon_i,\xi_i,\mathfrak s_i))_{i=1}^n
\]
write
\[
\operatorname{Wr}(\underline\xi)
=
(\xi_{i_1},\ldots,\xi_{i_m}),
\qquad
i_1<\cdots<i_m,\qquad \epsilon_{i_a}=\mathrm W,
\]
for the ordered subsequence of wrapped leaves.  The wrapped insertion
convention is the following.
\begin{enumerate}
\item At every binary vertex the wrapped order of the output is the
concatenation of the wrapped orders of the ordered inputs.
\item For a \(WW\) vertex with ordered source pair
\((\eta_1,\eta_2)\), the Hall correspondence is
\[
0\longrightarrow W_{\eta_2}\longrightarrow B_{\eta_1+\eta_2}
\longrightarrow W_{\eta_1}\longrightarrow0,
\]
and the ordered wrapped source list is \((\eta_1,\eta_2)\): first the
quotient colour, then the subobject colour.  The target wrapped colour
remembers this source order as the anchor-memory order before the
reduced \(E\)-quotient.
\item For \(LW\), \(WL\), and \(LWL\) vertices, local factors do not
change the order of the wrapped subsequence.
\item The comparison maps for the eight two-step flag stacks of
Proposition~\ref{prop:eight-two-step-flag-stacks} preserve the ordered
wrapped subsequence \(\operatorname{Wr}\) for each of the words
\[
LLL,\ LLW,\ LWL,\ WLL,\ LWW,\ WLW,\ WWL,\ WWW.
\]
\end{enumerate}
Therefore parenthesization changes in
\(\mathsf{Tree}^{E,\mathrm{hyb}}_R\) use the planar leaf order and
introduce no braid operator and no permutation of wrapped inputs.  The
wrapped-order input row in
Proposition~\ref{prop:higher-coloured-residual-vanishing-criterion}
has defect rank zero at fixed height before quotient.
\end{proposition}

\begin{proof}
The object set of
\(\mathsf{Tree}^{E,\mathrm{hyb}}_R\) is ordered by definition, and a
generating morphism is a planar rooted tree whose leaves are ordered by
the input profile.  Thus every internal vertex sees an ordered
subprofile.  The rule for the output colour is additive, and the type
rule sends an output to the wrapped sector as soon as one input below
the vertex is wrapped.  This gives a canonical ordered wrapped
subsequence: scan the leaves below the vertex from left to right and
discard the local leaves.

For an \(LW\) or \(WL\) vertex there is exactly one wrapped input, so
the ordered wrapped subsequence is unchanged.  For an \(LWL\) vertex
the two local leaves are discarded by the scan, and the wrapped
subsequence is again the single middle wrapped colour.  For a \(WW\)
vertex the ordered source pair in
Definition~\ref{def:wrapped-wrapped-correspondence-datum} is
\((\eta_1,\eta_2)\), while the Hall extension convention is
\[
0\to W_{\eta_2}\to B_{\eta_1+\eta_2}\to W_{\eta_1}\to0.
\]
Hence the first source factor is the quotient and the second source
factor is the subobject.  This fixes the ordered source-anchor list as
\((\eta_1,\eta_2)\).  No transposition is applied when the target
wrapped colour is formed.

It remains to check parenthesization changes.  The two-step flag stack
for a word \(w=\epsilon_1\epsilon_2\epsilon_3\) records a filtration
whose successive quotients are the three ordered leaves.  The left
comparison map contracts the first two leaves and then the resulting
intermediate object with the third leaf.  The right comparison map
contracts the last two leaves and then the first leaf with the
resulting intermediate object.  In both cases the ordered wrapped
subsequence is obtained by scanning the same ordered list
\((\epsilon_1,\epsilon_2,\epsilon_3)\) and discarding local entries.
Thus the comparison maps preserve \(\operatorname{Wr}\) for all eight
words.

Consequently every contraction in the planar tree category preserves
the wrapped subsequence by induction on the number of vertices.  Since
the category has no braid generator and no symmetric identification of
wrapped leaves, a braid word is not a hidden morphism in the
realization.  The only braid compatible with the convention is the
identity braid, which acts by the identity on the ordered source
coefficient.  Hence the wrapped insertion order row has defect rank
zero.

The proof does not prove symmetric descent for unordered wrapped
pairs, closed-chart refinement, descent from ordered charts to the
symmetric finite-colour colimit, overlap compatibility, quotient
descent, transition compatibility, or aggregate hybrid-carrier
population.
\end{proof}

The packet
\[
\texttt{certificates/hybrid/wrapped\_insertion\_order\_conventions}
\]
verifies
\(\texttt{WRAPPED\_INSERTION\_ORDER\_CONVENTIONS\_VERIFIED}\):
Proposition~\ref{prop:wrapped-insertion-order-conventions} fixes the
planar wrapped-source order for mixed and wrapped vertices, checks the
\(WW\) quotient/subobject convention, and proves that the eight
two-step flag comparisons preserve the same ordered wrapped subsequence
before the reduced \(E\)-quotient.  It does not prove unordered
wrapped-pair symmetric descent, ordered-chart descent, overlap
compatibility, quotient descent, transition compatibility, or aggregate
population.

\begin{definition}[Ordered two-sided mixed correspondence datum]
\label{def:two-sided-mixed-correspondence-datum}
Fix a finite HN height \(R\).  Let
\(\alpha_-:I_-\to\Gamma_R^{\mathrm{loc}}\) and
\(\beta_+:I_+\to\Gamma_R^{\mathrm{loc}}\) be finite local colour maps,
and let \(\eta\in\Gamma_R^{\mathrm{wr}}\).  Write
\(|\alpha_-|=\sum_{i\in I_-}\alpha_{-,i}\) and
\(|\beta_+|=\sum_{j\in I_+}\beta_{+,j}\).  A two-sided mixed
correspondence row is defined only for a retained wrapped target colour
\(\zeta\in\Gamma_R^{\mathrm{wr}}\) satisfying
\[
\zeta=|\alpha_-|+\eta+|\beta_+|.
\]
The ordered local--wrapped--local correspondence is the unreduced
diagram
\[
\mathfrak E^{\mathrm{LWL}}_{\alpha_-;\eta;\beta_+,R,I_-,I_+}
\xrightarrow{p^{\mathrm{LWL}}}
\mathfrak M^{\mathrm{loc,cl}}_{\alpha_-,R,I_-}
\times
\mathcal M_{\eta,R}^{\mathrm{wr,rig}}
\times
\mathfrak M^{\mathrm{loc,cl}}_{\beta_+,R,I_+},
\]
\[
\mathfrak E^{\mathrm{LWL}}_{\alpha_-;\eta;\beta_+,R,I_-,I_+}
\xrightarrow{q^{\mathrm{LWL}}}
\mathcal M_{\zeta,R}^{\mathrm{wr,rig}},
\]
classifying retained flags
\[
0\subset A_{\beta_+}\subset B_{\eta+\beta_+}\subset C_\zeta
\]
with associated graded pieces
\[
A_{\beta_+},\qquad W_\eta,\qquad A_{\alpha_-}
\]
from bottom to top.  Equivalently,
\[
0\to A_{\beta_+}\to B_{\eta+\beta_+}\to W_\eta\to0,
\qquad
0\to B_{\eta+\beta_+}\to C_\zeta\to A_{\alpha_-}\to0.
\]
The diagram lives before the reduced \(E\)-quotient.  It is a
three-input correspondence type, not the full eight-word two-step flag
atlas and not an associativity theorem.  Source and target anchor
memory, admissibility, base change, projection formula,
Thom--Sebastiani transport, and transition compatibility are separate
rows.
\end{definition}

The packet
\[
\texttt{certificates/hybrid/two\_sided\_mixed\_correspondence\_definition}
\]
verifies
\(\texttt{TWO\_SIDED\_MIXED\_CORRESPONDENCE\_DEFINITION\_VERIFIED}\):
the ordered \(LWL\) two-sided mixed correspondence type is defined,
while flag-stack rows, source/target-map rows, populated anchor-memory
rows, admissibility rows, base-change rows, projection-formula rows,
Thom--Sebastiani rows, and transition rows are recorded as missing open
obligations.

\begin{definition}[Source and target anchors before quotient]
\label{def:source-target-anchor-memory-datum}
Fix a finite HN height \(R\).  Assume the wrapped prequotient sector
has the prequotient stacks of
Definition~\ref{def:e-equivariant-wrapped-prequotient-stack} and
has the abstract anchors
\[
\lambda_{\eta,R}:\mathcal M_{\eta,R}^{\mathrm{wr,rig}}\longrightarrow E
\]
of Definition~\ref{def:hybrid-wrapped-factorization-datum}.  For every
unreduced correspondence with a wrapped input or wrapped target, the
source/target anchor memory is part of the correspondence datum before
the reduced \(E\)-quotient is applied.

For the one-sided mixed rows of
Definition~\ref{def:mixed-local-wrapped-correspondence-datum}, it is
\[
a^{\mathrm{LW}}_{\alpha,\eta;\zeta,R,I}:
\mathfrak E^{\mathrm{LW}}_{\alpha,\eta;\zeta,R,I}\longrightarrow E^2,
\qquad
\xi\longmapsto
\bigl(\lambda_{\eta,R}(W_\eta),
\lambda_{\zeta,R}(B_\zeta)\bigr),
\]
and
\[
a^{\mathrm{WL}}_{\eta,\alpha;\zeta,R,I}:
\mathfrak E^{\mathrm{WL}}_{\eta,\alpha;\zeta,R,I}\longrightarrow E^2,
\qquad
\xi\longmapsto
\bigl(\lambda_{\eta,R}(W_\eta),
\lambda_{\zeta,R}(B_\zeta)\bigr).
\]
For the wrapped/wrapped row of
Definition~\ref{def:wrapped-wrapped-correspondence-datum}, it is
\[
a^{\mathrm{WW}}_{\eta_1,\eta_2;\zeta,R}:
\mathfrak E^{\mathrm{WW}}_{\eta_1,\eta_2;\zeta,R}\longrightarrow E^3,
\qquad
\xi\longmapsto
\bigl(\lambda_{\eta_1,R}(W_1),
\lambda_{\eta_2,R}(W_2),
\lambda_{\zeta,R}(B_\zeta)\bigr).
\]
For the ordered two-sided mixed row of
Definition~\ref{def:two-sided-mixed-correspondence-datum}, it is
\[
a^{\mathrm{LWL}}_{\alpha_-;\eta;\beta_+,R,I_-,I_+}:
\mathfrak E^{\mathrm{LWL}}_{\alpha_-;\eta;\beta_+,R,I_-,I_+}
\longrightarrow E^2,
\qquad
\xi\longmapsto
\bigl(\lambda_{\eta,R}(W_\eta),
\lambda_{\zeta,R}(C_\zeta)\bigr).
\]
These maps record ordered source anchors and target anchors on the
unreduced stack.  The local/local correspondence has no wrapped input
and no wrapped anchor-memory row.  The determinant-anchor construction
and translation-weight computation are dedicated later packets; the
\(\chi=0\) degeneracy and the degree-zero Jordan--Hoelder extra anchor
are isolated by their own packets.  The anchor residual definition
\(o^\lambda_R\), its vanishing, anchor losslessness, quotient descent,
stabilizer linearization, and transition compatibility remain separate
rows.
Replacing the tuple above by a scalar Fock factor, a Borcherds
\(s\)-degree, or a post-quotient class is not this datum.
\end{definition}

The packet
\[
\texttt{certificates/hybrid/source\_target\_anchor\_memory\_definition}
\]
verifies
\(\texttt{SOURCE\_TARGET\_ANCHOR\_MEMORY\_DEFINITION\_VERIFIED}\):
the \(LW\), \(WL\), \(WW\), and \(LWL\) source/target anchor-memory
maps are defined before quotient.  The determinant anchor is constructed
by the next packet, and its translation weight is computed by the
following packet.  The \(\chi=0\) degeneracy is isolated by the
chi-zero packet.  The degree-zero Jordan--Hoelder extra-anchor datum is
recorded by
\(\texttt{certificates/hybrid/determinant\_anchor\_extra\_anchor\_data}\);
anchor-residual definition and vanishing, full losslessness,
quotient-descent, stabilizer linearization, populated finite-stack, and
transition rows remain separate obligations.

The packet
\[
\texttt{certificates/hybrid/determinant\_anchor\_construction}
\]
verifies
\(\texttt{DETERMINANT\_ANCHOR\_CONSTRUCTION\_VERIFIED}\):
Definition~\ref{def:determinant-anchor} constructs the normalized
determinant-of-cohomology anchor
\[
\lambda^{\det}_{\eta,R}(\mathcal F_{\eta,T})
=\det R\pi_{E,T,*}\mathcal F_{\eta,T}
\otimes\operatorname{pr}_E^*\mathcal O_E(-\chi_\eta\cdot0_E)
\in\operatorname{Pic}^0(E)
\]
as a functorial map to \(\operatorname{Pic}^0(E)\simeq E\).  The packet
checks the degree-zero normalization and determinant-functor
base-change; unit-weight normalization, the Jordan--Hoelder
extra-anchor datum, the residual definition \(o^\lambda_R\), residual
vanishing, losslessness, quotient descent, populated wrapped-prequotient
rows, and transition rows remain separate.  The \(\chi=0\) degeneracy
is isolated by the chi-zero packet below.
The translation-weight packet
\[
\texttt{certificates/hybrid/determinant\_anchor\_translation\_weight}
\]
verifies
\(\texttt{DETERMINANT\_ANCHOR\_TRANSLATION\_WEIGHT\_VERIFIED}\):
for the support-translation convention
\(a\cdot\mathcal F=(\operatorname{id}_S\times\tau_a)_*\mathcal F\),
Proposition~\ref{prop:determinant-anchor-translation-weight} gives
\[
\lambda^{\det}_{\eta,R}(a\cdot\mathcal F)
=\lambda^{\det}_{\eta,R}(\mathcal F)+\chi_\eta a
\quad\text{in }\operatorname{Pic}^0(E)\simeq E .
\]
The packet records the finite-stabilizer restriction
\(h\mapsto\chi_\eta h\).  It does not choose coherent stabilizer
linearizations, make the weight one, handle the \(\chi_\eta=0\)
degeneracy, identify the Jordan--Hoelder extra anchor, prove
losslessness, descend through the quotient, or prove transition
compatibility.
The chi-zero degeneracy packet
\[
\texttt{certificates/hybrid/determinant\_anchor\_chi\_zero\_degeneracy}
\]
verifies
\(\texttt{DETERMINANT\_ANCHOR\_CHI\_ZERO\_DEGENERACY\_VERIFIED}\):
Proposition~\ref{prop:determinant-anchor-chi-zero-degeneracy} proves
that when \(\chi_\eta=0\), the determinant anchor is constant on the
full \(E\)-translation orbit.  Freeness of the retained
\(E\)-translation action gives a rank-one separation defect for
determinant-anchor-only data.  The packet records the necessity of
extra anchor data; it does not itself construct that data, define
\(o^\lambda_R\), prove losslessness, descend through the quotient, or
prove transition compatibility.
The determinant-anchor extra-anchor-data packet
\[
\texttt{certificates/hybrid/determinant\_anchor\_extra\_anchor\_data}
\]
verifies
\(\texttt{DETERMINANT\_ANCHOR\_EXTRA\_ANCHOR\_DATA\_VERIFIED}\):
Proposition~\ref{prop:determinant-anchor-extra-data-separates-rank-two}
adds the degree-zero Jordan--Hoelder divisor of the elliptic
pushforward as anchor data.  The Abel sum is the determinant anchor,
and the rank-two pair \(\mathcal O_E^{\oplus2}\) and
\(L\oplus L^{-1}\), \(L\ne\mathcal O_E\), is separated despite equal
determinant.  This packet does not prove global anchor losslessness,
prove \(o^\lambda_R=0\), descend the repaired anchor through the
quotient, choose stabilizer linearizations, populate all finite
wrapped-prequotient rows, or prove transition compatibility.

\begin{definition}[Finite anchor residual]
\label{def:finite-anchor-residual}
Fix a finite HN height \(R\).  For
\(\eta\in\Gamma_R^{\mathrm{wr}}\), let
\[
\mathcal A_{\eta,R}:=
\operatorname{Sym}^{r_\eta}\operatorname{Pic}^0(E),
\qquad
\mathrm{Ab}_{\eta,R}:\mathcal A_{\eta,R}\longrightarrow
\operatorname{Pic}^0(E)\simeq E
\]
be the Abel-sum morphism.  A repaired wrapped anchor over the
prequotient row is a functorial map
\[
\Lambda_{\eta,R}:\mathcal M_{\eta,R}^{\mathrm{wr,rig}}
\longrightarrow \mathcal A_{\eta,R}
\]
such that
\[
\mathrm{Ab}_{\eta,R}\circ\Lambda_{\eta,R}
=\lambda^{\det}_{\eta,R},
\]
and whose restriction to the degree-zero vector-bundle row is the
Jordan--Hoelder divisor of
Proposition~\ref{prop:determinant-anchor-extra-data-separates-rank-two}.
For an exact sequence of semistable degree-zero elliptic pushforwards,
the anchor target has the divisor-union operation
\[
\mu_{\eta_1,\eta_2}:
\mathcal A_{\eta_1,R}\times\mathcal A_{\eta_2,R}
\longrightarrow \mathcal A_{\eta_1+\eta_2,R}.
\]
The finite anchor residual is the five-component defect
\[
\mathfrak o^\lambda_R=\bigl(
\mathfrak o^{\lambda,\mathrm{ex}}_R,\
\mathfrak o^{\lambda,\mathrm{unit}}_R,\
\mathfrak o^{\lambda,\mathrm{loss}}_R,\
\mathfrak o^{\lambda,\mathrm{multi}}_R,\
\mathfrak o^{\lambda,\mathrm{tr}}_{R'R}\bigr),
\]
with the following zero conditions.
The \(E\)-valued source/target anchors of
Definition~\ref{def:source-target-anchor-memory-datum} are the
Abel-sum shadows \(\mathrm{Ab}_{\eta,R}\circ\Lambda_{\eta,R}\); they
do not replace the repaired symmetric-product anchor before the
residual vanishes.
\[
\mathfrak o^{\lambda,\mathrm{ex}}_R=0
\]
if and only if every \(\eta\in\Gamma_R^{\mathrm{wr}}\) carries such a
functorial repaired anchor \(\Lambda_{\eta,R}\), including the
semistable vector-bundle locus needed for the Jordan--Hoelder divisor.
The component
\(\mathfrak o^{\lambda,\mathrm{unit}}_R\) vanishes if and only if the
chosen repaired anchors admit a finite-cover normalization with
translation weight one for the retained \(E\)-action.  The component
\(\mathfrak o^{\lambda,\mathrm{loss}}_R\) vanishes if and only if the
kernel pair of \(\Lambda_{\eta,R}\) on each retained wrapped row is the
retained equivalence relation, so that no extra wrapped position is
identified by the anchor.  The component
\(\mathfrak o^{\lambda,\mathrm{multi}}_R\) vanishes if and only if, on
every \(LW\), \(WL\), \(WW\), and \(LWL\) unreduced correspondence, the
target repaired anchor equals the divisor-union operation applied to
the ordered wrapped source anchors.  The component
\(\mathfrak o^{\lambda,\mathrm{tr}}_{R'R}\) vanishes if and only if,
for every \(R\le R'\), the repaired anchors commute with the HN
restriction \(R'\to R\).  This definition supplies the residual and
its zero semantics.  It does not prove any component vanishes, does
not descend the anchor through the reduced \(E\)-quotient, and does not
choose stabilizer linearizations.
\end{definition}

The packet
\[
\texttt{certificates/hybrid/anchor\_residual\_definition}
\]
verifies
\(\texttt{ANCHOR\_RESIDUAL\_DEFINITION\_VERIFIED}\):
Definition~\ref{def:finite-anchor-residual} defines the repaired
anchor target, its Abel-sum projection, and the five residual
components
\(\mathrm{ex}\), \(\mathrm{unit}\), \(\mathrm{loss}\),
\(\mathrm{multi}\), and \(\mathrm{tr}\).  The packet does not prove
first-window vanishing, quotient descent, stabilizer linearization,
finite-stage population, or transition compatibility.

\begin{definition}[Quotient-after-correspondence pseudofunctor]
\label{def:quotient-after-correspondence-pseudofunctor}
\textnormal{[definition only]}
Fix a finite HN height \(R\).  Let
\(\mathsf{Corr}^{E,\mathrm{hyb}}_R\) be the unreduced
\(E\)-equivariant hybrid correspondence category generated by the
retained local, mixed, wrapped, two-sided mixed, and flag
correspondence stacks before the reduced \(E\)-quotient.  A
quotient-after-correspondence datum consists of the following finite
rows.

For every retained object \(Y\) of
\(\mathsf{Corr}^{E,\mathrm{hyb}}_R\), including local closed
configuration charts, wrapped rigidified prequotients, correspondence
stacks, and flag stacks, choose the quotient stack
\[
\overline Y=[Y/E]
\]
for the retained diagonal translation action and the quotient map
\[
q_Y:Y\longrightarrow\overline Y .
\]
For every generating unreduced correspondence
\[
e=(Y_0\xleftarrow{p_e}\mathfrak E_e\xrightarrow{q_e}Y_1)
\]
whose maps are \(E\)-equivariant, choose the reduced correspondence
\[
\overline e=
(\overline Y_0\xleftarrow{\overline p_e}
[\mathfrak E_e/E]\xrightarrow{\overline q_e}\overline Y_1)
\]
induced after the same correspondence stack has been formed.  The
assignment
\[
\mathcal Q^{\mathrm{corr}}_{E,R}(Y)=\overline Y,\qquad
\mathcal Q^{\mathrm{corr}}_{E,R}(e)=\overline e
\]
is the quotient-after-correspondence pseudofunctor at the level of
objects and generating arrows.

For a reduced vanishing-cycle coefficient \(K_Y\) on \(Y\), the datum
also specifies the coefficient-descent interface
\[
\overline K_Y\in D^b_c(\overline Y),\qquad
q_Y^*\overline K_Y\simeq K_Y,
\]
whenever the relevant stabilizer linearization and orientation descent
rows are supplied.  The induced chain functor
\[
Q_{E,R}:\mathsf{BM}^{E,\mathrm{hyb}}_R
\longrightarrow
\mathsf{BM}^{\mathrm{red},\mathrm{hyb}}_R
\]
is not part of this definition; it is the later Borel--Moore chain
row built from these quotient and coefficient-descent interfaces.

The finite quotient residual attached to this datum is
\[
\mathfrak o^Q_R=
\bigl(
\mathfrak o^{Q,\mathrm{form}}_R,\
\mathfrak o^{Q,\mathrm{comp}}_R,\
\mathfrak o^{Q,\mathrm{bc}}_R,\
\mathfrak o^{Q,\mathrm{TS}}_R,\
\mathfrak o^{Q,\mathrm{flag}}_R,\
\mathfrak o^{Q,\mathrm{BM}}_R,\
\mathfrak o^{Q,\theta}_R,\
\mathfrak o^{Q,\mathrm{tr}}_{R'R}
\bigr).
\]
The first component says that objects, arrows, quotient maps, and
coefficient-descent interfaces are defined.  The remaining components
measure, respectively, preservation of composition, base-change,
Thom--Sebastiani transport, flag-stack coherences, Borel--Moore chain
functoriality, the comparison maps \(\theta^Q_{\mu,R}\), and
compatibility with \(R'\to R\) transitions.

This is a quotient-after-correspondence definition: the correspondence
stack \(\mathfrak E_e\) is formed before applying \([-/E]\).  It is
not a quotient-first Hall product on already reduced object stacks.
\end{definition}

The packet
\[
\texttt{certificates/hybrid/quotient\_after\_correspondence\_pseudofunctor\_definition}
\]
verifies
\(\texttt{QUOTIENT\_AFTER\_CORRESPONDENCE\_PSEUDOFUNCTOR\_DEFINED}\):
Definition~\ref{def:quotient-after-correspondence-pseudofunctor}
defines the object assignment, arrow assignment, coefficient-descent
interface, and eight quotient residual components of the
quotient-after-correspondence pseudofunctor before Borel--Moore
chains.  The composition component is supplied by
Proposition~\ref{prop:quotient-after-correspondence-preserves-composition}.
The base-change square component is supplied by
Proposition~\ref{prop:quotient-after-correspondence-preserves-base-change}.
The Thom--Sebastiani descent component is supplied by
Proposition~\ref{prop:quotient-after-correspondence-preserves-thom-sebastiani}.
The quotient-first exclusion theorem is supplied by
Proposition~\ref{prop:quotient-after-correspondence-excludes-quotient-first}.
The Borel--Moore chain functor is supplied by
Proposition~\ref{prop:quotient-after-correspondence-borel-moore-chain-functor}.
The operation-comparison maps \(\theta^Q_{\mu,R}\) are supplied by
Proposition~\ref{prop:quotient-after-correspondence-theta-mu-comparison}.
The quotient two-step flag associativity component is supplied by
Proposition~\ref{prop:quotient-after-correspondence-theta-mu-associativity}.
It does not prove compact-support pushforward base change, transition
compatibility, or aggregate population.

\begin{proposition}[Quotient-after-correspondence preserves composition]
\label{prop:quotient-after-correspondence-preserves-composition}
\textnormal{[proved for finite-height composition of spans]}
Fix \(R\) and the quotient-after-correspondence datum of
Definition~\ref{def:quotient-after-correspondence-pseudofunctor}.
Let
\[
e=(Y_0\xleftarrow{p_e}\mathfrak E_e\xrightarrow{q_e}Y_1),
\qquad
f=(Y_1\xleftarrow{p_f}\mathfrak E_f\xrightarrow{q_f}Y_2)
\]
be composable retained \(E\)-equivariant hybrid correspondences.  The
unreduced composite has middle stack
\[
\mathfrak E_e\times_{Y_1}\mathfrak E_f
\]
with the diagonal \(E\)-action.  There is a canonical isomorphism of
reduced correspondences
\[
\Xi_{f,e}\colon
\mathcal Q^{\mathrm{corr}}_{E,R}(f\circ e)
\xrightarrow{\ \sim\ }
\mathcal Q^{\mathrm{corr}}_{E,R}(f)
\circ
\mathcal Q^{\mathrm{corr}}_{E,R}(e),
\]
whose middle-stack component is
\[
\bigl[(\mathfrak E_e\times_{Y_1}\mathfrak E_f)/E\bigr]
\simeq
[\mathfrak E_e/E]\times_{[Y_1/E]}[\mathfrak E_f/E].
\]
The isomorphisms \(\Xi_{f,e}\) satisfy the associativity cocycle for
three composable correspondences.  Hence
\[
\mathfrak o^{Q,\mathrm{comp}}_R=0.
\]
This proves no base-change theorem for compact-support pushforwards,
no Thom--Sebastiani compatibility, no Borel--Moore chain functor
\(Q_{E,R}\), no comparison map \(\theta^Q_{\mu,R}\), no quotient-first
exclusion theorem, and no transition or aggregate statement.
\end{proposition}

\begin{proof}
The maps \(q_e\) and \(p_f\) are \(E\)-equivariant, so the fibre
product \(\mathfrak E_e\times_{Y_1}\mathfrak E_f\) carries the
diagonal \(E\)-action.  For a test scheme \(T\), an object of
\([\mathfrak E_e/E]\times_{[Y_1/E]}[\mathfrak E_f/E]\) over \(T\)
consists of \(E\)-torsors \(P_e,P_f\to T\), equivariant maps
\[
u:P_e\to\mathfrak E_e,\qquad v:P_f\to\mathfrak E_f,
\]
and an isomorphism \(\alpha:P_e\xrightarrow{\sim}P_f\) over the
common image in \([Y_1/E]\); equivalently \(p_fv\alpha=q_eu\) as maps
\(P_e\to Y_1\).  Transporting \(v\) along \(\alpha\) gives one torsor
\(P_e\to T\) with compatible equivariant maps to \(\mathfrak E_e\) and
\(\mathfrak E_f\).  This is precisely an \(E\)-equivariant map
\[
P_e\longrightarrow \mathfrak E_e\times_{Y_1}\mathfrak E_f,
\]
or equivalently an object of
\([(\mathfrak E_e\times_{Y_1}\mathfrak E_f)/E]\) over \(T\).
The converse sends one torsor with a map to the fibre product to the
two quotient-stack objects with the identity torsor isomorphism.
Morphisms are transported by the same torsor isomorphisms.  The two
constructions are inverse and functorial in \(T\), giving the
displayed isomorphism of quotient stacks.

The source and target maps to \([Y_0/E]\) and \([Y_2/E]\) are induced
from \(p_e\) and \(q_f\), so the middle-stack isomorphism identifies
the quotient of the unreduced composite with the composite of the two
reduced correspondences.  For three composable correspondences both
iterated comparisons are obtained from the same quotient of the
triple fibre product, and the two maps agree by the associativity of
2-fibre products and the universal property of quotient stacks.  Thus
the composition defect has rank zero, while all coefficient,
base-change, Thom--Sebastiani, Borel--Moore, transition, and aggregate
claims remain outside this proposition.
\end{proof}

The packet
\[
\texttt{certificates/hybrid/quotient\_after\_correspondence\_composition}
\]
verifies
\(\texttt{QUOTIENT\_AFTER\_CORRESPONDENCE\_COMPOSITION\_VERIFIED}\):
Proposition~\ref{prop:quotient-after-correspondence-preserves-composition}
proves that quotient-after-correspondence commutes with finite-height
composition of retained \(E\)-equivariant spans and that the
associativity cocycle has defect rank zero.  The base-change
component is supplied by
Proposition~\ref{prop:quotient-after-correspondence-preserves-base-change}.
The Thom--Sebastiani descent component is supplied by
Proposition~\ref{prop:quotient-after-correspondence-preserves-thom-sebastiani}.
The quotient-first exclusion theorem is supplied by
Proposition~\ref{prop:quotient-after-correspondence-excludes-quotient-first}.
The Borel--Moore chain functor is supplied by
Proposition~\ref{prop:quotient-after-correspondence-borel-moore-chain-functor}.
The operation-comparison maps \(\theta^Q_{\mu,R}\) are supplied by
Proposition~\ref{prop:quotient-after-correspondence-theta-mu-comparison}.
The quotient two-step flag associativity component is supplied by
Proposition~\ref{prop:quotient-after-correspondence-theta-mu-associativity}.
The HN-transition compatibility of \(Q_{E,R}\) is supplied by
Proposition~\ref{prop:quotient-after-correspondence-hn-transition-compatibility}.
It does not prove aggregate population.

\begin{proposition}[Quotient-after-correspondence preserves base-change squares]
\label{prop:quotient-after-correspondence-preserves-base-change}
\textnormal{[proved for finite-height cartesian squares of stacks]}
Fix \(R\) and the quotient-after-correspondence datum of
Definition~\ref{def:quotient-after-correspondence-pseudofunctor}.
Let
\[
\begin{array}{ccc}
X' & \xrightarrow{\widetilde g} & X\\
\scriptstyle h'\downarrow && \downarrow\scriptstyle h\\
Y' & \xrightarrow{g} & Y
\end{array}
\qquad(\square)
\]
be a retained \(E\)-equivariant cartesian square of stacks appearing
as a source, target, open, or compactified square of the finite
hybrid correspondence calculus.  Then the quotient square
\[
\begin{array}{ccc}
[X'/E] & \xrightarrow{[\widetilde g/E]} & [X/E]\\
\scriptstyle [h'/E]\downarrow && \downarrow\scriptstyle [h/E]\\
{[Y'/E]} & \xrightarrow{[g/E]} & {[Y/E]}
\end{array}
\]
is cartesian.  Equivalently, the canonical morphism
\[
[X'/E]\longrightarrow [X/E]\times_{[Y/E]}[Y'/E]
\]
is an isomorphism.  Hence quotient-after-correspondence preserves the
retained base-change squares and
\[
\mathfrak o^{Q,\mathrm{bc}}_R=0.
\]
This proves only the stack-level cartesian-square component of the
quotient residual.  It does not prove compact-support pushforward
base change, coefficient pullback for vanishing cycles, the
projection formula, Thom--Sebastiani compatibility, the chain functor
\(Q_{E,R}\), the comparison maps \(\theta^Q_{\mu,R}\), transition
compatibility, or aggregate population.
\end{proposition}

\begin{proof}
Since the square \((\square)\) is \(E\)-equivariant and cartesian,
there is an \(E\)-equivariant isomorphism
\[
X'\simeq X\times_Y Y'.
\]
For a test scheme \(T\), an object of
\([X/E]\times_{[Y/E]}[Y'/E]\) over \(T\) consists of \(E\)-torsors
\(P_X,P_{Y'}\to T\), equivariant maps
\[
u:P_X\to X,\qquad v:P_{Y'}\to Y',
\]
and an isomorphism
\(\alpha:P_X\xrightarrow{\sim}P_{Y'}\) over the common image in
\([Y/E]\).  After transporting \(v\) along \(\alpha\), this is one
\(E\)-torsor \(P_X\to T\) with maps \(u:P_X\to X\) and
\(v\alpha:P_X\to Y'\) satisfying
\[
h\,u=g\,v\,\alpha .
\]
By the cartesian identity \(X'\simeq X\times_Y Y'\), this is the same
as an \(E\)-equivariant map \(P_X\to X'\), hence an object of
\([X'/E]\) over \(T\).  The converse sends an equivariant map
\(P\to X'\) to its two projections with the identity torsor
isomorphism.  Morphisms are carried by the same torsor isomorphisms,
and the two constructions are inverse and functorial in \(T\).
Therefore the quotient square is cartesian.

Applying this argument to each retained source, target, open, and
compactified square in the finite hybrid correspondence calculus gives
zero defect for preservation of base-change squares.  The proof has
not applied \(h_!\), \(h^*\), reduced Thom--Sebastiani transport, or
Borel--Moore chains, so the corresponding functorial and chain-level
claims remain separate obligations.
\end{proof}

The packet
\[
\texttt{certificates/hybrid/quotient\_after\_correspondence\_base\_change}
\]
verifies
\(\texttt{QUOTIENT\_AFTER\_CORRESPONDENCE\_BASE\_CHANGE\_VERIFIED}\):
Proposition~\ref{prop:quotient-after-correspondence-preserves-base-change}
proves that quotient-after-correspondence preserves retained
cartesian base-change squares of finite-height \(E\)-equivariant
stacks.  It does not prove compact-support pushforward base change,
coefficient pullback, or aggregate population.  The
Thom--Sebastiani descent component is supplied by
Proposition~\ref{prop:quotient-after-correspondence-preserves-thom-sebastiani}.
The quotient-first exclusion theorem is supplied by
Proposition~\ref{prop:quotient-after-correspondence-excludes-quotient-first}.
The Borel--Moore chain functor is supplied by
Proposition~\ref{prop:quotient-after-correspondence-borel-moore-chain-functor}.
The operation-comparison maps \(\theta^Q_{\mu,R}\) are supplied by
Proposition~\ref{prop:quotient-after-correspondence-theta-mu-comparison}.
The quotient two-step flag associativity component is supplied by
Proposition~\ref{prop:quotient-after-correspondence-theta-mu-associativity}.
The HN-transition compatibility of \(Q_{E,R}\) is supplied by
Proposition~\ref{prop:quotient-after-correspondence-hn-transition-compatibility}.

\begin{proposition}[Quotient-after-correspondence preserves Thom--Sebastiani transport]
\label{prop:quotient-after-correspondence-preserves-thom-sebastiani}
\textnormal{[proved for equivariant finite-height Thom--Sebastiani rows]}
Fix \(R\) and the quotient-after-correspondence datum of
Definition~\ref{def:quotient-after-correspondence-pseudofunctor}.
Let
\[
e=(Y_{\mathrm{src}}\xleftarrow{p_e}\mathfrak E_e
\xrightarrow{q_e}Y_{\mathrm{tar}})
\]
be a retained \(E\)-equivariant correspondence whose unreduced
reduced vanishing-cycle coefficients carry an \(E\)-equivariant
Thom--Sebastiani transport
\[
\operatorname{TS}^{\mathrm{red}}_e:
p_e^*K_{Y_{\mathrm{src}}}\xrightarrow{\sim}K_{\mathfrak E_e}.
\]
Assume the coefficient-descent interfaces
\[
q_{Y_{\mathrm{src}}}^{*}\overline K_{Y_{\mathrm{src}}}
\simeq K_{Y_{\mathrm{src}}},\qquad
q_{\mathfrak E_e}^{*}\overline K_{\mathfrak E_e}
\simeq K_{\mathfrak E_e}
\]
are supplied, including the stabilizer linearizations and orientation
descent needed for the reduced vanishing-cycle complexes.  Then there
is a unique reduced Thom--Sebastiani transport
\[
\overline{\operatorname{TS}}^{\mathrm{red}}_{\overline e}:
\overline p_e^{\,*}\overline K_{Y_{\mathrm{src}}}
\xrightarrow{\sim}\overline K_{\mathfrak E_e}
\]
on the quotient correspondence
\[
\overline e=([Y_{\mathrm{src}}/E]\xleftarrow{\overline p_e}
[\mathfrak E_e/E]\xrightarrow{\overline q_e}[Y_{\mathrm{tar}}/E])
\]
whose pullback along \(q_{\mathfrak E_e}\) is
\(\operatorname{TS}^{\mathrm{red}}_e\) under the descent
identifications and the cartesian quotient square of
Proposition~\ref{prop:quotient-after-correspondence-preserves-base-change}.
Hence
\[
\mathfrak o^{Q,\mathrm{TS}}_R=0.
\]
This proves only descent of the reduced Thom--Sebastiani coefficient
transport through the quotient-after-correspondence construction.  It
does not prove the existence of the unreduced BBDJS
Thom--Sebastiani row, the Joyce--Upmeier orientation row,
compact-support pushforward base change, flag-stack coherences, the
Borel--Moore chain functor \(Q_{E,R}\), the comparison maps
\(\theta^Q_{\mu,R}\), transition compatibility, or aggregate
population.
\end{proposition}

\begin{proof}
The quotient square for the source map is cartesian by
Proposition~\ref{prop:quotient-after-correspondence-preserves-base-change}:
\[
\begin{array}{ccc}
\mathfrak E_e & \xrightarrow{q_{\mathfrak E_e}} & [\mathfrak E_e/E]\\
\scriptstyle p_e\downarrow && \downarrow\scriptstyle \overline p_e\\
Y_{\mathrm{src}} & \xrightarrow{q_{Y_{\mathrm{src}}}} &
[Y_{\mathrm{src}}/E].
\end{array}
\]
Therefore
\[
q_{\mathfrak E_e}^*\overline p_e^{\,*}
\overline K_{Y_{\mathrm{src}}}
\simeq
p_e^*q_{Y_{\mathrm{src}}}^*
\overline K_{Y_{\mathrm{src}}}
\simeq p_e^*K_{Y_{\mathrm{src}}}.
\]
The second descent interface gives
\[
q_{\mathfrak E_e}^*\overline K_{\mathfrak E_e}
\simeq K_{\mathfrak E_e}.
\]
Thus the unreduced Thom--Sebastiani isomorphism is a morphism between
the pullbacks of the two reduced coefficient objects on
\([\mathfrak E_e/E]\).

Constructible sheaves on the quotient stack \([\mathfrak E_e/E]\) are
equivalent to \(E\)-equivariant constructible sheaves on
\(\mathfrak E_e\), with the stated stabilizer linearizations.  Under
this equivalence, morphisms on the quotient stack are exactly
equivariant morphisms on \(\mathfrak E_e\).  Since
\(\operatorname{TS}^{\mathrm{red}}_e\) is \(E\)-equivariant by
hypothesis, it descends uniquely to a morphism
\[
\overline{\operatorname{TS}}^{\mathrm{red}}_{\overline e}:
\overline p_e^{\,*}\overline K_{Y_{\mathrm{src}}}
\longrightarrow\overline K_{\mathfrak E_e}.
\]
Its pullback is an isomorphism, hence the descended morphism is an
isomorphism.  The same equivalence of quotient sheaves proves
uniqueness and functoriality for every retained finite-height
Thom--Sebastiani row.  No pushforward, Borel--Moore chain functor, or
flag comparison has been applied, so those coherences remain outside
this proposition.
\end{proof}

The packet
\[
\texttt{certificates/hybrid/quotient\_after\_correspondence\_thom\_sebastiani}
\]
verifies
\(\texttt{QUOTIENT\_AFTER\_CORRESPONDENCE\_THOM\_SEBASTIANI\_VERIFIED}\):
Proposition~\ref{prop:quotient-after-correspondence-preserves-thom-sebastiani}
proves that an equivariant reduced Thom--Sebastiani transport descends
uniquely through the quotient-after-correspondence construction once
the coefficient-descent interfaces are supplied.  It does not prove
aggregate population.
The operation-comparison maps \(\theta^Q_{\mu,R}\) are supplied by
Proposition~\ref{prop:quotient-after-correspondence-theta-mu-comparison}.
The quotient-first exclusion theorem is
Proposition~\ref{prop:quotient-after-correspondence-excludes-quotient-first}.
The Borel--Moore chain functor is
Proposition~\ref{prop:quotient-after-correspondence-borel-moore-chain-functor}.
The quotient two-step flag associativity component is
Proposition~\ref{prop:quotient-after-correspondence-theta-mu-associativity}.
The HN-transition compatibility of \(Q_{E,R}\) is
Proposition~\ref{prop:quotient-after-correspondence-hn-transition-compatibility}.

\begin{proposition}[Quotient-first construction is excluded]
\label{prop:quotient-after-correspondence-excludes-quotient-first}
\textnormal{[proved as a finite-height image characterization]}
Fix \(R\) and the quotient-after-correspondence datum of
Definition~\ref{def:quotient-after-correspondence-pseudofunctor}.
Every generating arrow in the image of
\(\mathcal Q^{\mathrm{corr}}_{E,R}\) is represented by a retained
unreduced \(E\)-equivariant correspondence
\[
e=(Y_0\xleftarrow{p_e}\mathfrak E_e\xrightarrow{q_e}Y_1)
\]
together with the quotient presentation
\[
\overline e=
([Y_0/E]\xleftarrow{\overline p_e}
[\mathfrak E_e/E]\xrightarrow{\overline q_e}[Y_1/E]).
\]
Equivalently, the two squares
\[
\begin{array}{ccc}
\mathfrak E_e & \longrightarrow & [\mathfrak E_e/E]\\
\scriptstyle p_e\downarrow && \downarrow\scriptstyle\overline p_e\\
Y_0 & \longrightarrow & [Y_0/E]
\end{array}
\qquad
\begin{array}{ccc}
\mathfrak E_e & \longrightarrow & [\mathfrak E_e/E]\\
\scriptstyle q_e\downarrow && \downarrow\scriptstyle\overline q_e\\
Y_1 & \longrightarrow & [Y_1/E]
\end{array}
\]
are part of the arrow data.  A quotient-first span, obtained by first
passing to \([Y_0/E]\) and \([Y_1/E]\) and then choosing an independent
middle stack over their product, is not an arrow of
\(\mathcal Q^{\mathrm{corr}}_{E,R}\) unless it is equipped with a
specified isomorphism to \([\mathfrak E_e/E]\) compatible with these
two squares.  In that exceptional case it is the same quotient-after-
correspondence arrow after transport along the specified isomorphism.
Thus no quotient-first Hall product is used in the finite quotient
correspondence calculus, and the quotient-first exclusion row has
defect rank zero.
\end{proposition}

\begin{proof}
The object assignment of
Definition~\ref{def:quotient-after-correspondence-pseudofunctor}
sends an unreduced retained object \(Y\) to \([Y/E]\).  Its arrow
assignment does not choose a new reduced extension stack.  It sends
the already formed unreduced correspondence
\[
Y_0\xleftarrow{p_e}\mathfrak E_e\xrightarrow{q_e}Y_1
\]
to the reduced correspondence with middle stack
\([\mathfrak E_e/E]\).  Hence the image of
\(\mathcal Q^{\mathrm{corr}}_{E,R}\) is the full finite-height family
of reduced spans carrying such an unreduced quotient presentation.

Let
\[
[Y_0/E]\xleftarrow{a}Z\xrightarrow{b}[Y_1/E]
\]
be a quotient-first candidate obtained after quotienting the two
object stacks.  If \(Z\) is not supplied with an isomorphism
\[
Z\simeq[\mathfrak E_e/E]
\]
for a retained unreduced \(E\)-equivariant correspondence \(e\), and
if \(a,b\) are not the transported maps \(\overline p_e,\overline q_e\),
then this candidate lacks the quotient map
\(\mathfrak E_e\to[\mathfrak E_e/E]\) used in the definition.  It is
therefore not in the image of
\(\mathcal Q^{\mathrm{corr}}_{E,R}\).  If such an isomorphism is
supplied, transporting the span structure along it identifies the
candidate with the quotient-after-correspondence span assigned to
\(e\); no independent quotient-first construction remains.

The composition comparison of
Proposition~\ref{prop:quotient-after-correspondence-preserves-composition}
uses the quotient of the unreduced fibre product
\(\mathfrak E_e\times_{Y_1}\mathfrak E_f\).  The base-change
comparison of
Proposition~\ref{prop:quotient-after-correspondence-preserves-base-change}
uses quotient presentations of the unreduced cartesian squares.  The
Thom--Sebastiani descent of
Proposition~\ref{prop:quotient-after-correspondence-preserves-thom-sebastiani}
pulls coefficients back along
\(\mathfrak E_e\to[\mathfrak E_e/E]\).  None of these rows can be
interpreted on an independent quotient-first middle stack unless it
has first been identified with the quotient of the unreduced middle
stack.  The exclusion is therefore a statement about the finite
correspondence calculus itself, not a Borel--Moore, theta-comparison,
transition, or aggregate-population claim.
\end{proof}

The packet
\[
\texttt{certificates/hybrid/quotient\_after\_correspondence\_quotient\_first\_exclusion}
\]
verifies
\(\texttt{QUOTIENT\_AFTER\_CORRESPONDENCE\_QUOTIENT\_FIRST\_EXCLUSION\_VERIFIED}\):
Proposition~\ref{prop:quotient-after-correspondence-excludes-quotient-first}
proves that every reduced generating arrow used by
\(\mathcal Q^{\mathrm{corr}}_{E,R}\) carries a quotient presentation
from a preformed unreduced \(E\)-equivariant correspondence.  An
object-quotient-first Hall product is excluded unless it is identified
with that quotient presentation.  This packet does not prove transition
compatibility for Hall products or aggregate population.  The
Borel--Moore chain functor is supplied by
Proposition~\ref{prop:quotient-after-correspondence-borel-moore-chain-functor}.
The operation-comparison maps \(\theta^Q_{\mu,R}\) are supplied by
Proposition~\ref{prop:quotient-after-correspondence-theta-mu-comparison}.
The quotient two-step flag associativity component is supplied by
Proposition~\ref{prop:quotient-after-correspondence-theta-mu-associativity}.
The HN-transition compatibility of \(Q_{E,R}\) is supplied by
Proposition~\ref{prop:quotient-after-correspondence-hn-transition-compatibility}.

\begin{proposition}[The quotient \(E\)-descent functor on Borel--Moore chains]
\label{prop:quotient-after-correspondence-borel-moore-chain-functor}
\textnormal{[proved for quotient-presented finite-height coefficient objects]}
Fix \(R\) and the quotient-after-correspondence datum of
Definition~\ref{def:quotient-after-correspondence-pseudofunctor}.
For every retained object \(Y\) of
\(\mathsf{Corr}^{E,\mathrm{hyb}}_R\), suppose the reduced
vanishing-cycle coefficient \(K_Y\in D^b_{c,E}(Y)\) is supplied with
the descent interface
\[
\overline K_Y\in D^b_c([Y/E]),\qquad
q_Y^*\overline K_Y\simeq K_Y .
\]
Let
\[
C_*^{\mathrm{BM}}(Z,L):=
R\Gamma\!\left(Z,\mathbb D_Z L\right)
\]
denote Borel--Moore chains with constructible coefficient \(L\), with
\(\mathbb D_Z\) the Verdier duality functor on the retained finite
stack \(Z\).  Define
\[
C_*^{\mathrm{BM},E}(Y,K_Y)
:=
C_*^{\mathrm{BM}}([Y/E],\overline K_Y).
\]
Then the assignments
\[
(Y,K_Y)\longmapsto
([Y/E],\overline K_Y),\qquad
C_*^{\mathrm{BM},E}(Y,K_Y)\longmapsto
C_*^{\mathrm{BM}}([Y/E],\overline K_Y)
\]
extend to a functor
\[
Q_{E,R}:\mathsf{BM}^{E,\mathrm{hyb}}_R
\longrightarrow
\mathsf{BM}^{\mathrm{red},\mathrm{hyb}}_R .
\]
On an \(E\)-equivariant coefficient morphism
\(\alpha:K_Y\to L_Y\), the functor sends \(\alpha\) to the descended
morphism
\[
\overline\alpha:\overline K_Y\longrightarrow\overline L_Y
\]
under the quotient-stack descent equivalence, and then applies
\(C_*^{\mathrm{BM}}([Y/E],-)\).  Hence the Borel--Moore chain
component of the quotient residual vanishes:
\[
\mathfrak o^{Q,\mathrm{BM}}_R=0.
\]
This proves only the chain-level descent functor.  It does not prove
flag coherences, compact-support pushforward base change, the
comparison isomorphisms
\(\theta^Q_{\mu,R}\), compatibility with products or coproducts,
transition compatibility, protected-integration compatibility, or
aggregate population.
\end{proposition}

\begin{proof}
The coefficient-descent interface identifies the supplied
\(E\)-equivariant coefficient \(K_Y\) with the pullback of a
constructible coefficient on the quotient stack.  The standard
descent equivalence
\[
D^b_c([Y/E])\simeq D^b_{c,E}(Y)
\]
is fully faithful on morphisms: an equivariant morphism
\(\alpha:K_Y\to L_Y\) is uniquely the pullback of a morphism
\(\overline\alpha:\overline K_Y\to\overline L_Y\) whenever both
coefficients carry the prescribed descent data.  Identities and
composites descend to identities and composites because this is an
equivalence of categories.

Borel--Moore chains are functorial for coefficient morphisms on the
fixed stack:
\[
\overline\alpha:\overline K_Y\to\overline L_Y
\quad\Longmapsto\quad
R\Gamma([Y/E],\mathbb D\overline\alpha):
C_*^{\mathrm{BM}}([Y/E],\overline K_Y)
\to
C_*^{\mathrm{BM}}([Y/E],\overline L_Y).
\]
The object assignment \(Y\mapsto [Y/E]\) and the coefficient
assignment \(K_Y\mapsto\overline K_Y\) therefore define a functor on
the finite additive chain category generated by the retained
quotient-presented coefficient objects.  This is \(Q_{E,R}\).

No map of the form \(q_!\operatorname{TS}p^*\) has been compared with
its reduced counterpart.  The present functor supplies the chain
objects and descended coefficient morphisms on which the later
\(\theta^Q_{\mu,R}\) comparison maps may act; it does not construct
those maps.
\end{proof}

The packet
\[
\texttt{certificates/hybrid/quotient\_after\_correspondence\_bm\_chain\_functor}
\]
verifies
\(\texttt{QUOTIENT\_AFTER\_CORRESPONDENCE\_BM\_CHAIN\_FUNCTOR\_VERIFIED}\):
Proposition~\ref{prop:quotient-after-correspondence-borel-moore-chain-functor}
constructs \(Q_{E,R}\) on quotient-presented Borel--Moore chain
objects and proves \(\mathfrak o^{Q,\mathrm{BM}}_R=0\).  This packet
does not prove compact-support pushforward base change, protected
integration, or aggregate population.  The operation-comparison maps
\(\theta^Q_{\mu,R}\) are
supplied by
Proposition~\ref{prop:quotient-after-correspondence-theta-mu-comparison}.
The quotient two-step flag associativity component is supplied by
Proposition~\ref{prop:quotient-after-correspondence-theta-mu-associativity}.
The HN-transition compatibility of \(Q_{E,R}\) is supplied by
Proposition~\ref{prop:quotient-after-correspondence-hn-transition-compatibility}.

\begin{proposition}[Quotient comparison maps for retained operations]
\label{prop:quotient-after-correspondence-theta-mu-comparison}
\textnormal{[proved for quotient-presented finite-height operation rows]}
Fix \(R\) and the quotient-after-correspondence datum of
Definition~\ref{def:quotient-after-correspondence-pseudofunctor}.
Let
\[
e=\bigl(Y_1\times\cdots\times Y_k
\xleftarrow{\ p_e\ }\mathfrak E_e\xrightarrow{\ q_e\ }Y_0\bigr)
\]
be a retained \(k\)-ary local, mixed, wrapped, or flag operation row.
Assume the input, middle, and target coefficient objects carry
quotient-descent interfaces, the external product coefficient descends
to
\[
\overline K_1\boxtimes\cdots\boxtimes\overline K_k
\quad\text{on}\quad
[Y_1/E]\times\cdots\times [Y_k/E],
\]
the reduced Thom--Sebastiani transport descends as in
Proposition~\ref{prop:quotient-after-correspondence-preserves-thom-sebastiani},
and the chosen compact-support pushforward row for \(q_e\) is
quotient-presented by an isomorphism
\[
\delta_{q,e}\colon
\overline{\,q_{e,!}^{\mathrm{cs}}L_e\,}
\xrightarrow{\ \sim\ }
\overline q_{e,!}^{\mathrm{cs}}\overline L_e ,
\]
where
\[
L_e=
\operatorname{TS}^{\mathrm{red}}_e
p_e^*(K_1\boxtimes\cdots\boxtimes K_k),
\qquad
\overline L_e=
\overline{\operatorname{TS}}^{\mathrm{red}}_{\overline e}
\overline p_e^{\,*}
(\overline K_1\boxtimes\cdots\boxtimes\overline K_k).
\]
Then there is a canonical chain isomorphism
\[
\theta^Q_{\mu,e,R}\colon
Q_{E,R}\,
q_{e,!}^{\mathrm{cs}}\operatorname{TS}^{\mathrm{red}}_e p_e^*
\xrightarrow{\ \sim\ }
\overline q_{e,!}^{\mathrm{cs}}\,
\overline{\operatorname{TS}}^{\mathrm{red}}_{\overline e}
\overline p_e^{\,*}(Q_{E,R}^{\boxtimes k}).
\]
For every retained operation row with the displayed descent data,
\[
\mathfrak o^{Q,\theta}_R=0 .
\]
This proves only the existence of the operation-comparison
isomorphisms \(\theta^Q_{\mu,R}\).  It does not prove their
compatibility with two-step flag associativity, coproducts,
primitive projection, \(R'\to R\) transitions, protected integration,
or aggregate hybrid-carrier population.
\end{proposition}

\begin{proof}
By Proposition~\ref{prop:quotient-after-correspondence-preserves-base-change},
the quotient square for \(p_e\) is cartesian.  Hence the pullback of
the descended external product along the quotient middle stack is the
descent of
\[
p_e^*(K_1\boxtimes\cdots\boxtimes K_k).
\]
Proposition~\ref{prop:quotient-after-correspondence-preserves-thom-sebastiani}
then identifies the descended middle coefficient with
\[
\overline L_e=
\overline{\operatorname{TS}}^{\mathrm{red}}_{\overline e}
\overline p_e^{\,*}
(\overline K_1\boxtimes\cdots\boxtimes\overline K_k).
\]
The quotient-presented compact-support row for the target map gives
the specified isomorphism
\[
\delta_{q,e}\colon
\overline{\,q_{e,!}^{\mathrm{cs}}L_e\,}
\xrightarrow{\sim}
\overline q_{e,!}^{\mathrm{cs}}\overline L_e
\]
on \([Y_0/E]\).

Applying the Borel--Moore chain functor
of Proposition~\ref{prop:quotient-after-correspondence-borel-moore-chain-functor}
to \(\delta_{q,e}\) gives the displayed chain isomorphism.  Its
source is, by definition,
\[
Q_{E,R}\,
q_{e,!}^{\mathrm{cs}}\operatorname{TS}^{\mathrm{red}}_e p_e^*
(K_1,\ldots,K_k),
\]
and its target is the reduced operation
\[
\overline q_{e,!}^{\mathrm{cs}}\,
\overline{\operatorname{TS}}^{\mathrm{red}}_{\overline e}
\overline p_e^{\,*}
(Q_{E,R}K_1,\ldots,Q_{E,R}K_k).
\]
Uniqueness follows from full faithfulness of constructible descent on
quotient stacks.  The proof compares one retained operation with its
reduced quotient operation; it does not compare two parenthesized
operations, a product with a coproduct, a primitive projection, or two
HN heights.
\end{proof}

The packet
\[
\texttt{certificates/hybrid/quotient\_after\_correspondence\_theta\_mu\_comparison}
\]
verifies
\(\texttt{QUOTIENT\_AFTER\_CORRESPONDENCE\_THETA\_MU\_COMPARISON\_VERIFIED}\):
Proposition~\ref{prop:quotient-after-correspondence-theta-mu-comparison}
constructs the operation-comparison isomorphisms
\(\theta^Q_{\mu,R}\) for quotient-presented finite-height operation
rows and proves \(\mathfrak o^{Q,\theta}_R=0\) on those rows.  This
packet imports associativity compatibility from
Proposition~\ref{prop:quotient-after-correspondence-theta-mu-associativity}.
It imports coproduct compatibility from
Proposition~\ref{prop:quotient-after-correspondence-theta-mu-coproduct}.
It imports primitive-projection compatibility from
Proposition~\ref{prop:quotient-after-correspondence-theta-mu-primitives}.
It imports the HN-transition compatibility of \(Q_{E,R}\) from
Proposition~\ref{prop:quotient-after-correspondence-hn-transition-compatibility}.
It does not prove Hall-operation transition compatibility, protected
integration, or aggregate population.

\begin{proposition}[Associativity compatibility of the quotient comparison maps]
\label{prop:quotient-after-correspondence-theta-mu-associativity}
\textnormal{[proved for the eight two-step flag associativity rows]}
Fix \(R\), and let
\[
w=\epsilon_1\epsilon_2\epsilon_3\in\{\mathrm L,\mathrm W\}^3
\]
be one of the eight local/wrapped words with retained ordered colours
\((\xi_1,\xi_2,\xi_3)\).  Assume the two-step flag stack
\(\mathfrak F^{(2),w}_{\xi_1,\xi_2,\xi_3,R}\) and the wordwise
associativity row of
Propositions~\ref{prop:lll-associativity}--\ref{prop:www-associativity}.
Let
\[
A_w=
\mu_{\xi_1+\xi_2,\xi_3,R}
\circ(\mu_{\xi_1,\xi_2,R}\otimes 1),
\qquad
B_w=
\mu_{\xi_1,\xi_2+\xi_3,R}
\circ(1\otimes\mu_{\xi_2,\xi_3,R})
\]
be the two unreduced parenthesized operations, and let
\[
a_w:A_w\xrightarrow{\sim}B_w
\]
be the associativity isomorphism represented by the common flag stack.
Let \(\overline A_w,\overline B_w\), and
\(\overline a_w:\overline A_w\xrightarrow{\sim}\overline B_w\), be
the corresponding reduced operations after quotient-after-
correspondence.  If every binary and composite operation occurring in
\(A_w\) and \(B_w\) carries the quotient-presented data required in
Proposition~\ref{prop:quotient-after-correspondence-theta-mu-comparison},
then the square
\[
\begin{tikzcd}
Q_{E,R}A_w \ar[r,"\Theta^L_w"]
\ar[d,"Q_{E,R}(a_w)"']&
\overline A_w\,(Q_{E,R}^{\boxtimes3})
\ar[d,"\overline a_w"]\\
Q_{E,R}B_w \ar[r,"\Theta^R_w"']&
\overline B_w\,(Q_{E,R}^{\boxtimes3})
\end{tikzcd}
\]
commutes.  Here \(\Theta^L_w\) and \(\Theta^R_w\) are obtained by
composing the binary \(\theta^Q_{\mu,R}\) maps of
Proposition~\ref{prop:quotient-after-correspondence-theta-mu-comparison}
with the composition comparison of
Proposition~\ref{prop:quotient-after-correspondence-preserves-composition}.
Consequently the quotient two-step flag associativity component
vanishes:
\[
\mathfrak o^{Q,\mathrm{flag}}_R=0 .
\]
This proves only compatibility of \(\theta^Q_{\mu,R}\) with the
eight length-three associativity rows.  It does not prove coproduct
compatibility, primitive compatibility, \(R'\to R\) transition
compatibility, protected-integration compatibility, or aggregate
hybrid-carrier population.
\end{proposition}

\begin{proof}
For a fixed \(w\), both parenthesizations are represented before
quotient by the same retained two-step flag stack
\[
\mathfrak F^{(2),w}_{\xi_1,\xi_2,\xi_3,R}.
\]
The maps \(\lambda^w\) and \(\rho^w\) of
Proposition~\ref{prop:eight-two-step-flag-stacks} identify this flag
stack with the two iterated correspondence fibre products used by
\(A_w\) and \(B_w\).  The wordwise associativity propositions state
that the two pull-push operations obtained from this common flag stack
agree after the required base-change, projection-formula, and reduced
Thom--Sebastiani identifications.

Apply the quotient-after-correspondence construction to the same
diagram.  Proposition~\ref{prop:quotient-after-correspondence-preserves-composition}
identifies the quotient of each iterated correspondence fibre product
with the corresponding iterated quotient correspondence.
Proposition~\ref{prop:quotient-after-correspondence-preserves-base-change}
identifies the quotient of the common flag square with the reduced
cartesian square, and
Proposition~\ref{prop:quotient-after-correspondence-preserves-thom-sebastiani}
descends the coefficient transports.  Thus the reduced associator
\(\overline a_w\) is the image of \(a_w\) under quotient descent.

The maps \(\Theta^L_w\) and \(\Theta^R_w\) are induced by the same
quotient-presented compact-support pushforward descent isomorphisms
\(\delta_{q,e}\) used in
Proposition~\ref{prop:quotient-after-correspondence-theta-mu-comparison},
pulled back to the common quotient flag stack.  Full faithfulness of
constructible descent on quotient stacks makes the two morphisms in
the displayed square equal after pullback to
\(\mathfrak F^{(2),w}_{\xi_1,\xi_2,\xi_3,R}\).  Since descent is
faithful, the square commutes on the quotient stack.  The argument is
wordwise for the eight length-three words; it does not involve the
Hall coproduct, primitive projection, finite-height transitions, or
protected integration.
\end{proof}

The packet
\[
\texttt{certificates/hybrid/quotient\_after\_correspondence\_theta\_mu\_associativity}
\]
verifies
\(\texttt{QUOTIENT\_AFTER\_CORRESPONDENCE\_THETA\_MU\_ASSOCIATIVITY\_VERIFIED}\):
Proposition~\ref{prop:quotient-after-correspondence-theta-mu-associativity}
proves that the quotient comparison maps \(\theta^Q_{\mu,R}\) commute
with the eight wordwise two-step flag associativity rows, hence
\(\mathfrak o^{Q,\mathrm{flag}}_R=0\) for those rows.  This packet
imports coproduct compatibility from
Proposition~\ref{prop:quotient-after-correspondence-theta-mu-coproduct}.
It imports primitive-projection compatibility from
Proposition~\ref{prop:quotient-after-correspondence-theta-mu-primitives}.
It imports the HN-transition compatibility of \(Q_{E,R}\) from
Proposition~\ref{prop:quotient-after-correspondence-hn-transition-compatibility}.
It does not prove Hall-operation transition compatibility, protected
integration, or aggregate population.

\begin{proposition}[Coproduct compatibility of the quotient comparison maps]
\label{prop:quotient-after-correspondence-theta-mu-coproduct}
\textnormal{[proved for supplied quotient-presented Hall bialgebra rows]}
Fix \(R\).  Let \(\mu_R\) be a retained binary Hall operation row, and
let
\[
\Delta_R:
\mathcal H_R^{E,\mathrm{hyb}}
\longrightarrow
\mathcal H_R^{E,\mathrm{hyb}}\boxtimes
\mathcal H_R^{E,\mathrm{hyb}}
\]
be a retained collision coproduct row.  Assume the product
correspondence, the splitting correspondence defining \(\Delta_R\),
and the four-leg bialgebra comparison correspondence carry the
quotient-presented compact-support, base-change, and reduced
Thom--Sebastiani data required in
Proposition~\ref{prop:quotient-after-correspondence-theta-mu-comparison}.
Assume also that the source finite Hall bialgebra identity
\[
\Delta_R\mu_R
=
(\mu_R\boxtimes\mu_R)(1\boxtimes\tau\boxtimes1)
(\Delta_R\boxtimes\Delta_R)
\]
has defect rank zero on that retained correspondence, with Koszul sign
braiding \(\tau\).  Let \(\overline\mu_R\), \(\overline\Delta_R\), and
the reduced bialgebra identity be the quotient-after-correspondence
images of these rows.

Then the two morphisms
\[
(Q_{E,R}\boxtimes Q_{E,R})\Delta_R\mu_R
\rightrightarrows
(\overline\mu_R\boxtimes\overline\mu_R)
(1\boxtimes\tau\boxtimes1)
(\overline\Delta_R\boxtimes\overline\Delta_R)
Q_{E,R}^{\boxtimes2}
\]
coincide.  The first morphism applies the quotient comparison for the
coproduct row, then \(\overline\Delta_R(\theta^Q_{\mu,R})\), and then
the reduced bialgebra identity.  The second applies the source
bialgebra identity, then the two coproduct comparison maps and the two
product comparison maps \(\theta^Q_{\mu,R}\).  Consequently the
quotient product-coproduct compatibility defect vanishes:
\[
\mathfrak o^{Q,\Delta}_R=0 .
\]
This proves only compatibility of \(\theta^Q_{\mu,R}\) with supplied
finite coproduct/bialgebra rows.  It does not construct the compact
Hall coproduct in the aggregate carrier, prove primitive compatibility,
\(R'\to R\) transition compatibility, protected integration, or
aggregate hybrid-carrier population.
\end{proposition}

\begin{proof}
The source bialgebra identity is represented by the retained
four-leg extension-splitting correspondence: both
\(\Delta_R\mu_R\) and
\((\mu_R\boxtimes\mu_R)(1\boxtimes\tau\boxtimes1)
(\Delta_R\boxtimes\Delta_R)\) are its two pull-push descriptions.
The zero-defect source row asserts equality after the listed
base-change, projection-formula, compact-support, and
Thom--Sebastiani identifications.

Apply quotient-after-correspondence to the same correspondence.
Proposition~\ref{prop:quotient-after-correspondence-preserves-composition}
identifies the quotient of the composite correspondence with the
composite of the quotient correspondences.
Proposition~\ref{prop:quotient-after-correspondence-preserves-base-change}
identifies the quotient of the bialgebra comparison squares with the
reduced cartesian squares, and
Proposition~\ref{prop:quotient-after-correspondence-preserves-thom-sebastiani}
descends the product and inverse-coproduct coefficient transports.
Thus the reduced bialgebra identity is the quotient image of the
source identity.

The comparison maps for the product and coproduct rows are induced by
the same quotient-presented compact-support descent isomorphisms
\(\delta_{q,e}\) used in
Proposition~\ref{prop:quotient-after-correspondence-theta-mu-comparison},
with source and target legs interchanged for the coproduct
correspondence.  After pullback to the common equivariant
four-leg correspondence, the two displayed composites are the same
morphism because the source bialgebra square has defect rank zero.
Full faithfulness of constructible descent on quotient stacks then
identifies the two descended morphisms on the quotient
correspondence.  The argument uses only the supplied finite
product-coproduct square; it does not form primitive kernels,
finite-height transition maps, or the protected integration map.
\end{proof}

The packet
\[
\texttt{certificates/hybrid/quotient\_after\_correspondence\_theta\_mu\_coproduct}
\]
verifies
\(\texttt{QUOTIENT\_AFTER\_CORRESPONDENCE\_THETA\_MU\_COPRODUCT\_VERIFIED}\):
Proposition~\ref{prop:quotient-after-correspondence-theta-mu-coproduct}
proves that the quotient comparison maps \(\theta^Q_{\mu,R}\) commute
with supplied finite Hall coproduct/bialgebra rows.  This packet
imports primitive-projection compatibility from
Proposition~\ref{prop:quotient-after-correspondence-theta-mu-primitives}.
It imports the HN-transition compatibility of \(Q_{E,R}\) from
Proposition~\ref{prop:quotient-after-correspondence-hn-transition-compatibility}.
It does not prove Hall-coproduct transition compatibility, protected
integration, or aggregate population.

\begin{proposition}[Primitive-projection compatibility of the quotient comparison maps]
\label{prop:quotient-after-correspondence-theta-mu-primitives}
\textnormal{[proved for supplied finite primitive kernel and projection rows]}
Fix \(R\).  Let
\[
\widetilde\Delta_R
=\Delta_R-\eta_R\epsilon_R\boxtimes\mathrm{id}
-\mathrm{id}\boxtimes\eta_R\epsilon_R,
\qquad
\widetilde{\overline\Delta}_R
=\overline\Delta_R-\overline\eta_R\overline\epsilon_R\boxtimes\mathrm{id}
-\mathrm{id}\boxtimes\overline\eta_R\overline\epsilon_R
\]
be the reduced coproducts attached to the retained finite Hall
bialgebra rows.  Assume the finite primitive rows supply the kernel
subobjects and idempotent projections
\[
P_R^E=\ker(\widetilde\Delta_R)\cap\ker(\epsilon_R),
\qquad
\overline P_R=\ker(\widetilde{\overline\Delta}_R)
\cap\ker(\overline\epsilon_R),
\]
\[
\pi_R^{\Prim}\colon\mathcal H_R^{E,\mathrm{hyb}}\to P_R^E,
\qquad
\overline\pi_R^{\Prim}\colon
\mathcal H_R^{\mathrm{red},\mathrm{hyb}}\to\overline P_R,
\]
and assume that the quotient-after-correspondence descent of the
unit, counit, and reduced coproduct rows identifies
\[
Q_{E,R}P_R^E\simeq \overline P_R,\qquad
Q_{E,R}\pi_R^{\Prim}
=\overline\pi_R^{\Prim}Q_{E,R}
\]
on the supplied primitive kernel/projection row.  For a retained
binary Hall product row \(\mu_R\), let \(\theta^Q_{\mu,R}\) be the
comparison map of
Proposition~\ref{prop:quotient-after-correspondence-theta-mu-comparison},
and assume the coproduct compatibility of
Proposition~\ref{prop:quotient-after-correspondence-theta-mu-coproduct}
for the same product-coproduct row.

Then \(\theta^Q_{\mu,R}\) commutes with the primitive projection on
the output Hall object:
\[
\overline\pi_R^{\Prim}\,
\overline\mu_R\,Q_{E,R}^{\boxtimes2}\,\theta^Q_{\mu,R}
=
\theta^{Q,\Prim}_{\mu,R}\,
Q_{E,R}\pi_R^{\Prim}\,\mu_R ,
\]
where \(\theta^{Q,\Prim}_{\mu,R}\colon Q_{E,R}P_R^E\to
\overline P_R\) is the quotient identification on the primitive
summand.  Equivalently, the square
\[
\begin{tikzcd}
Q_{E,R}\mu_R \ar[r,"\theta^Q_{\mu,R}"]
\ar[d,"Q_{E,R}\pi_R^{\Prim}"'] &
\overline\mu_R Q_{E,R}^{\boxtimes2}
\ar[d,"\overline\pi_R^{\Prim}"]\\
Q_{E,R}P_R^E \ar[r,"\theta^{Q,\Prim}_{\mu,R}"'] &
\overline P_R
\end{tikzcd}
\]
commutes.  Consequently the quotient primitive-projection defect
vanishes on the supplied finite rows:
\[
\mathfrak o^{Q,\Prim}_R=0 .
\]
This proves only compatibility of \(\theta^Q_{\mu,R}\) with the
supplied primitive kernel/projection rows.  It does not construct
source compact Hall primitive matrices, prove \(R'\to R\) transition
compatibility, prove protected-integration compatibility, or populate the
aggregate hybrid carrier.
\end{proposition}

\begin{proof}
The primitive subobject is the intersection of two kernels: the kernel
of the reduced coproduct and the kernel of the counit.  By
Proposition~\ref{prop:quotient-after-correspondence-theta-mu-coproduct},
the quotient comparison sends the source reduced-coproduct equation to
the reduced quotient equation.  The primitive row additionally
supplies the unit and counit descent, so \(Q_{E,R}\) sends
\(\ker(\widetilde\Delta_R)\cap\ker(\epsilon_R)\) into
\(\ker(\widetilde{\overline\Delta}_R)\cap
\ker(\overline\epsilon_R)\).  This gives the kernel identification
\(Q_{E,R}P_R^E\simeq\overline P_R\) recorded in the hypothesis.

The two composites in the displayed square are morphisms obtained from
the same quotient-presented binary Hall correspondence after adjoining
the supplied primitive idempotent on the output.  On the coproduct
kernel their reduced-coproduct components agree by the row-254
coproduct square; on the augmentation kernel their counit components
agree by the supplied counit descent.  Full faithfulness of
constructible descent on the quotient stack identifies the two
idempotent-split morphisms.  Hence the primitive projection square
commutes and the finite primitive defect has rank zero.
\end{proof}

The packet
\[
\texttt{certificates/hybrid/quotient\_after\_correspondence\_theta\_mu\_primitives}
\]
verifies
\(\texttt{QUOTIENT\_AFTER\_CORRESPONDENCE\_THETA\_MU\_PRIMITIVES\_VERIFIED}\):
Proposition~\ref{prop:quotient-after-correspondence-theta-mu-primitives}
proves that the quotient comparison maps \(\theta^Q_{\mu,R}\) commute
with the supplied finite primitive kernel and projection rows.  This
packet imports the comparison, associativity, and coproduct packets.
It imports the HN-transition compatibility of \(Q_{E,R}\) from
Proposition~\ref{prop:quotient-after-correspondence-hn-transition-compatibility}.
It does not prove HN-transition preservation of primitive subspaces,
protected-integration compatibility, or aggregate population.

\begin{proposition}[HN-transition compatibility of the quotient descent functor]
\label{prop:quotient-after-correspondence-hn-transition-compatibility}
\textnormal{[proved for supplied finite HN transition rows]}
Let \(R\le R'\).  Suppose the finite quotient-after-correspondence
data at heights \(R'\) and \(R\) are supplied, and suppose every
retained object row \(Y_{R'}\to Y_R\) in the HN transition system is
\(E\)-equivariant.  Write
\[
T^E_{R'R}:\mathsf{BM}^{E,\mathrm{hyb}}_{R'}
\longrightarrow
\mathsf{BM}^{E,\mathrm{hyb}}_{R},
\qquad
\overline T_{R'R}:\mathsf{BM}^{\mathrm{red},\mathrm{hyb}}_{R'}
\longrightarrow
\mathsf{BM}^{\mathrm{red},\mathrm{hyb}}_{R}
\]
for the supplied source and reduced Borel--Moore transition functors.
Assume that each transition row is quotient-presented: the square
\[
\begin{array}{ccc}
Y_{R'} & \longrightarrow & Y_R\\
\downarrow && \downarrow\\
\left[Y_{R'}/E\right] & \longrightarrow & \left[Y_R/E\right]
\end{array}
\]
commutes, the lower map is the quotient of the upper \(E\)-equivariant
map, and the vanishing-cycle coefficient transport at height
\(R'\to R\) descends to the quotient coefficient transport.  Equivalently,
the source transition chain map and the reduced transition chain map
are induced from the same equivariant transition correspondence.

Then the quotient descent functors commute with the HN transition:
\[
\overline T_{R'R}\,Q_{E,R'}
=
Q_{E,R}\,T^E_{R'R}
\]
as functors on the quotient-presented finite Borel--Moore chain
category.  Consequently the transition component of the quotient
residual vanishes on the supplied rows:
\[
\mathfrak o^{Q,\mathrm{tr}}_{R'R}=0.
\]
This proves only the HN-transition compatibility of \(Q_{E,R}\) on
quotient-presented finite transition rows.  It does not prove that HN
transitions preserve the Hall product, Hall coproduct, primitive
subspaces, protected-integration compatibility, or aggregate
hybrid-carrier population.
\end{proposition}

\begin{proof}
For a retained transition object \(Y_{R'}\to Y_R\), the map is
\(E\)-equivariant.  The quotient-stack universal property therefore
gives a unique reduced transition map
\[
[Y_{R'}/E]\longrightarrow [Y_R/E]
\]
for which the quotient square commutes.  The coefficient transport is
assumed to be an \(E\)-equivariant morphism between quotient-presented
vanishing-cycle coefficients, and hence descends uniquely to the
corresponding coefficient transport on the two quotient stacks.

Apply the Borel--Moore chain construction of
Proposition~\ref{prop:quotient-after-correspondence-borel-moore-chain-functor}.
The composite
\(\overline T_{R'R}Q_{E,R'}\) first replaces the source object by the
quotient stack \([Y_{R'}/E]\) and then applies the descended transition
map to \([Y_R/E]\).  The composite \(Q_{E,R}T^E_{R'R}\) first applies
the equivariant source transition and then descends the result to the
same quotient stack \([Y_R/E]\).  Both composites are therefore the
same functor
\[
C_*^{\mathrm{BM}}([Y_{R'}/E],\overline K_{R'})
\longrightarrow
C_*^{\mathrm{BM}}([Y_R/E],\overline K_R)
\]
with the same descended coefficient morphism.  Full faithfulness of
constructible descent on quotient stacks identifies the two chain maps
and their composites over successive height restrictions.  Thus the
finite transition defect of \(Q_{E,R}\) has rank zero.
\end{proof}

The packet
\[
\texttt{certificates/hybrid/quotient\_after\_correspondence\_hn\_transition\_compatibility}
\]
verifies
\(\texttt{QUOTIENT\_AFTER\_CORRESPONDENCE\_HN\_TRANSITION\_COMPATIBILITY\_VERIFIED}\):
Proposition~\ref{prop:quotient-after-correspondence-hn-transition-compatibility}
proves that \(Q_{E,R}\) commutes with supplied finite HN transition
rows.  It does not prove Hall product, Hall coproduct,
primitive-subspace, protected-integration, or aggregate transition
compatibility.

\begin{definition}[Finite protected integration datum]
\label{def:finite-protected-integration-datum}
Fix \(R\), the reduced Borel--Moore chain object
\[
\mathcal H_R^{\mathrm{red},\mathrm{hyb}}
:=
Q_{E,R}\mathcal F_{X,\sigma,S,\le R}^{\mathrm{hyb}},
\]
and its finite charge decomposition
\[
\mathcal H_R^{\mathrm{red},\mathrm{hyb}}
=\bigoplus_{\gamma\in\Gamma_R}
\mathcal H_{\gamma,R}^{\mathrm{red},\mathrm{hyb}} .
\]
A finite protected integration datum at height \(R\) consists of:
\begin{enumerate}
\item a finite commutative monoid \(\mathsf T_R\) and labels
\(\omega_R(\gamma)\in\mathsf T_R\) for the retained charges;
\item the finite subspace
\[
\C[\mathsf T_R]_{\le R}
:=\sum_{\gamma\in\Gamma_R}\C e_{\omega_R(\gamma)}
\subset \C[\mathsf T_R];
\]
\item for every \(\gamma\in\Gamma_R\), a degree-zero chain map
\[
\int^{\mathrm{prot}}_{\gamma,R}\colon
\mathcal H_{\gamma,R}^{\mathrm{red},\mathrm{hyb}}
\longrightarrow
\C e_{\omega_R(\gamma)} .
\]
\end{enumerate}
The associated protected integration map is the finite direct sum
\[
I_R^{\mathrm{prot}}
:=
\bigoplus_{\gamma\in\Gamma_R}
\int^{\mathrm{prot}}_{\gamma,R}\colon
\mathcal H_R^{\mathrm{red},\mathrm{hyb}}
\longrightarrow
\C[\mathsf T_R]_{\le R}.
\]
Thus, for \(x\in\mathcal H_{\gamma,R}^{\mathrm{red},\mathrm{hyb}}\),
\[
I_R^{\mathrm{prot}}(x)\in \C e_{\omega_R(\gamma)}.
\]
The label \(\omega_R(\gamma)\) is not yet identified with
\(\overline\Pi_X(\gamma)\), and \(e_{\omega_R(\gamma)}\) is not yet
identified with \(q^nr^ls^m\).  Those are the Gram-degree and
\(\overline\Pi_X\)-separation rows.  This datum also does not assert
multiplicativity, comultiplicativity, primitive-projection
compatibility, HN-transition compatibility, the closed
\(K3\times E\to\mathrm{pt}\) protected trace, or a Pfaffian line.
\end{definition}

\begin{proposition}[Definition row for protected integration]
\label{prop:protected-integration-definition}
\textnormal{[proved for supplied finite chargewise integration maps]}
The datum of Definition~\ref{def:finite-protected-integration-datum}
defines a degree-zero chain map
\[
I_R^{\mathrm{prot}}\colon
Q_{E,R}\mathcal F_{X,\sigma,S,\le R}^{\mathrm{hyb}}
\longrightarrow
\C[\mathsf T_R]_{\le R}.
\]
Equivalently, the chain-map component of the protected-integration
residual vanishes on the supplied finite rows:
\[
\mathfrak o^{I,\mathrm{ch}}_R=0.
\]
This proves only the definition and well-typedness of
\(I_R^{\mathrm{prot}}\).  It does not prove product compatibility,
coproduct compatibility, primitive-projection compatibility,
HN-transition compatibility, the Gram exponent law, the use of
\(\overline\Pi_X\) rather than \(b_R^{\mathrm{geom}}\), or aggregate
hybrid-carrier population.
\end{proposition}

\begin{proof}
The charge set \(\Gamma_R\) is finite at height \(R\).  Each summand
\(\int^{\mathrm{prot}}_{\gamma,R}\) is a degree-zero chain map to the
one-dimensional complex \(\C e_{\omega_R(\gamma)}\).  The finite direct
sum of these chain maps is therefore a degree-zero chain map from the
finite direct sum
\(\bigoplus_{\gamma\in\Gamma_R}
\mathcal H_{\gamma,R}^{\mathrm{red},\mathrm{hyb}}\)
to the finite subspace of the monoid algebra spanned by the labels
\(e_{\omega_R(\gamma)}\).  This map is exactly
\(I_R^{\mathrm{prot}}\).  No Hall operation, coproduct, primitive
idempotent, transition functor, Gram-coordinate comparison, or closed
trace functor enters the construction.
\end{proof}

The packet
\[
\texttt{certificates/hybrid/protected\_integration\_definition}
\]
verifies
\(\texttt{PROTECTED\_INTEGRATION\_DEFINITION\_VERIFIED}\):
Proposition~\ref{prop:protected-integration-definition} defines
\(I_R^{\mathrm{prot}}\) as the finite chargewise direct sum of supplied
degree-zero protected integration maps.  Product compatibility is
supplied separately by
Proposition~\ref{prop:protected-integration-product-compatibility}.
Coproduct compatibility is supplied separately by
Proposition~\ref{prop:protected-integration-coproduct-compatibility}.
Primitive-projection compatibility is supplied separately by
Proposition~\ref{prop:protected-integration-primitive-compatibility}.
Gram degree is supplied separately by
Proposition~\ref{prop:protected-integration-gram-degree}.
The normal-ordered label comparison is supplied separately by
Proposition~\ref{prop:protected-integration-pi-x-separation}.  The
packet records transition compatibility as a separate open row.

\begin{proposition}[Product compatibility of protected integration]
\label{prop:protected-integration-product-compatibility}
\textnormal{[proved for supplied finite binary product rows]}
Fix \(R\) and the finite protected integration datum of
Definition~\ref{def:finite-protected-integration-datum}.  Let
\[
w=\epsilon_1\epsilon_2\in\{\mathrm L,\mathrm W\}^2
\]
be a retained binary local/wrapped word, and let
\[
\overline\mu_{w,\gamma_1,\gamma_2,R}\colon
\mathcal H_{\gamma_1,R}^{\mathrm{red},\mathrm{hyb}}
\otimes
\mathcal H_{\gamma_2,R}^{\mathrm{red},\mathrm{hyb}}
\longrightarrow
\mathcal H_{\gamma_1+\gamma_2,R}^{\mathrm{red},\mathrm{hyb}}
\]
be the corresponding reduced Hall product row after
quotient-after-correspondence and the comparison map
\(\theta^Q_{\mu,R}\).  Assume the target monoid labels satisfy the
supplied product row
\[
m_{\mathsf T_R}\bigl(
e_{\omega_R(\gamma_1)}\otimes e_{\omega_R(\gamma_2)}
\bigr)
=e_{\omega_R(\gamma_1+\gamma_2)}
\]
with defect rank zero.  Assume also that the chargewise protected
integration functionals satisfy the supplied square
\[
\begin{tikzcd}
\mathcal H_{\gamma_1,R}^{\mathrm{red},\mathrm{hyb}}
\otimes
\mathcal H_{\gamma_2,R}^{\mathrm{red},\mathrm{hyb}}
\ar[r,"\overline\mu_{w,\gamma_1,\gamma_2,R}"]
\ar[d,"\int^{\mathrm{prot}}_{\gamma_1,R}\otimes
\int^{\mathrm{prot}}_{\gamma_2,R}"']&
\mathcal H_{\gamma_1+\gamma_2,R}^{\mathrm{red},\mathrm{hyb}}
\ar[d,"\int^{\mathrm{prot}}_{\gamma_1+\gamma_2,R}"]\\
\C e_{\omega_R(\gamma_1)}\otimes
\C e_{\omega_R(\gamma_2)}
\ar[r,"m_{\mathsf T_R}"']&
\C e_{\omega_R(\gamma_1+\gamma_2)}
\end{tikzcd}
\]
with defect rank zero.  Then the protected integration map is
multiplicative on this product row:
\[
I_R^{\mathrm{prot}}\,
\overline\mu_{w,\gamma_1,\gamma_2,R}
=
m_{\mathsf T_R}\,
(I_R^{\mathrm{prot}}\otimes I_R^{\mathrm{prot}})
\]
on
\(\mathcal H_{\gamma_1,R}^{\mathrm{red},\mathrm{hyb}}\otimes
\mathcal H_{\gamma_2,R}^{\mathrm{red},\mathrm{hyb}}\).  Consequently
\[
\mathfrak o^{I,\mathrm{prod}}_{w,R}=0
\]
for every supplied finite binary product row of word \(w\).
This proves only product compatibility of
\(I_R^{\mathrm{prot}}\).  It does not prove coproduct compatibility,
primitive-projection compatibility, HN-transition compatibility, the
Gram exponent law, the use of \(\overline\Pi_X\) rather than
\(b_R^{\mathrm{geom}}\), the level-\(\mathsf Z\) protected trace, or a
Pfaffian line.
\end{proposition}

\begin{proof}
Restrict \(I_R^{\mathrm{prot}}\) to the two source charge summands and
to the target charge summand.  By
Definition~\ref{def:finite-protected-integration-datum}, these
restrictions are exactly
\[
\int^{\mathrm{prot}}_{\gamma_1,R},\qquad
\int^{\mathrm{prot}}_{\gamma_2,R},\qquad
\int^{\mathrm{prot}}_{\gamma_1+\gamma_2,R}.
\]
The displayed component square is therefore precisely the restriction
of the global multiplicativity square for \(I_R^{\mathrm{prot}}\) to
the summand
\(\mathcal H_{\gamma_1,R}^{\mathrm{red},\mathrm{hyb}}\otimes
\mathcal H_{\gamma_2,R}^{\mathrm{red},\mathrm{hyb}}\).  Since
\(\Gamma_R\) is finite, the global product square is the finite direct
sum of these component squares.  Each component has defect rank zero by
hypothesis, and the target label multiplication identifies the
one-dimensional target line with
\(\C e_{\omega_R(\gamma_1+\gamma_2)}\).  Hence the product defect of
protected integration vanishes on the supplied row.  The proof uses no
coproduct, primitive projection, transition functor, Gram-coordinate
identification, closed trace functor, or Pfaffian line.
\end{proof}

The packet
\[
\texttt{certificates/hybrid/protected\_integration\_product\_compatibility}
\]
verifies
\(\texttt{PROTECTED\_INTEGRATION\_PRODUCT\_COMPATIBILITY\_VERIFIED}\):
Proposition~\ref{prop:protected-integration-product-compatibility}
proves that \(I_R^{\mathrm{prot}}\) respects every supplied finite
binary local/wrapped product row.  Coproduct compatibility is supplied
separately by
Proposition~\ref{prop:protected-integration-coproduct-compatibility}.
Primitive-projection compatibility is supplied separately by
Proposition~\ref{prop:protected-integration-primitive-compatibility}.
Gram degree is supplied separately by
Proposition~\ref{prop:protected-integration-gram-degree}.
The normal-ordered label comparison is supplied separately by
Proposition~\ref{prop:protected-integration-pi-x-separation}.  The
packet records transition compatibility as a separate open row.

\begin{proposition}[Coproduct compatibility of protected integration]
\label{prop:protected-integration-coproduct-compatibility}
\textnormal{[proved for supplied finite binary coproduct rows]}
Fix \(R\) and the finite protected integration datum of
Definition~\ref{def:finite-protected-integration-datum}.  Let
\[
w=\epsilon_1\epsilon_2\in\{\mathrm L,\mathrm W\}^2
\]
be a retained binary local/wrapped output word, and let
\[
\overline\Delta_{w,\gamma_1,\gamma_2,R}\colon
\mathcal H_{\gamma_1+\gamma_2,R}^{\mathrm{red},\mathrm{hyb}}
\longrightarrow
\mathcal H_{\gamma_1,R}^{\mathrm{red},\mathrm{hyb}}
\otimes
\mathcal H_{\gamma_2,R}^{\mathrm{red},\mathrm{hyb}}
\]
be the corresponding reduced Hall coproduct row after
quotient-after-correspondence and the coproduct comparison row
\(\theta^Q_{\Delta,R}\).  Assume the target label coalgebra has the
supplied component row
\[
\Delta_{\mathsf T_R,w}\bigl(e_{\omega_R(\gamma_1+\gamma_2)}\bigr)
=e_{\omega_R(\gamma_1)}\otimes e_{\omega_R(\gamma_2)}
\]
with defect rank zero.  Assume also that the chargewise protected
integration functionals satisfy the supplied square
\[
\begin{tikzcd}
\mathcal H_{\gamma_1+\gamma_2,R}^{\mathrm{red},\mathrm{hyb}}
\ar[r,"\overline\Delta_{w,\gamma_1,\gamma_2,R}"]
\ar[d,"\int^{\mathrm{prot}}_{\gamma_1+\gamma_2,R}"']&
\mathcal H_{\gamma_1,R}^{\mathrm{red},\mathrm{hyb}}
\otimes
\mathcal H_{\gamma_2,R}^{\mathrm{red},\mathrm{hyb}}
\ar[d,"\int^{\mathrm{prot}}_{\gamma_1,R}\otimes
\int^{\mathrm{prot}}_{\gamma_2,R}"]\\
\C e_{\omega_R(\gamma_1+\gamma_2)}
\ar[r,"\Delta_{\mathsf T_R,w}"']&
\C e_{\omega_R(\gamma_1)}\otimes
\C e_{\omega_R(\gamma_2)}
\end{tikzcd}
\]
with defect rank zero.  Then the protected integration map is
comultiplicative on this coproduct row:
\[
(I_R^{\mathrm{prot}}\otimes I_R^{\mathrm{prot}})\,
\overline\Delta_{w,\gamma_1,\gamma_2,R}
=
\Delta_{\mathsf T_R,w}\,I_R^{\mathrm{prot}}
\]
on
\(\mathcal H_{\gamma_1+\gamma_2,R}^{\mathrm{red},\mathrm{hyb}}\).
Consequently
\[
\mathfrak o^{I,\mathrm{co}}_{w,R}=0
\]
for every supplied finite binary coproduct row of word \(w\).
This proves only coproduct compatibility of
\(I_R^{\mathrm{prot}}\).  It does not prove primitive-projection
compatibility, HN-transition compatibility, the Gram exponent law, the
use of \(\overline\Pi_X\) rather than \(b_R^{\mathrm{geom}}\), the
level-\(\mathsf Z\) protected trace, or a Pfaffian line.
\end{proposition}

\begin{proof}
Restrict \(I_R^{\mathrm{prot}}\) to the source charge summand and to
the two target charge summands.  By
Definition~\ref{def:finite-protected-integration-datum}, these
restrictions are
\[
\int^{\mathrm{prot}}_{\gamma_1+\gamma_2,R},\qquad
\int^{\mathrm{prot}}_{\gamma_1,R},\qquad
\int^{\mathrm{prot}}_{\gamma_2,R}.
\]
The displayed component square is therefore the restriction of the
global coproduct square for \(I_R^{\mathrm{prot}}\) to the charge
\(\gamma_1+\gamma_2\) and the output split
\((\gamma_1,\gamma_2)\).  Since \(\Gamma_R\) is finite, the supplied
global coproduct square is the finite direct sum of these component
squares over the retained split rows.  Each component has defect rank
zero by hypothesis, and the target label coproduct identifies the
one-dimensional source label line with
\(\C e_{\omega_R(\gamma_1)}\otimes\C e_{\omega_R(\gamma_2)}\).  Hence
the coproduct defect of protected integration vanishes on the supplied
row.  The proof uses no primitive projection, transition functor,
Gram-coordinate identification, closed trace functor, or Pfaffian line.
\end{proof}

The packet
\[
\texttt{certificates/hybrid/protected\_integration\_coproduct\_compatibility}
\]
verifies
\(\texttt{PROTECTED\_INTEGRATION\_COPRODUCT\_COMPATIBILITY\_VERIFIED}\):
Proposition~\ref{prop:protected-integration-coproduct-compatibility}
proves that \(I_R^{\mathrm{prot}}\) respects every supplied finite
binary local/wrapped coproduct row.  Primitive-projection compatibility
is supplied separately by
Proposition~\ref{prop:protected-integration-primitive-compatibility}.
Gram degree is supplied separately by
Proposition~\ref{prop:protected-integration-gram-degree}.
The normal-ordered label comparison is supplied separately by
Proposition~\ref{prop:protected-integration-pi-x-separation}.  The
packet records transition compatibility as a separate open row.

\begin{proposition}[Primitive-projection compatibility of protected integration]
\label{prop:protected-integration-primitive-compatibility}
\textnormal{[proved for supplied finite primitive projection rows]}
Fix \(R\) and the finite protected integration datum of
Definition~\ref{def:finite-protected-integration-datum}.  Assume the
quotient primitive-projection row of
Proposition~\ref{prop:quotient-after-correspondence-theta-mu-primitives}
has supplied the reduced primitive subobject
\[
\overline P_R
=\ker(\widetilde{\overline\Delta}_R)\cap
\ker(\overline\epsilon_R)
\subset \mathcal H_R^{\mathrm{red},\mathrm{hyb}}
\]
and an idempotent projection
\[
\overline\pi_R^{\Prim}\colon
\mathcal H_R^{\mathrm{red},\mathrm{hyb}}\longrightarrow \overline P_R .
\]
Assume also that the target label coalgebra has a supplied primitive
projector
\[
\pi_{\mathsf T_R}^{\Prim}\colon
\C[\mathsf T_R]_{\le R}
\longrightarrow
\C[\mathsf T_R]_{\le R}^{\Prim}
\]
with defect rank zero, and that the protected integration functionals
satisfy the supplied square
\[
\begin{tikzcd}
\mathcal H_R^{\mathrm{red},\mathrm{hyb}}
\ar[r,"\overline\pi_R^{\Prim}"]
\ar[d,"I_R^{\mathrm{prot}}"']&
\overline P_R
\ar[d,"I_R^{\mathrm{prot}}\vert_{\overline P_R}"]\\
\C[\mathsf T_R]_{\le R}
\ar[r,"\pi_{\mathsf T_R}^{\Prim}"']&
\C[\mathsf T_R]_{\le R}^{\Prim}
\end{tikzcd}
\]
with defect rank zero.  Then \(I_R^{\mathrm{prot}}\) respects the
primitive projection:
\[
I_R^{\mathrm{prot}}\,
\overline\pi_R^{\Prim}
=
\pi_{\mathsf T_R}^{\Prim}\,I_R^{\mathrm{prot}}
\]
on \(\mathcal H_R^{\mathrm{red},\mathrm{hyb}}\), after identifying
\(I_R^{\mathrm{prot}}\vert_{\overline P_R}\) with its target primitive
subspace.  Consequently
\[
\mathfrak o^{I,\Prim}_R=0
\]
on the supplied finite primitive projection row.  This proves only
primitive-projection compatibility of \(I_R^{\mathrm{prot}}\).  It does
not prove HN-transition compatibility, the Gram exponent law, the use
of \(\overline\Pi_X\) rather than \(b_R^{\mathrm{geom}}\), the
level-\(\mathsf Z\) protected trace, or a Pfaffian line.
\end{proposition}

\begin{proof}
The reduced primitive subobject \(\overline P_R\) is already the
idempotent-split summand supplied by
Proposition~\ref{prop:quotient-after-correspondence-theta-mu-primitives};
the present statement does not reconstruct it.  Restrict the
chargewise direct sum \(I_R^{\mathrm{prot}}\) to this summand and
compose with the target primitive projector.  The displayed square is
exactly the assertion that these two composites agree.  Its defect rank
is zero by hypothesis, hence the primitive-projection component of the
protected-integration residual vanishes on the supplied finite row.  No
transition functor, Gram-coordinate identification, closed trace
functor, or Pfaffian line enters the argument.
\end{proof}

The packet
\[
\texttt{certificates/hybrid/protected\_integration\_primitive\_compatibility}
\]
verifies
\(\texttt{PROTECTED\_INTEGRATION\_PRIMITIVE\_COMPATIBILITY\_VERIFIED}\):
Proposition~\ref{prop:protected-integration-primitive-compatibility}
proves that \(I_R^{\mathrm{prot}}\) respects the supplied finite
primitive projection row.  Gram degree is supplied separately by
Proposition~\ref{prop:protected-integration-gram-degree}.
The normal-ordered label comparison is supplied separately by
Proposition~\ref{prop:protected-integration-pi-x-separation}.  The
packet records transition compatibility as a separate open row.

\begin{proposition}[Gram degree of protected integration]
\label{prop:protected-integration-gram-degree}
\textnormal{[proved for supplied finite Gram-monomial label rows]}
Fix \(R\) and the finite protected integration datum of
Definition~\ref{def:finite-protected-integration-datum}.  Assume a
finite Gram-monomial target row supplies a map
\[
g_R:\mathsf T_R\longrightarrow \Z^3,\qquad
g_R(\omega_R(\gamma))=(n_R(\gamma),l_R(\gamma),m_R(\gamma)),
\]
and a monomialization map
\[
\chi_R:\C[\mathsf T_R]_{\le R}\longrightarrow
\C[q^{\pm1},r^{\pm1},s^{\pm1}]
\]
such that
\[
\chi_R\bigl(e_{\omega_R(\gamma)}\bigr)
=q^{n_R(\gamma)}r^{l_R(\gamma)}s^{m_R(\gamma)}
\]
on every retained charge label, with defect rank zero.  Then, for
every
\[
x\in\mathcal H_{\gamma,R}^{\mathrm{red},\mathrm{hyb}},
\]
one has
\[
(\chi_R I_R^{\mathrm{prot}})(x)
\in
\C\,q^{n_R(\gamma)}r^{l_R(\gamma)}s^{m_R(\gamma)} .
\]
Equivalently, the \(s\)-degree of the protected integration image is
\[
\deg_s(\chi_R I_R^{\mathrm{prot}}(x))=m_R(\gamma)
\]
whenever the scalar coefficient is nonzero.  Consequently
\[
\mathfrak o^{I,\deg}_R=0
\]
on the supplied finite Gram-monomial rows.  This proves only the
target Gram-degree statement for protected integration.  It does not
prove that \(g_R\omega_R\) is the normal-ordered charge map
\(\overline\Pi_X\), does not prove that the third coordinate is
different from \(b_R^{\mathrm{geom}}\), does not prove HN-transition
compatibility, and does not construct the level-\(\mathsf Z\) protected
trace or a Pfaffian line.
\end{proposition}

\begin{proof}
By Definition~\ref{def:finite-protected-integration-datum},
\[
I_R^{\mathrm{prot}}(x)\in \C e_{\omega_R(\gamma)}
\]
for \(x\in\mathcal H_{\gamma,R}^{\mathrm{red},\mathrm{hyb}}\).  Applying
the supplied monomialization map \(\chi_R\) sends the one-dimensional
target line \(\C e_{\omega_R(\gamma)}\) to the one-dimensional monomial
line
\[
\C q^{n_R(\gamma)}r^{l_R(\gamma)}s^{m_R(\gamma)} .
\]
Thus the image has the displayed \(q,r,s\)-degree, and in particular
\(s\)-degree \(m_R(\gamma)\) when the coefficient is nonzero.  The
argument uses only the target label row and the chargewise definition of
\(I_R^{\mathrm{prot}}\); it does not use a comparison between
\(\omega_R\) and \(\overline\Pi_X\), does not compare with
\(b_R^{\mathrm{geom}}\), and does not use transition or trace data.
\end{proof}

The packet
\[
\texttt{certificates/hybrid/protected\_integration\_gram\_degree}
\]
verifies
\(\texttt{PROTECTED\_INTEGRATION\_GRAM\_DEGREE\_VERIFIED}\):
Proposition~\ref{prop:protected-integration-gram-degree} proves that
the supplied finite monomialization sends the protected integration
image of charge \(\gamma\) to the line
\(\C q^{n_R(\gamma)}r^{l_R(\gamma)}s^{m_R(\gamma)}\).  The packet
records transition compatibility and
leaves the normal-ordered label comparison to
Proposition~\ref{prop:protected-integration-pi-x-separation}.

\begin{proposition}[Normal-ordered Gram label of protected integration]
\label{prop:protected-integration-pi-x-separation}
\textnormal{[proved for supplied finite normal-ordered label comparison rows]}
Fix \(R\) and the finite protected integration datum of
Definition~\ref{def:finite-protected-integration-datum}.  Assume the
Gram-degree row of
Proposition~\ref{prop:protected-integration-gram-degree}.  Assume also
that a finite protected-label comparison row supplies a lift
\[
\widehat\omega_R:\Gamma_R\longrightarrow \widehat\Gamma_R
\]
to the domain of the normal-ordered Gram map and a zero-defect equality
\[
g_R\omega_R
=
\overline\Pi_X\widehat\omega_R
\colon \Gamma_R\longrightarrow \Z^3 .
\]
Write
\[
\overline\Pi_X(\widehat\omega_R(\gamma))
=
\bigl(n_R^\Pi(\gamma),l_R^\Pi(\gamma),m_R^\Pi(\gamma)\bigr).
\]
Then, for every
\[
x\in\mathcal H_{\gamma,R}^{\mathrm{red},\mathrm{hyb}},
\]
one has
\[
(\chi_R I_R^{\mathrm{prot}})(x)
\in
\C q^{n_R^\Pi(\gamma)}r^{l_R^\Pi(\gamma)}s^{m_R^\Pi(\gamma)}
\]
and
\[
\deg_s(\chi_R I_R^{\mathrm{prot}}(x))
=m_R^\Pi(\gamma)
=\operatorname{pr}_3\overline\Pi_X(\widehat\omega_R(\gamma))
\]
whenever the scalar coefficient is nonzero.  The exponent of protected
integration is therefore the normal-ordered Gram coordinate
\[
m_R^{\mathrm{Bch}}
:=\operatorname{pr}_3\overline\Pi_X\widehat\omega_R,
\]
not the geometric support degree \(b_R^{\mathrm{geom}}\).  Equivalently,
the protected-integration label residual satisfies
\[
\mathfrak o^{I,\Pi}_R=0
\]
on the supplied finite comparison rows.  This proves only the
\(\overline\Pi_X\)-label comparison for protected integration.  It does
not prove HN-transition compatibility, a branch formula relating
\(m_R^{\mathrm{Bch}}\) to \(b_R^{\mathrm{geom}}\), the
level-\(\mathsf Z\) protected trace, or a Pfaffian line.
\end{proposition}

\begin{proof}
By Proposition~\ref{prop:protected-integration-gram-degree}, the
image of a homogeneous class of charge \(\gamma\) lies in the monomial
line
\[
\C q^{n_R(\gamma)}r^{l_R(\gamma)}s^{m_R(\gamma)}
\]
where
\[
g_R(\omega_R(\gamma))=(n_R(\gamma),l_R(\gamma),m_R(\gamma)).
\]
The supplied comparison row identifies this triple with
\[
\overline\Pi_X(\widehat\omega_R(\gamma))
=
\bigl(n_R^\Pi(\gamma),l_R^\Pi(\gamma),m_R^\Pi(\gamma)\bigr).
\]
Substitution gives the displayed monomial line and the displayed
\(s\)-degree.  Proposition~\ref{prop:geometric-borcherds-degree-separation}
separates the support map
\[
b_R^{\mathrm{geom}}(\gamma)
=\operatorname{coeff}_{[E]}(p_E)_*\operatorname{cyc}_1(\gamma)
\]
from the Borcherds coordinate
\[
\operatorname{pr}_3\overline\Pi_X(\widehat\omega_R(\gamma)).
\]
Thus \(b_R^{\mathrm{geom}}\) controls the local/wrapped support split,
whereas the protected-integration monomial is labelled by the
normal-ordered Gram coordinate.  No equality or affine relation between
these two integers is used; such a relation would require the separate
branch-comparison datum of
Definition~\ref{def:geometric-borcherds-degree-comparison-datum}.
\end{proof}

The packet
\[
\texttt{certificates/hybrid/protected\_integration\_pi\_x\_separation}
\]
verifies
\(\texttt{PROTECTED\_INTEGRATION\_PI\_X\_SEPARATION\_VERIFIED}\):
Proposition~\ref{prop:protected-integration-pi-x-separation} proves
that the protected-integration monomial exponent is
\(\overline\Pi_X(\widehat\omega_R(\gamma))\), with \(s\)-degree
\(\operatorname{pr}_3\overline\Pi_X(\widehat\omega_R(\gamma))\), and
not the geometric support degree \(b_R^{\mathrm{geom}}(\gamma)\).
The packet records transition compatibility as a separate open row.

\begin{proposition}[First-window anchor residual vanishing criterion]
\label{prop:first-window-anchor-residual-vanishing-criterion}
\textnormal{[proved as a criterion; unverified for the present first-window scaffold]}
Fix a first relation-closed finite window \(W_1\) and a height \(R\)
whose wrapped colours lie in \(W_1\).  The restriction
\[
\mathfrak o^\lambda_R|_{W_1}=0
\]
is equivalent to the simultaneous vanishing on \(W_1\) of the five
component defects in Definition~\ref{def:finite-anchor-residual}:
existence of repaired anchors \(\Lambda_{\eta,R}\), unit translation
weight after the chosen finite cover, kernel-pair losslessness,
multiplicativity on the \(LW\), \(WL\), \(WW\), and \(LWL\)
unreduced correspondences, and compatibility with the \(R'\to R\)
HN-transition maps.  In the current first-window scaffold
\[
\texttt{certificates/first\_window/k3e\_relation\_closed\_window},
\]
the compact-source window rows, target-source representative rows,
correspondence theorem rows, and strict transition rows are empty.
Consequently the present data do not prove
\(\mathfrak o^\lambda_R|_{W_1}=0\).
\end{proposition}

\begin{proof}
The equivalence is the product-zero condition for the five residual
components in Definition~\ref{def:finite-anchor-residual}.  Each
component is defined by an equality of finite data on the retained
prequotient, correspondence, or transition row; the product residual
vanishes exactly when all five equalities hold.  The first-window
obstruction ledger records no compact-source
\(\texttt{window\_closure.csv}\) rows, no
\(\texttt{target\_source\_representatives.csv}\) rows, no theorem rows
in \(\texttt{first\_window\_theorems.csv}\), and no strict transition
rows.  These missing tables contain the required first-window
existence, losslessness, multiplicativity, and transition data.  Hence
the criterion is defined, but its hypotheses are not supplied by the
current first-window scaffold.
\end{proof}

The packet
\[
\texttt{certificates/hybrid/anchor\_residual\_first\_window\_obstruction}
\]
verifies
\(\texttt{ANCHOR\_RESIDUAL\_FIRST\_WINDOW\_OBSTRUCTION\_VERIFIED}\):
Proposition~\ref{prop:first-window-anchor-residual-vanishing-criterion}
gives the exact first-window zero criterion and records that the
present first-window scaffold does not contain the rows needed to prove
it.  This packet is not a proof of row \(218\); it is the obstruction
ledger for that proof.

\subsubsection*{Finite hybrid local/wrapped factorization datum}

The full carrier at finite stage is the following supplied chain datum,
with named obstruction classes that must be trivialized.

\begin{definition}[Finite hybrid local/wrapped factorization datum]
\label{def:hybrid-wrapped-factorization-datum}
Fix a Bridgeland-style stability sector $(\sigma,S)$ in the K3-surface
lineage \cite{BridgelandK3Stability,BayerMacriProjectivity} and a
finite HN height $R$.  A finite-stage hybrid datum is
\[
\mathcal F_{X,\sigma,S,\le R}^{\mathrm{hyb}}
=\bigl(
\mathcal C_{\sigma,S,\le R},\,
\operatorname{Ran}^{\mathrm{hyb}}_{R,\mathrm{pre}}(E),\,
b_R^{\mathrm{geom}},\,
\mathcal F_R^{\mathrm{loc}},\,
\mathcal G_R^{\mathrm{wr}},\,
\mathfrak E_R^{\mathrm{loc}},\,
\mathfrak E_R^{\mathrm{hyb}},\,
\mathfrak E_R^{\mathrm{wr}},\,
\mathfrak{Flag}_R^{\mathrm{hyb}},\,
\mathsf{Corr}^{E,\mathrm{hyb}}_R,\,
\lambda_R,\,
Q_{E,R},\,
\theta^Q_R,\,
o_R,\,
I_R^{\mathrm{prot}}\bigr)
\]
with the seven clauses below.  Every stack, sheaf, operation, and
coherence lives at the fixed height $R$; the inverse limit appears
only after all finite stages and transition maps have been supplied.

\begin{enumerate}
\item \emph{Finite Hall stage, hybrid prestack, and elliptic degree.}
$\mathcal C_{\sigma,S,\le R}$ is a supplied extension-closed
charge-support HN quotient of the candidate compact Hall category
under the finite-type hypothesis $\textup{[K3$\times$E-fin]}$,
with finite-type substacks
$\mathfrak M_{\gamma,R}$ for
$\gamma\in\Gamma_{\sigma,S}^{\mathrm{HN}}(R)$,
vanishing-cycle coefficients $\Phi_{\gamma,R}$, and orientation lines
$o_{\gamma,R}$.  The carrier prestack is
\(\operatorname{Ran}^{\mathrm{hyb}}_{R,\mathrm{pre}}(E)\) of
Definition~\ref{def:hybrid-ran-prestack}; its stackification is not
used unless the descent rows of that definition are supplied.  The
geometric elliptic-degree component
\[
b_R^{\mathrm{geom}}:\Gamma_{\sigma,S}^{\mathrm{HN}}(R)\longrightarrow\Z_{\ge0}
\]
is the map of Definition~\ref{def:geometric-elliptic-degree-map}; it
splits the finite charge support into local and wrapped pieces:
\begin{equation}
\label{eq:hybrid-loc-wr-windows}
\Gamma_R^{\mathrm{loc}}=\{b_R^{\mathrm{geom}}=0\},
\qquad
\Gamma_R^{\mathrm{wr}}=\{b_R^{\mathrm{geom}}>0\}.
\end{equation}
The \(b=0\) carrier is
\(\operatorname{Ran}^{\mathrm{loc}}_{R,\mathrm{pre}}(E)\) of
Definition~\ref{def:hybrid-local-stratum-prestack}; it is the
\(J=\varnothing\) subprestack of the hybrid carrier, not the later
local coefficient sheaf.
The \(b>0\) carrier is
\(\operatorname{Ran}^{\mathrm{wr}}_{R,\mathrm{pre}}(E)\) of
Definition~\ref{def:hybrid-wrapped-stratum-prestack}; it is the
\(I=\varnothing\) subprestack whose points are wrapped prequotient
families with anchors, not the later wrapped coefficient object
\(\mathcal G_R^{\mathrm{wr}}\).
Additivity on retained extensions is the datum and criterion of
Definition~\ref{def:geometric-elliptic-degree-additivity-datum} and
Proposition~\ref{prop:geometric-elliptic-degree-additivity-criterion};
terms of height $>R$ do not enter the retained quotient.
The Borcherds trace coordinate
\(m_R^{\mathrm{Bch}}=\operatorname{pr}_3\overline\Pi_X\) is not used
to form the split in \eqref{eq:hybrid-loc-wr-windows}; any relation
between \(m_R^{\mathrm{Bch}}\) and \(b_R^{\mathrm{geom}}\) requires the
comparison datum of
Definition~\ref{def:geometric-borcherds-degree-comparison-datum}.

\item \emph{Projection-finite local sector through closed
configurations.}  Over the local carrier
\(\operatorname{Ran}^{\mathrm{loc}}_{R,\mathrm{pre}}(E)\), and for a
finite set $I$, set
\[
Z_I=\bigcup_{i\in I}\{(e,(e_j)_j)\in E\times E^I\mid e=e_i\}.
\]
For $\alpha\in\Gamma_R^{\mathrm{loc}}$ the stage supplies a
closed-configuration incidence stack
\[
\mathfrak M_{\alpha,R,I}^{\mathrm{loc,cl}}
=\{(A,\mathbf e)\in\mathfrak M_{\alpha,R}\times E^I\mid
p_E(\operatorname{Supp}A)\subset Z_I(\mathbf e)\}
\]
compatible with diagonals $E^J\to E^I$, symmetric-group descent, and
the d-critical/vanishing-cycle structure.  Open-support locality is
recovered as a limit
\[
\mathcal M_{\alpha,R}^{\mathrm{loc}}(U)
\simeq\operatorname*{colim}_{I}
\bigl(\mathfrak M_{\alpha,R,I}^{\mathrm{loc,cl}}\times_{E^I}U^I\bigr),
\]
and the configuration map $\pi_{\alpha,R,I}$ produces
\[
\mathscr H_{\alpha,R,I}^{\mathrm{loc}}
=(\pi_{\alpha,R,I})_!(\Phi_{\alpha,R,I}^{\mathrm{loc}}\otimes
o_{\alpha,R,I}^{\mathrm{loc}}).
\]
For an open $U\subset E$,
\[
\mathcal F_{\alpha,R}^{\mathrm{loc}}(U)
=\bigoplus_{I/S_I}R\Gamma_c\!\left(U^I,
j_{U,I}^{!}\mathscr H_{\alpha,R,I}^{\mathrm{loc}}\right),
\]
and $\mathcal F_R^{\mathrm{loc}}(U)$ is the direct sum over
$\Gamma_R^{\mathrm{loc}}$.  Locality is $!$-restriction on closed
configuration spaces; the open-support Borel--Moore formula is a
theorem to be derived from this atlas, not the primary datum.

\item \emph{Wrapped prequotient sector and rigidification.}  Over
\(\operatorname{Ran}^{\mathrm{wr}}_{R,\mathrm{pre}}(E)\), and for
$\eta\in\Gamma_R^{\mathrm{wr}}$, the scalar-rigidified prequotient
stack of Definition~\ref{def:e-equivariant-wrapped-prequotient-stack}
\[
\mathcal M_{\eta,R}^{\mathrm{wr,rig}}
\xrightarrow{\rho_{\eta,R}}
\mathcal M_{\eta,R}^{\mathrm{wr}}\subset\mathfrak M_{\eta,R}
\]
is supplied with a fixed origin $0_E\in E\).  The map \(\rho_{\eta,R}\)
forgets the scalar rigidification; the reduced \(E\)-quotient is not
taken in this clause.  The normalized determinant construction of
Definition~\ref{def:determinant-anchor} gives the
candidate anchor
\[
\lambda^{\det}_{\eta,R}(A)
=\det Rp_{E,*}A\otimes\mathcal O_E(-\chi(A)\cdot0_E)
\in \operatorname{Pic}^0(E)\simeq E .
\]
Its translation weight is \(\chi(A)\) by
Proposition~\ref{prop:determinant-anchor-translation-weight}.
If \(\chi(A)=0\), Proposition~\ref{prop:determinant-anchor-chi-zero-degeneracy}
shows that this determinant anchor is constant along the retained
\(E\)-translation orbit and therefore has rank-one separation defect.
The degree-zero Jordan--Hoelder divisor of
Proposition~\ref{prop:determinant-anchor-extra-data-separates-rank-two}
separates the determinant-zero rank-two test pair
\(\mathcal O_E^{\oplus2}\) and \(L\oplus L^{-1}\).  Unit-weight
finite-cover normalization, global anchor losslessness, quotient
descent, stabilizer linearization, transition compatibility, and
finite-stage population are separate rows before it becomes the final
wrapped anchor
\(\lambda_{\eta,R}:\mathcal M_{\eta,R}^{\mathrm{wr,rig}}\to E\).
The finite anchor residual \(\mathfrak o^\lambda_R\) is
Definition~\ref{def:finite-anchor-residual}; its vanishing, not its
definition, is required before the repaired anchors become the final
wrapped anchors.  After the final wrapped anchor rows are supplied, set
\[
\mathcal G_{\eta,R}^{\mathrm{wr}}
=H_*^{\mathrm{BM}}\!\left(
\mathcal M_{\eta,R}^{\mathrm{wr,rig}},
\Phi_{\eta,R}^{\mathrm{wr}}\otimes o_{\eta,R}^{\mathrm{wr}}\right),
\qquad
\mathcal G_R^{\mathrm{wr}}=\bigoplus_{\eta}\mathcal G_{\eta,R}^{\mathrm{wr}}.
\]

\item \emph{Local, mixed, and wrapped extension correspondences.}  For
disjoint opens $(U_i)\subset E$, local charges $\alpha_i$, and wrapped
charges $\eta,\eta_1,\eta_2$, the finite stage contains the
correspondence stacks
\[
\mathfrak E_{\alpha_1,\ldots,\alpha_r,R}^{\mathrm{loc}}(U_1,\ldots,U_r),
\quad
\mathfrak E_{\alpha_1,\ldots,\alpha_r;\eta,R}^{\mathrm{hyb}},
\quad
\mathfrak E_{\eta_1,\eta_2,R}^{\mathrm{wr}},
\]
where the local symbol is the \(LL\) correspondence datum of
Definition~\ref{def:local-local-correspondence-datum}.  The one-sided
mixed symbol is the \(LW/WL\) correspondence datum
of Definition~\ref{def:mixed-local-wrapped-correspondence-datum} after
collecting the finite local colours into a local colour map
\(\alpha:I\to\Gamma_R^{\mathrm{loc}}\).  It is still before the
reduced \(E\)-quotient.  The wrapped symbol is the ordered \(WW\)
correspondence datum of
Definition~\ref{def:wrapped-wrapped-correspondence-datum}.  These
stacks are supplied together with the ordered \emph{two-sided} mixed
stacks of Definition~\ref{def:two-sided-mixed-correspondence-datum},
\[
\mathfrak E_{\alpha_-;\eta;\beta_+,R}^{\mathrm{hyb}}((U_i),(V_j))
\]
classifying ordered insertions of finite local supports on both sides
of one wrapped input.  The defects of closed-configuration finite-type
properness, pull-push admissibility, and Thom--Sebastiani transport
are recorded by
\[
\mathfrak o^{\mathrm{prop},\mathrm{LL}}_R,\
\mathfrak o^{\mathrm{prop},\mathrm{mix}}_R,\
\mathfrak o^{\mathrm{prop},\mathrm{WW}}_R,
\]
and aggregate residuals
$\mathfrak o^{\mathrm{conf}}_R$,
$\mathfrak o^{\mathrm{conf{-}prop}}_R$,
$\mathfrak o^{\mathrm{adm}}_R$,
$\mathfrak o^{\mathrm{TS}}_R$.  Mixed and wrapped stacks remember
source and target anchors before quotienting by the maps
\(a^{\mathrm{LW}},a^{\mathrm{WL}},a^{\mathrm{WW}}\), and
\(a^{\mathrm{LWL}}\) of
Definition~\ref{def:source-target-anchor-memory-datum}.  Replacing
this memory by a scalar Fock factor, a Hecke factor, a Borcherds trace
degree, or a post-quotient class is not the hybrid datum.  These stacks
generate a finite $E$-equivariant correspondence category
$\mathsf{Corr}^{E,\mathrm{hyb}}_R$ with operation maps
$\mu_R^\bullet=q_!\circ\operatorname{TS}_R\circ p^*$.

\item \emph{Eight-word two-step binary flag atlas.}  Write
$\mathrm L$ for a projection-finite local input and $\mathrm W$ for a
rigidified wrapped input.  For every word
$w=(\epsilon_1,\epsilon_2,\epsilon_3)\in\{\mathrm L,\mathrm W\}^3$ and
every compatible ordered triple $(\xi_1,\xi_2,\xi_3)$ of charges and
local opens, the datum supplies a two-step flag stack
$\mathfrak{Flag}^{w}_{\xi_1,\xi_2,\xi_3,R}$ classifying filtrations
with the prescribed associated graded pieces, with comparison maps to
the iterated correspondence fibre products.  The eight words
\[
\mathrm{LLL},\ \mathrm{LLW},\ \mathrm{LWL},\ \mathrm{WLL},\
\mathrm{LWW},\ \mathrm{WLW},\ \mathrm{WWL},\ \mathrm{WWW}
\]
are all part of the datum; in mixed and wrapped words the flag stack
retains initial, intermediate, and terminal anchors before $E$-quotient.  The
two parenthesizations give the same map in the finite quotient,
\[
\mu_{\xi_1+\xi_2,\xi_3,R}\circ(\mu_{\xi_1,\xi_2,R}\otimes 1)
=\mu_{\xi_1,\xi_2+\xi_3,R}\circ(1\otimes\mu_{\xi_2,\xi_3,R}),
\]
with defect classes $\mathfrak o^{\mathrm{assoc}}_{w,R}$ required to
vanish for all eight words; the chosen $2$-morphisms must satisfy the
four-input pentagon on the four-step flag-of-flags stacks.  The
higher-coloured residual
\(\mathfrak o^{\mathrm{col}}_R\) of
Definition~\ref{def:higher-coloured-hybrid-tree-category} records the
remaining coloured configuration data, units, symmetry conventions,
refinement maps, descent, and overlap coherences for all finite HN
windows.  The eight-word atlas gives only binary associativity; full
hybrid factorization requires the higher-coloured residual to vanish
in addition.

\item \emph{Quotient-after-correspondence functor and reduced
$E$-quotient.}  All stacks and correspondences above are
$E$-equivariant before quotienting.  The quotient datum supplies a
pseudofunctor
\[
\mathsf{Corr}^{E,\mathrm{hyb}}_R\longrightarrow
\mathsf{Corr}^{\mathrm{red},\mathrm{hyb}}_R
\]
as in Definition~\ref{def:quotient-after-correspondence-pseudofunctor}.
Proposition~\ref{prop:quotient-after-correspondence-preserves-composition}
supplies the coherent \(2\)-morphism for composition.  The later
Proposition~\ref{prop:quotient-after-correspondence-preserves-base-change}
supplies preservation of the stack-level cartesian base-change
squares.  The later quotient rows supply compact-support pushforward
base change, Thom--Sebastiani transport, and chain-level comparison;
applying vanishing-cycle Borel--Moore chains gives the finite
correspondence module categories
$\mathsf{BM}^{E,\mathrm{hyb}}_R$ and
$\mathsf{BM}^{\mathrm{red},\mathrm{hyb}}_R$, and the induced functor
$Q_{E,R}$ between them.  For each operation
$\mu_R^\bullet=q_!\operatorname{TS}p^*$ the datum carries a comparison
isomorphism
\[
\theta^Q_{\mu,R}:Q_{E,R}\,q_!\operatorname{TS}p^*
\xrightarrow{\sim}
\overline q_!\overline{\operatorname{TS}}\,\overline p^*
(Q_{E,R}^{\boxtimes k}),
\]
compatible with the flag-stack coherences.  The quotient residual
lanes are refined by
\[
\mathfrak o^Q_R=
\bigl(
\mathfrak o^{Q,\mathrm{form}}_R,\
\mathfrak o^{Q,\mathrm{comp}}_R,\
\mathfrak o^{Q,\mathrm{bc}}_R,\
\mathfrak o^{Q,\mathrm{TS}}_R,\
\mathfrak o^{Q,\mathrm{flag}}_R,\
\mathfrak o^{Q,\mathrm{BM}}_R,\
\mathfrak o^{Q,\theta}_R,\
\mathfrak o^{Q,\mathrm{tr}}_{R'R}
\bigr),
\]
and these components must vanish before the reduced hybrid product is
used.  A quotient-first construction (object-stack quotient followed
by a new Hall product) is \emph{not} this datum.

\item \emph{Protected integration with Gram-trace degree.}
The finite stage carries a protected integration datum
of Definition~\ref{def:finite-protected-integration-datum} and hence a
chain map
\[
I_R^{\mathrm{prot}}:Q_{E,R}\mathcal F_{X,\sigma,S,\le R}^{\mathrm{hyb}}
\longrightarrow\C[\mathsf T_R]_{\le R}.
\]
For homogeneous \(x\) of charge \(\gamma\),
\[
I_R^{\mathrm{prot}}(Q_{E,R}x)\in
\C e_{\omega_R(\gamma)} .
\]
The finite Gram-degree row identifies
\[
e_{\omega_R(\gamma)}=q^nr^ls^m,\qquad
g_R\omega_R(\gamma)=(n,l,m),
\]
in the OP/Borcherds variables of
Proposition~\ref{conv:op-chamber}.  The normal-ordered label
comparison of
Proposition~\ref{prop:protected-integration-pi-x-separation}
identifies \(g_R\omega_R\) with \(\overline\Pi_X\widehat\omega_R\) and
excludes \(b_R^{\mathrm{geom}}\) as protected-integration exponent.
The integration defects
\[
\mathfrak o^{I,\mathrm{ch}}_R,\
\{\mathfrak o^{I,\mathrm{prod}}_{w,R}\}_w,\
\mathfrak o^{I,\mathrm{co}}_R,\
\mathfrak o^{I,\Prim}_R,\
\mathfrak o^{I,\mathrm{tr}}_{R'R},\
\mathfrak o^{I,\deg}_R,\
\mathfrak o^{I,\Pi}_R
\]
record chain-map well-typedness, multiplicativity on each word,
comultiplicativity, primitive compatibility, HN transition, the
Gram-degree law, and the exclusion of \(b_R^{\mathrm{geom}}\) as the
Borcherds exponent.
\end{enumerate}

The aggregate finite obstruction tuple of the hybrid lane is
\[
\mathfrak O_{\mathrm{hyb},R}
=(\mathfrak o^{\mathrm{stage}}_R,
\mathfrak o^{\lambda}_R,
\mathfrak o^{\mathrm{corr},\mathrm{LL}}_R,
\mathfrak o^{\mathrm{corr},\mathrm{mix}}_R,
\mathfrak o^{\mathrm{corr},\mathrm{WW}}_R,
\mathfrak o^{\mathrm{conf}}_R,
\{\mathfrak o^{\mathrm{assoc}}_{w,R}\}_w,
\mathfrak o^{\mathrm{col}}_R,
\mathfrak o^{Q,\bullet}_R,
\mathfrak o^{I,\bullet}_R,
\mathfrak o^{\mathrm{tr}}_{\mathrm{hyb},R'R}).
\]
At table level the hybrid datum is supplied by
\[
\mathrm{hybrid\_ran\_prestack\_definition},\quad
\mathrm{local\_stratum\_b0\_prestack\_definition},\quad
\mathrm{wrapped\_stratum\_bpositive\_prestack\_definition},\quad
\mathrm{geometric\_elliptic\_degree\_map},\quad
\mathrm{geometric\_elliptic\_degree\_additivity},\quad
\mathrm{geometric\_borcherds\_degree\_separation},\quad
\mathrm{positive\_elliptic\_degree\_projection},\quad
\mathrm{hybrid\_base},\quad
\mathrm{local\_configurations},\quad
\mathrm{local\_local\_correspondence\_definition},\quad
\mathrm{mixed\_local\_wrapped\_correspondence\_definition},\quad
\mathrm{wrapped\_wrapped\_correspondence\_definition},\quad
\mathrm{two\_sided\_mixed\_correspondence\_definition},\quad
\mathrm{wrapped\_prequotients},\quad
\mathrm{correspondences},\quad
\mathrm{flag\_atlas},\quad
\mathrm{higher\_coloured},\quad
\mathrm{quotient\_pseudofunctor},\quad
\mathrm{protected\_integration},\quad
\mathrm{transitions},\quad
\mathrm{scalar\_firewall}.
\]
The scalar firewall excludes ordinary Ran locality alone,
determinant anchors alone, quotient-first Hall products, Fock traces,
Borcherds \(s\)-degree alone, and scalar Pfaffian products as hybrid
carrier data.
\end{definition}

\begin{proposition}[Finite consequences of the hybrid datum;
\textit{conditional on
$\mathfrak O_{\mathrm{hyb},R}=0$}]
\label{prop:hybrid-datum-scope}
Assume the data of
Definition~\ref{def:hybrid-wrapped-factorization-datum} are supplied
for all $R$ with compatible transition maps, and
$\mathfrak O_{\mathrm{hyb},R}=0$.  Then:
\begin{enumerate}
\item the $b=0$ restriction is an ordinary support-local factorization
algebra on $\operatorname{Ran}(E)$ in the sense of Beilinson--Drinfeld
\cite[\S3.4]{BD}, equivalently a Costello--Gwilliam factorization
algebra over the smooth real curve $E_{\mathrm{top}}$
\cite[Vol.~II~\S10]{CostelloGwilliam2};
\item every operation with a wrapped input lands in a wrapped
stratum;
\item the ordered mixed maps, including the two-sided
$\mu_{\alpha_-;\eta;\beta_+,R}^{\mathrm{hyb}}$, are the only Hall
actions of projection-finite insertions on wrapped sectors in
$\mathsf{Corr}^{E,\mathrm{hyb}}_R$ before reduced $E$-quotient;
\item the eight-word atlas plus pentagon coherence yield an
associative finite-stage reduced Hall product after vanishing cycles,
orientations, HN quotienting, the quotient-after-correspondence
pseudofunctor, and the $\theta^Q_{\mu,R}$;
\item for homogeneous $x$ of charge $\gamma$, the protected integration
map is chargewise typed by
\[
I_R^{\mathrm{prot}}(Q_{E,R}x)\in\C e_{\omega_R(\gamma)};
\]
the Gram-degree and \(\overline\Pi_X\)-separation rows give
\(e_{\omega_R(\gamma)}=q^nr^ls^m\) for
\(\overline\Pi_X(\widehat\omega_R(\gamma))=(n,l,m)\), while the
local/wrapped type of \(x\) is still measured by
\(b_R^{\mathrm{geom}}(\gamma)\).
\end{enumerate}
\end{proposition}

\begin{proof}
Items (i) and (ii) are part of the elliptic-degree stratification
and Proposition~\ref{prop:positive-elliptic-degree-projects}.  Item (iii) is the
definition of the mixed operations: before reduced $E$-quotient they
are the only displayed operations remembering relative $E$-position
between finite local supports and a wrapped leg.  Items (iv) and (v)
are the eight-word coherence and the protected integration clause.
\end{proof}

\begin{corollary}[No quotient-first universal hybrid repair;
\textit{proved}]
\label{cor:no-quotient-first-hybrid-repair}
A candidate finite correspondence category for the wrapped sector
gives the source object $\mathcal F_{X,\sigma,S,\le R}^{\mathrm{hyb}}$
only if it contains the unreduced local, wrapped, and ordered mixed
source stacks of
Definition~\ref{def:hybrid-wrapped-factorization-datum}, their
anchors, the eight-word flag atlas, the quotient-after-correspondence
pseudofunctor inducing $Q_{E,R}$ after Borel--Moore chains, the comparison maps
$\theta^Q_{\mu,R}$, and protected integration with
$\mathfrak O_{\mathrm{hyb},R}=0$.  A quotient of object stacks formed
\emph{before} the mixed and wrapped correspondences, or a
determinant/Fock category carrying only the trace degree
$s^{m_R^{\mathrm{Bch}}}$, does not define the hybrid source.
\end{corollary}

\begin{proof}
By Proposition~\ref{prop:positive-elliptic-degree-projects}, a
positive elliptic-degree sector with an effective support-realization
row is not a finite-support Ran sector; only the rigidified wrapped
prequotient with anchor, together with the unreduced mixed and wrapped
correspondences, retains the relative \(E\)-position used by the source
data.  A quotient-first construction loses these anchors.
\end{proof}

\subsubsection*{Comparison with the two-stage chiral construction}

The hybrid carrier is the finite \(K3\times E\) specialization datum
attached to the two-stage construction
$\textup{CY}_d\textup{-Cat}\to E_d$-$\textup{HolFA}(X)\to
\textup{ChirAlg}(C)$ of Vol III Calabi--Yau quantum groups, read here
as a chosen specialization datum rather than
a constructed functor.  The projection-finite local
stratum is the retained holomorphic \(E_3\)-prefactorization input
restricted to the smooth elliptic factor, and the wrapped stratum is the chain-level
$K3$-fusion datum that the local stratum cannot see.  Beilinson and
Drinfeld's Ran construction \cite[\S3.4--\S3.5]{BD} provides the
$b=0$ part; the $b>0$ part is the Hall--Borcherds residual.
The finite \(E_3\)-source and specialization rows
\(\mathrm{formal\_target\_charts}\),
\(\mathrm{field\_complexes}\), \(\mathrm{e3\_operations}\),
\(\mathrm{bv\_qme}\), \(\mathrm{anomaly\_framing}\),
\(\mathrm{compact\_support}\), \(\mathrm{factorization\_descent}\),
\(\mathrm{k3\_to\_e\_specialization}\), \(\mathrm{transitions}\), and
\(\mathrm{scalar\_firewall}\) are upstream of this hybrid carrier.
They are not supplied by the hybrid carrier itself.
The eight-word two-step binary flag
atlas $w\in\{\mathrm L,\mathrm W\}^3$ is the smallest closed coherence
calculus that detects the local/wrapped split through binary
associativity; it is independent of the higher-coloured residual
$\mathfrak o^{\mathrm{col}}_R$ which closes higher symmetric arities,
units, and overlaps.

\subsubsection*{The \(D0\) specialization}

The finite-stage components of
$\mathfrak o^{\mathrm{stage}}_R$ and $\mathfrak o^{\lambda}_R$ are
constructed in
Appendix~\ref{app:G-D0-discharge}.  The wrapped-leg construction and
the eight-word atlas at finite height remain conditional on the named
hybrid obstruction tuple; the $b=0$ stratum is the classical
projection-finite Ran sector of \cite{BD}.
The hybrid carrier is the projection-finite Ran stratum on the \(b=0\)
component and the wrapped stratum relative to the named finite
obstruction tuple \(\mathfrak O_{\mathrm{hyb},R}\).

Chiral operations on $\mathrm{Ran}^{\mathrm{hyb}}(E)$ require an
orientation line on the cosection-reduced d-critical heart.  On
$K3\times E$ this orientation descends through the reduced
$E$-quotient only after the retained quotient-orientation datum is
supplied; on rank-$2$ $E$-strata the Klein-four calculation gives the
edge criterion and leaves the degree-one linearization character to be
set to zero.  Entry $(\mathrm L 2)$, constructed in
\S\ref{subsec:cosection-twisted-orientation}, is this orientation datum.
